% concatenated from main.tex
\documentclass[12pt]{amsart}
\usepackage{amsfonts}
\usepackage{mathrsfs}
\usepackage{bbm}%
\usepackage{bm}
\usepackage{bbding}

\usepackage{fontawesome5}
\usepackage[mathcal]{euscript}
\usepackage{graphicx}
\usepackage{amsthm,amscd}
\usepackage{calc}
\usepackage{mathtools}
\usepackage{mathrsfs,dsfont}
\usepackage{CJK,fancyhdr,amscd}
\usepackage[dvipsnames]{xcolor}
\usepackage{anysize}
\usepackage{amsmath,amssymb,amsfonts}
\usepackage{epsfig,enumerate}
\usepackage{latexsym}
\usepackage{indentfirst, latexsym}
\usepackage{graphics}
\usepackage{tikz-cd}
\usepackage[right]{lineno}
\usepackage[all,poly,knot]{xy}
\usepackage{xpatch}
%\usepackage{hyperref}
\usepackage[pagebackref]{hyperref}

%\usepackage[notcite, notref]{showkeys}

\hypersetup{
 	colorlinks=true,
 	linkcolor=Blue,
 	filecolor=Yellowblue,
 	citecolor = WildStrawberry,      
 	urlcolor=cyan,
}

\xpatchcmd{\step}{%
	\normalfont\scshape\centering}{%
	\normalfont\scshape}{\typeout{Success}}{\typeout{Failure}}%

\allowdisplaybreaks[2]

\providecommand{\U}[1]{\protect\rule{.1in}{.1in}}

\numberwithin{equation}{section}

\newtheorem{theorem} {Theorem} [section]
\newtheorem{proposition}[theorem]{Proposition}
\newtheorem{corollary}  [theorem]     {Corollary}
\newtheorem{lemma}  [theorem]     {Lemma}
\theoremstyle{definition}
\newtheorem{example}  [theorem]     {Example}
\newtheorem{question}  [theorem]     {Question}
\newtheorem{prob}[theorem]{Problem}
\newtheorem{step}{Step}
\theoremstyle{plain}
%\addtocounter{step}{-1}
\renewcommand{\thestep}{$($\Roman{step}$)$}

\newtheorem{definition}  [theorem]     {Definition}
\newtheorem{notation}  [theorem]     {Notation}
\newtheorem{Convention}  [theorem]     {Convention}
\newtheorem{Claim}  [theorem]     {Claim}
\newtheorem{Construction}  [theorem]     {Construction}
\newtheorem{conjecture}  [theorem]     {Conjecture}
\newtheorem{observation}  [theorem]     {Observation}
\newtheorem{main}  [theorem]{Main Theorem}
\newtheorem*{thm A*}{Theorem A}
\newtheorem*{thm B*}{Theorem B}
\newtheorem*{thm C*}{Theorem C}
\newtheorem*{thm D*}{Theorem D}

\font\strange=msbm10 \pagestyle{plain}
\theoremstyle{definition}
\newtheorem{remark}  [theorem]     {Remark}
\newtheorem{type}{Type}
\newcommand{\dz}{d \overline{z}}
\newcommand{\pz}{\partial\overline{z}}
\newcommand{\w}{\wedge}
\newcommand{\im}{\mathrm{im}}
\renewcommand{\1}{\mathds{1}}
\newcommand{\G}{\mathbb{G}}
\newcommand{\db}{\overline{\partial}}
\newcommand{\dbs}{\overline{\partial}^*}
\newcommand{\nb}{\nabla}
\newcommand{\lc}{\lrcorner}
\newcommand{\pts}{\partial^*}
\newcommand{\vol}{\mathrm{vol}}
\newcommand{\sqa}{\square}
\newcommand{\osqa}{\overline{\square}}
\newcommand{\lk}{\left(}
\newcommand{\rk}{\right)}
\newcommand{\cL}{\mathcal{L}}
\newcommand{\Om}{\Omega}
\newcommand{\btheorem}{\begin{theorem}}
	\newcommand{\etheorem}{\end{theorem}}
\newcommand{\bproposition}{\begin{proposition}}
	\newcommand{\eproposition}{\end{proposition}}
\newcommand{\bdefinition}{\begin{definition}}
	\newcommand{\edefinition}{\end{definition}}
\newcommand{\bcorollary}{\begin{corollary}}
	\newcommand{\ecorollary}{\end{corollary}}
\newcommand{\bproof}{\begin{proof}}
	\newcommand{\eproof}{\end{proof}}
\newcommand{\beq}{\begin{equation}}
	\newcommand{\eeq}{\end{equation}}
\newcommand{\ee}{\end{eqnarray*}}
\newcommand{\be}{\begin{eqnarray*}}
\newcommand{\hooklongrightarrow}{\lhook\joinrel\longrightarrow}
\newcommand{\elemma}{\end{lemma}}
\newcommand{\blemma}{\begin{lemma}}
\newcommand{\ts}{\otimes}
\newcommand{\om}{\omega}
\newcommand{\bbeta}{\boldsymbol{\beta}}
\renewcommand{\>}{\rightarrow}
\newcommand{\sL}{{\mathcal L}}
\newcommand{\p}{\partial}
\newcommand{\R}{\mathbb R}
\newcommand{\Q}{\mathbb Q}
\newcommand{\bp}{\overline\partial}
\renewcommand{\H}{\mathbb H}
\newcommand{\la}{\langle}
\newcommand{\ra}{\rangle}
\newcommand{\ssL}{{\mathfrak L}}
\newcommand{\ssX}{{\mathfrak X}}
\newcommand{\bd}{\begin{enumerate} }
	\newcommand{\ed}{\end{enumerate}}
\renewcommand{\labelenumi}{$ ( \arabic{enumi} )$}


\setlength{\hoffset}{0in} \setlength{\voffset}{0in}
\setlength{\oddsidemargin}{0in} \setlength{\evensidemargin}{0in}
\setlength{\marginparsep}{0in} \setlength{\topmargin}{0in}
\setlength{\headheight}{0in} \setlength{\headsep}{0in}
\setlength{\footskip}{0.375in} \setlength{\textwidth}{6.268in}
\setlength{\textheight}{9.693in}

\def\p{\partial}
\def\o{\overline}
\def\b{\overline}
\def\mb{\mathbb}
\def\mc{\mathcal}
\def\n{\nabla}
\def\v{\varphi}
\def\w{\wedge}
\def\om{\omega}
\def\t{\triangle}
\def\l{\lrcorner}

\usepackage{adjustbox}
\usepackage{lineno}
%\linenumbers
\usepackage{float}
\providecommand*{\propositionautorefname}{Proposition}

\begin{document}

\title{On Pseudo-Effectivity and Volumes of Adjoint Classes in K\"ahler Families with Projective Central Fiber}


\begin{abstract}
This paper is devoted to studying the deformation behavior of pseudo-effective canonical divisors and volumes of adjoint classes in K\"ahler families. Based on recent developments in the K\"ahler minimal model program, for flat families with fiberwise canonical singularities, we establish the global stability of the pseudo-effectivity of canonical divisors and uniruledness, assuming in addition that one fiber is projective, while the same conclusion for K\"ahler threefolds is also true without the projectivity assumption of the central fiber.  For a smooth K\"ahler family whose central fiber is projective with a big adjoint class, we show that its volume remains locally constant. Finally, using the (relative) minimal model program for K\"ahler threefolds, we verify the deformation invariance of volumes of adjoint classes and plurigenera for smooth families of K\"ahler threefolds, thereby confirming Siu’s invariance of plurigenera conjecture in dimension three.
\end{abstract}
	
	\author{Christopher D. Hacon}
	\address{Department of Mathematics, University of Utah, Salt Lake City, UT 84112,
USA}
\email{hacon@math.utah.edu}
 \author{Yi Li}
    \address{Department of Mathematics, Wuhan University, Hubei, Wuhan 430074, China}
    \email{yilimath@whu.edu.cn}
    \author{Sheng Rao}
    \address{Department of Mathematics, Wuhan University, Hubei, Wuhan 430074, China}
    \email{likeanyone@whu.edu.cn}
	\date{\today}
	\thanks{The first author was partially supported by  NSF research grant   DMS-2301374 and
by a grant from the Simons Foundation SFI-MPS-MOV-00006719-07. The second and third authors were partially supported by NSFC (Grant No. 12271412, W2441003) and Hubei Provincial Innovation Research Group Project (Grant No. 2025AFA044).}
    \keywords{K\"ahler minimal model program, volume, pseudo-effectivity}
	
	\maketitle
	\setcounter{tocdepth}{1}
	\tableofcontents

%%%%%%%%%%%%%%%%%%%%%%%%%%%%%%%%%%%%%%%%%%%%%%%%%%%%%%%%%%%%%%%%%%%%%%%%%%%%%%%%%%%%%%%%%%%%%%%%%%%%%%%%%%%%%%%%%%%%%%%%%%%%%%%%%
%%%%%%%%%%%%%%%%%%%%%%%%%%%%%%%%%%%%%%%%%%%%%%%%%%%% Section Introduction %%%%%%%%%%%%%%%%%%%%%%%%%%%%%%%%%%%%%%%%%%%%%%%%%%%%%%%
%%%%%%%%%%%%%%%%%%%%%%%%%%%%%%%%%%%%%%%%%%%%%%%%%%%%%%%%%%%%%%%%%%%%%%%%%%%%%%%%%%%%%%%%%%%%%%%%%%%%%%%%%%%%%%%%%%%%%%%%%%%%%%%%%
\section{Introduction}\label{Introduction}

One of the central problems in complex algebraic geometry is the classification of compact complex varieties. Plurigenera, volumes, and positivity properties are fundamental invariants associated to a variety and play a crucial role in classification problems. The aim of this paper is to study the deformation behavior of pseudo-effective canonical divisors, plurigenera and volumes of adjoint classes. 

The first part of this paper is devoted to studying the deformation behavior of pseudo-effective canonical divisors. Given a smooth projective family, H. Tsuji \cite{Tsuji02} showed that if the central fiber has a pseudo-effective canonical divisor, then the nearby fibers also have pseudo-effective canonical divisors. It is natural to ask whether similar results hold for K\"ahler families. Recently, W. Ou \cite{Ou25} proved the BDPP conjecture \cite[Conjecture 1.2]{BDPP}, which characterizes K\"ahler uniruledness:
\begin{theorem}[{\cite[Theorem 1.1]{Ou25}}]\label{Theorem Ou25}
Let X be a compact K\"ahler manifold. Then X is uniruled if and only if the canonical line bundle $K_X$ is not pseudo-effective.
\end{theorem}
Combining this with Fujiki--Levine's deformation stability of uniruledness \cite{Fujiki,Levine}, one obtains the deformation stability of pseudo-effectivity of canonical divisors for smooth K\"ahler families (cf. \cite[Corollary 5.5]{CP25}). 

Combining recent developments in the transcendental MMP \cite{DH24MMP}, Das--Hacon's transcendental base point free theorem for projective varieties \cite{DH24MMP}, and Koll\'ar's extension of MMP techniques in \cite{DefGeneralType}, we prove the local stability of pseudo-effective canonical divisors for K\"ahler families whose fibers are allowed to have canonical singularities in Theorem \ref{Deformation invariance of pseudo-effectivity}. Then based on Theorem \ref{Deformation invariance of pseudo-effectivity}, and using a relative Douady space argument in Proposition \ref{pseflocalglobal}, we can prove the following global stability result for the pseudo-effectivity of the canonical divisors. 
\begin{theorem}[Theorem \ref{GlobalPseff}]\label{MainPSEF}
Let $f: X \to S$ be a family from a K\"ahler variety $X$ onto a smooth, connected, and relatively compact curve $S$. Assume that $X_0$ is projective, and $X_t$ have canonical singularities for all $t\in S$. Then $K_{X_0}$ is pseudo-effective if and only if $K_{X_t} $ are pseudo-effective for all $t \in S \setminus \{0\}$.
\end{theorem}

With a similar approach to Theorem \ref{MainPSEF} and the characterization of uniruledness for K\"ahler varieties with canonical singularities given in Theorem~\ref{Theorem Ou25}, we prove the global deformation stability of uniruled structures with canonical singularities.
\begin{corollary}[Corollary \ref{DefUniruled}, \cite{DefGeneralType}]\label{MainDefUniruled}
Let $f: X \to S$ be a family from a K\"ahler variety $X$ onto a smooth, connected, and relatively compact curve $S$. Assume that $X_0$ is projective, and $X_t$ have canonical singularities for all $t\in S$. Then $X_0$ is uniruled if and only if $X_t$ are uniruled for all $t \in S \setminus \{0\}$.
\end{corollary}

Note that if the MMP and the transcendental base point free conjecture hold on the central fiber, then the same proofs of Theorem~\ref{MainPSEF} and Corollary~\ref{MainDefUniruled} apply. Consequently, the same conclusions of Theorem~\ref{MainPSEF} and Corollary~\ref{MainDefUniruled} for K\"ahler families of threefolds remain valid without assuming that the central fiber is projective.
It is reasonable to propose the more general:
\begin{conjecture}Let $f: X \to S$ be a family from a K\"ahler variety $X$ onto a smooth, connected, and relatively compact curve $S$. Assume that $X_t$ have canonical singularities for all $t\in S$. Then $K_{X_0}$ is pseudo-effective if and only if $K_{X_t} $ are pseudo-effective for all $t \in S \setminus \{0\}$, and $X_0$ is uniruled if and only if $X_t$ are uniruled for all $t \in S \setminus \{0\}$.
\end{conjecture}


Next, we turn our attention to the deformation invariance problem for plurigenera and volumes of adjoint classes. There is a long history of results related to the deformation invariance of plurigenera for projective or Moishezon morphisms, see \cite{Siu98,Kawamata99,Kawa99,Siu02,Tsuji02,Naka04,Taka06,Paun07,BP12,HMX13,FS20,RT1,RT2,LW26} etc. In \cite{Siu98,Siu02}, Siu proved the deformation invariance of plurigenera for smooth projective families of manifolds using Ohsawa--Takegoshi extension techniques. 
Later, Kawamata \cite{Kawa99} and Nakayama \cite{Naka04} gave algebraic proofs of a singular version of Siu's theorem for projective families of varieties of general type with canonical singularities and terminal singularities, respectively. 
Takayama \cite{Taka06} proved the invariance of plurigenera for Moishezon morphisms whose fibers have canonical singularities using a modified version of the Ohsawa--Takegoshi extension theorem. Rao--Tsai \cite{RT1,RT2} developed a bimeromorphic embedding theorem, and using this, they proved that smooth fiberwise Moishezon families and flat Moishezon families with uncountably many fibers of general type with canonical singularities both satisfy the deformation invariance of plurigenera. A  (possibly incomplete) introduction on the progress of the deformation invariance of plurigenera can be found in \cite[Introduction]{RT2}.

Nevertheless, the problem is widely open for smooth K\"ahler families, as conjectured by Siu. See the recent progress in \cite{CP23Pluri}. 
\begin{conjecture}[{\cite[Conjecture 0.4]{Siu02}, \cite[Conjecture 1.4]{CP23Pluri}}]\label{SiuConj}
Let $\pi: X \rightarrow \Delta$ be a smooth family of compact K\"ahler manifolds over the unit disk in $\mathbb{C}$ with fiber $X_t$. Then for any positive integer $m\ge 1$, the $m$-genus $
P_m\left(X_t\right):=\operatorname{dim} H^0\left(X_t, m K_{X_t}\right)
$ is independent of $t\in \Delta$.
\end{conjecture}

The cohomological methods used in the projective setting do not directly apply to transcendental classes on a K\"ahler variety. However, thanks to the recent developments in \cite{DH24MMP,DHY23}, we are able to treat transcendental classes in a way similar to boundary divisors. Using these techniques, we establish the deformation invariance of the volumes of adjoint classes under the assumption that the central fiber is projective with a big adjoint class.
\begin{theorem}[{Theorem \ref{InvarVol}}]\label{0InvarVol}
Let $f: X \to S$ be a smooth family from a K\"ahler manifold $X$ onto a smooth, connected, and relatively compact curve $S$. For a point $0\in S$, assume that $(X,B + \bbeta)$ is a generalized K\"ahler pair such that the boundary divisor $B$ does not contain the fiber $X_0$ and let $\omega$ be a K\"ahler class on $X$ with $\omega_{X_t}:=\omega|_{X_t}$ its restriction to the fiber $X_t$. 

Assume that there exists $\delta>0$ such that $K_{X_0} + B_0 +  \boldsymbol{\beta}_{X_0}+\delta {\omega}_{X_0}$ is big and $(X_0, B_0+  { \boldsymbol{\beta}}_{X_0}+ \delta\boldsymbol{\omega}_{X_0})$ is projective, with canonical singularities, $\lfloor{ B_0\rfloor} = 0$ and $N(K_{X_0 }+B_0 + \bbeta_{X_0})\wedge B_0  = 0$. Then the volume  function 
$$t \longmapsto \mathrm{vol}(K_{X_t}+B_t +  \boldsymbol{\beta}_{X_t} +\delta{\omega}_{X_t})$$ 
is constant on a Euclidean open neighborhood $U \subset S$ of $0$.
\end{theorem}

We also obtain the following variant under the assumption that $(X,B)$ is log smooth over $S$, without assuming that $N(K_{X_0} + B_0 + \bbeta_{X_0})\wedge B_0  = 0$. 
\begin{theorem}[Theorem \ref{InvarVol2}]\label{0InvarVol2}
Let $f:X\to S$ be a smooth family from a K\"ahler manifold $X$ onto a smooth, connected, relatively compact curve $S$. Assume that $(X,B+\bbeta)$ is a generalized klt pair such that $(X,B)$ is log smooth over $S$ and $\bbeta =\overline \beta $ where $\beta$ is nef over $S$, and $X_0$ is projective. If $K_{X_0}+B_0+\bbeta _{X_0}$ is big, then the volume function
\[t\longmapsto {\rm vol}(K_{X_t}+B_t+\bbeta _{X_t})\]
is constant on some Euclidean neighborhood $U\subset S$ of $0$.
\end{theorem}


Since the (relative) minimal model program and the abundance conjecture are completely established for K\"ahler threefolds by the series of works \cite{HP16,CHP16,CHP23Erratum,DHY23,DH24,DHP24,DO24,DO24II,DH20}, etc, we prove the deformation invariance of volumes (as well as plurigenera) for smooth families of K\"ahler threefolds without any additional assumption on $X_0$. In particular, this partially resolves Conjecture \ref{SiuConj} in the case of threefolds.
\begin{theorem}[Theorem \ref{DefVol3fold}]\label{0DefVol3fold}
Let $f: X \to S$ be a smooth family of compact K\"ahler threefolds over a smooth, connected, and relatively compact curve $S$. Then for any integer $m\ge 1$, the $m$-genus $P_m(X_t)$ is constant for any $t\in S$. Furthermore, assume that there exists a K\"ahler class $\omega_X$ on $X$ such that $\omega_{X_t} :  =  \omega_X |_{X_t}$. Then for every $\delta\ge 0$, the function $$t \longmapsto \operatorname{vol}(K_{X_t}+\delta\,\omega_{X_t})$$is constant for any $t\in S$.
\end{theorem}
Let us highlight some key points of our proofs. To study the deformation stability of pseudo-effectivity of canonical divisors in Theorem \ref{MainPSEF}, we will make use of some MMP arguments. %There are two major difficulties we will face.
%First, unlike the standard setting of \cite{BirkarCHM}, we do not know whether the MMP terminates. To overcome this difficulty, we argue by contradiction that the relative log  is not pseudo-effective, and then introduce a pseudo-effective threshold, in a manner similar to the proof of the existence of Mori fiber spaces (\cite[Corollary 1.3.2]{BirkarCHM} and also \cite[Theorem 3.6.6]{KawamataJaing24}). Using that positive pseudo-effective threshold, we can guarantee the termination of the MMP. 
A key complication is that, without the projectivity assumption of the family, it is not a priori clear whether the classical MMP framework can be employed here. However, because we assume that the central fiber is projective, thanks to \cite{DH24MMP}, it is possible to run the MMP on the central fiber. We then use the extension technique from \cite[Theorem 2]{DefGeneralType} as in Theorem \ref{ExtendMMP} to extend the MMP to nearby fibers. Note that the transcendental base point free Theorem \ref{DHTransbpf} plays an important role in guaranteeing the termination of the MMP for nearby fibers, provided termination holds for the central fiber.

In order to prove the deformation invariance of the volumes in Theorem \ref{0InvarVol}, we make use of a MMP type argument via \cite{DH24MMP} as in the proof of Theorem \ref{MainPSEF}. Since we assume that the central fiber has a big adjoint class, the MMP with scaling on the central fiber terminates. By base point freeness, the MMP also terminates on nearby fibers. Hence, without loss of generality, we may assume that all fibers have nef adjoint class, so that the volume can be computed via top intersection numbers. In the case of divisors, deformation invariance of the volume follows from the constancy of the Hilbert polynomial in flat families; however, we do not have a good substitute for this in the case of transcendental classes. Fortunately, after running the MMP, the singularities are still concentrated on a codimension 2 subset, and this allows us to apply a singular version of Stokes’ formula to deduce the constancy of the volume for adjoint classes. Using a similar argument, we can prove the deformation invariance of the volume of adjoint classes for smooth families of K\"ahler threefolds as in Theorem \ref{0DefVol3fold}.


In this paper, a \emph{family} $\pi: X\to S$ will refer to a proper surjective morphism from a normal complex analytic space
to a smooth connected curve whose fibers $X_t:=\pi^{-1}(t)$ $(t\in S)$ are connected (possibly reducible) complex analytic spaces, as in \cite{CRC25}. In particular, we call $\pi$ a \emph{smooth family} if it is a flat \emph{submersion}  (i.e., the relative cotangent sheaf $\Omega^1_{X|S}$ is locally
free of rank $\dim X -\dim S$). Note that in a smooth family, both the total space and each fiber are smooth. We also assume that the general fiber of a family is irreducible. The “general fiber” of a family refers to the one over a nonempty analytic Zariski open subset of the base. 
% Note that the fibers of a family are compact and thus automatically countable at infinity; the total
% space of a family is also automatically countable at infinity by the properness condition and the
% fact that a smooth curve is either (compact) projective or Stein. 

\textbf{Organization of the paper.} 
The paper is organized as follows. Section \ref{Preliminaries} establishes the foundations by reviewing generalized pairs, positivity notions such as nef, big, and pseudo-effective classes, the volume of cohomology classes, and basic results on bigness and pseudo-effectivity in families. Section \ref{Section: SingandPositivty} introduces singularities and positivity of line bundles in families, including deformation of singularities, relative ampleness, relative pseudo-effectivity, and criteria for global generation and semi-ampleness. Section \ref{Section: ExtMMP} develops the main tools needed for the main results, namely the extension of the analytic minimal model program to nearby fibers and the deformation openness of nef and big adjoint classes. Section \ref{Section: Deformation pseff} proves the global stability of  pseudo-effectivity of canonical divisors and uniruledness in K\"ahler families with one projective fiber. Section \ref{Section: DeformationVol} establishes invariance of the volumes of adjoint classes, first when the projective central fiber has a big adjoint class and then, in the case of K\"ahler threefolds without projectivity assumption.\\
\textbf{Acknowledgment.} 
The authors would like to express their gratitude to Professor Mihai P\u{a}un for informing us of the singular version of Stokes’ formula, Professors J\'anos Koll\'ar, Wenhao Ou and Omprokash Das for answering our questions and pointing out several errors in our preliminary version, and Professor Song Sun for bringing Junsheng Zhang’s preprint \cite{JunSheng} to our attention. The last two authors also thank Professors Junyan Cao and Xiaojun Wu for a useful discussion on pseudo-effectivity and uniruledness during the conference ``Families of K\"ahler spaces" at CIRM of France in last April.
%%%%%%%%%%%%%%%%%%%%%%%%%%%%%%%%%%%%%%%%%%%%%%%%%%%%%%%%%%%%%%%%%%%%%%%%%%%%%%%%%%%%%%%%%%%%%%%%%%%%%%%%%%%%%%%%%%%%%%%%%%%%%%%%%
%%%%%%%%%%%%%%%%%%%%%%%%%%%%%%%%%%%%%%%%%%%%%%%%%%%% Section2. Preliminaries %%%%%%%%%%%%%%%%%%%%%%%%%%%%%%%%%%%%%%%%%%%%%%%%%%%%
%%%%%%%%%%%%%%%%%%%%%%%%%%%%%%%%%%%%%%%%%%%%%%%%%%%%%%%%%%%%%%%%%%%%%%%%%%%%%%%%%%%%%%%%%%%%%%%%%%%%%%%%%%%%%%%%%%%%%%%%%%%%%%%%%
\section{Preliminaries: generalized K\"ahler pairs and positivities of classes}\label{Preliminaries}
We first introduce some preparatory notions and results to be used in the proofs of our main theorems in Sections \ref{Section: Deformation pseff} and \ref{Section: DeformationVol}.
\subsection{Generalized pairs}
Generalized pairs were first introduced in \cite{BZ16} in the study of the effective Iitaka fibration conjecture and later generalized to the K\"ahler setting in \cite{DHY23}. With this additional flexibility, we can run the transcendental minimal model program for generalized K\"ahler pairs. Let us first recall some basic definitions. 

Let $X$ be a \emph{complex analytic variety}, i.e., an irreducible reduced complex space. Let $\mathcal{H}_X$ be the sheaf of real parts of holomorphic functions multiplied with $\sqrt{-1}$. A $(1,1)$-form (resp. $(1,1)$-current) with local potentials on $X$ is a global section of the quotient sheaf $\mathcal{A}_X^0 / \mathcal{H}_X$ (resp. $\mathcal{D}_X / \mathcal{H}_X$), where $\mathcal{A}_X^0$ is the sheaf of smooth real valued functions and $\mathcal{D}_X$ is the sheaf of  distributions. We define the \emph{Bott--Chern cohomology} to be $$H^{1,1}_{\rm{BC}}(X,\mathbb{R}) : = H^1(X,\mathcal{H}_X).$$

Given a closed $(1,1)$-form $\alpha$ with local potentials (resp.\ a closed $(1,1)$-current $T$ with local potentials), we denote by $[\alpha]$ (resp.\ $[T]$) the associated Bott--Chern class. Without causing confusion, given an $\mathbb{R}$-Cartier divisor $D$, we will also denote by $[D]$ the associated cohomology class in $H^{1,1}_{\mathrm{BC}}(X,\mathbb{R})$. When there is no danger of ambiguity, we will sometimes write a class $\alpha \in H^{1,1}_{\mathrm{BC}}(X,\mathbb{R})$ without explicitly specifying its representative.

A complex analytic variety $X$ is a \emph{K\"ahler variety} if there exists a K\"ahler form $\omega$ on $X$, i.e., a positive closed real $(1,1)$-form $\omega \in \mathcal{A}_{\mathbb{R}}^{1,1}(X)$ such that for every singular point $x \in X$ there exists an open neighborhood $U \subseteq X$ of $x$ and a closed embedding $\iota_U: U \hookrightarrow V$ into an open subset $V \subseteq \mathbb{C}^N$, and a strictly plurisubharmonic $C^{\infty}$ function $f: V \rightarrow \mathbb{R}$ with $\left.\omega\right|_{U \cap X_{\mathrm{sm}}}=\left.(i \partial \bar{\partial} f)\right|_{U \cap X_{\mathrm{sm}}}$.


Let $f: X \rightarrow S$ be a proper morphism of normal K\"ahler varieties with relatively compact $S$, and $\beta$  a closed $(1,1)$-current with local potentials. We say that the class $[\beta] \in H_{\mathrm{BC}}^{1,1}(X,\mathbb{R})$ is \emph{relatively K\"ahler} (also referred to as \emph{$f$-K\"ahler} or \emph{K\"ahler over $S$}) if $\left[\beta+f^* \omega_S\right] \in H_{\mathrm{BC}}^{1,1}(X,\mathbb{R})$ is a K\"ahler class for some K\"ahler form $\omega_S$ on $S$.

Let $X$, $Y$ and $S$ be complex analytic varieties. A mapping $\varphi$ of $X$ into the power set of $Y$ over $S$ is called a \emph{meromorphic $S$-mapping} of $X$ into $Y$ (we shall write $\varphi: X \dashrightarrow Y/S$), if $\varphi$ satisfies the following conditions:\\
(1) The graph $G_{\varphi}=\{(x, y) \in X \times_S Y \mid y \in \varphi(x)\}$ of $\varphi$ is an irreducible analytic subset in ${X} \times_S {Y}$;\\
(2) The projection map ${p}_{{X}}: {G}_{\varphi} \to {X}$ is a proper modification. \\
A meromorphic $S$-mapping $\varphi: {X} \dashrightarrow {Y}/S$ of complex varieties is called a \emph{bimeromorphic $S$-mapping} if ${p}_{{Y}}: {G}_{\varphi} \to {Y}$ is also a proper modification. 

In what follows, without causing any confusion, we will use the terms `birational' and `bimeromorphic' interchangeably in the analytic setting.

A compact complex analytic variety is said to be in the \emph{Fujiki class $\mathcal{C}$} if it is bimeromorphic to a compact K\"ahler manifold. A compact complex analytic variety is said to be a \emph{Moishezon variety} if it is bimeromorphic to a projective manifold. 

Let $X$ be a normal complex analytic variety. A sheaf $\mathcal{F}$ on $X$ is called a \emph{divisorial sheaf} if $\mathcal{F}$ is a reflexive sheaf of rank~$1$. The set of all divisorial sheaves, denoted by $W(X)$, forms a group under the reflexive tensor product (that is, for $\mathcal{F}, \mathcal{G} \in W(X)$, we define their \emph{reflexive tensor product} by
$\mathcal{F} \,\hat{\otimes}\, \mathcal{G} := (\mathcal{F} \otimes \mathcal{G})^{**} \in W(X)$). A normal complex analytic variety $X$ is said to be \emph{$\mathbb{Q}$-factorial} if every divisorial sheaf is $\mathbb{Q}$-Cartier; that is, for every $\mathcal{F} \in W(X)$ there exists $m \in \mathbb{N}$ such that the reflexive power $\mathcal{F}^{[m]}$ is locally free.
 

As explained in \cite[II.2, p. 54]{Naka04}, for an $n$-dimensional normal complex analytic variety $X$, the \emph{canonical sheaf} $\omega_X$ is the unique reflexive sheaf whose restriction to $X_{\text {reg }}$ is isomorphic to the sheaf $\Omega_{X_{\text {reg }}}^n$ of germs of holomorphic differential $n$-forms. A non-zero meromorphic $n$-form $\eta$ on $X_{\text {reg }}$ is regarded as a meromorphic section of $\omega_X$. The associated divisor $\operatorname{div}(\eta)$ is called the \emph{canonical divisor} of $X$ and is denoted by $K_X$ even though it depends on the choice of $\eta$. Some complex analytic varieties do not admit any non-zero meromorphic section of $\omega_X$. However, we use the symbol $K_X$ as a formal divisor class with an isomorphism $\mathcal{O}_X(K_X)\cong \omega_X$ and call it the \emph{canonical divisor} of $X$. 

\begin{definition}[{Generalized pair, \cite[Definition 2.7]{DHY23}}]
Let $f: X \rightarrow S$ be a proper surjective morphism from a normal complex variety $X$ to a relatively compact complex variety $S$, $\nu: X^{\prime} \rightarrow X$ a resolution, and  $B^{\prime}$ an $\mathbb{R}$-divisor on $X^{\prime}$ with simple normal crossing support such that $B:=\nu_* B^{\prime} \ge 0$, and $ \boldsymbol{\beta}$ a closed b-$(1,1)$ current. We say that $(X, B+ { \boldsymbol{ \beta}})$ is a \emph{generalized pair} if\\
\emph{(1)} $ { \boldsymbol{\beta}}$ is a positive closed b-$(1,1)$ current that descends to $X^{\prime}$,\\
\emph{(2)} $\left[ { \boldsymbol{\beta}}_{X^{\prime}}\right] \in H_{\mathrm{BC}}^{1,1}\left(X^{\prime},\mathbb{R}\right)$ is nef over $S$, and\\
\emph{(3)} $\left[K_{X^{\prime}}+B^{\prime}+ { \boldsymbol{\beta}}_{X^{\prime}}\right]=\nu^* \gamma$ for some $\gamma \in H_{\mathrm{BC}}^{1,1}(X,\mathbb{R})$.\\
If additionally $X$ is a K\"ahler variety, then we call $(X, B + \boldsymbol{\beta})$ a \emph{generalized K\"ahler pair}. Here a \emph{b-$(1,1)$ current} is a collection of closed bi-degree $(1,1)$ currents $\boldsymbol{\beta}_{X^{\prime}}$ on all proper bimeromorphic models $X^{\prime} \rightarrow X$ such that if $p: X_1 \rightarrow X_2$ is a bimeromorphic morphism of proper models of $X$, then $p_* \boldsymbol{\beta}_{X_1}=\boldsymbol{\beta}_{X_2}$. We say that $\bbeta$ \emph{descends to $X'$} or equivalently $\bbeta =\overline { \bbeta _{X'}}$ if on any higher model $\pi:X''\to X'$ we have $\bbeta_{X''}=\pi^* \bbeta_{X'}$. 
\end{definition}
Following the usual convention, we will sometimes denote by $\boldsymbol{\beta}_X$ the trace of the b-$(1,1)$ current $\boldsymbol{\beta}$ on $X$, by $\boldsymbol{\beta}_{X'}$ its trace on $X'$, and by $\boldsymbol{\beta}$ the actual b-$(1,1)$-current.

\begin{definition}[Generalized discrepancies]\label{logdiskrep}
Let $\left({X}, {B}+ { \boldsymbol{\beta}}\right)$ be a generalized pair. Let $E$ be a prime divisor on some birational model of ${X}$. We define the {generalized  discrepancy} of $E$ with respect to the above generalized pair as follows. Let $f:X'\to X$ be a log resolution, such that:\\ 
\emph{(a)} $B'$ is a $\mathbb{Q}$-divisor with simple normal crossing support;\\ 
\emph{(b)} $ { \boldsymbol{\beta}}$ is the positive b-(1,1) current that descends to $X'$ so that  $\bbeta _{X'}$ is nef and $\bbeta =\overline { \bbeta _{X'}}$;\\ 
%(c) $[ { {\bbeta'}}]\in H^{1,1}_{\rm{BC}}(X', \mathbb{R})$ is nef;\\
\emph{(c)} $[K_{X'}+ B'+  { {\bbeta}}_{X'}]  =f^* \gamma$ for some $\gamma\in H^{1,1}_{\rm{BC}}(X, \mathbb{R})$. We can then write
$$
{K}_{{X}'}+{B}'+ { {\bbeta}}_{X'}=f^*\left({K}_{{X}}+{B}+ { {\bbeta}_X}\right).
$$
The \emph{generalized discrepancy} $a(X,B, { {\boldsymbol{\beta}}} , E)$ of $E$ is defined to be $-\text{coeff}_E(B')$.
\end{definition}

Using generalized discrepancies, we define various types of singularities for generalized pairs:

A generalized  pair $(X,B +  { { \boldsymbol{\beta}}} )$ is said to have \emph{canonical singularities} if %$K_X+B$ is a $\mathbb{Q}$-Cartier divisor and 
the generalized discrepancy for any exceptional divisors satisfies $a(X,B ,  { \boldsymbol{\beta}} , E) \ge 0$. 

A generalized pair $(X,B +  { \boldsymbol{\beta}} )$ is said to have \emph{generalized klt singularities} (resp. \emph{generalized lc singularities}) or $(X,B +  { \boldsymbol{\beta}} )$ is said to be \emph{gklt} (resp. \emph{glc}) if %$K_X+B$ is a $\mathbb{Q}$-Cartier divisor and 
the generalized discrepancy of any prime divisors $E$ over $X$ is $a(X,B ,  { \boldsymbol{\beta}} ,E)>-1$ (resp. $a(X,B, { \boldsymbol{\beta}};E)\ge -1$).

A generalized lc pair $(X,B +  { \boldsymbol{\beta}} )$ is said to have \emph{generalized dlt singularities} or $(X,B +  { \boldsymbol{\beta}} )$ is said to be \emph{gdlt} if there is an open subset $U \subset X$ such that $\left(U,\left.(B+\boldsymbol{\beta})\right|_U\right)$ is a log resolution (of itself) and $-1 \leq a(P, X, B+\boldsymbol{\beta}) \leq 0$ for any prime divisor $P$ on $U$ and $-1<a(P, X, B+\boldsymbol{\beta}) \leq 0$ for any prime divisor $P$ over $X$ with center contained in $X \backslash U$.

When ${\boldsymbol{\beta}}$ descends to $X$ (i.e., $\boldsymbol{\beta}_X$ is nef and $\bbeta =\overline {\boldsymbol{\beta}_X} $), the term $ { \boldsymbol{\beta}}$ in the generalized pair $(X,B + \boldsymbol{\beta})$ does not contribute to the singularities. 
\begin{proposition}[{\cite[Lemma 2.13]{DHY23}}]
Let $(X,B+  { \boldsymbol{\beta}})$ be a generalized pair. Suppose that $ { \boldsymbol{\beta}}$ descends to $X$ and $K_X+B$ is $\mathbb{Q}$-Cartier, then $(X,B)$ is klt (resp. dlt, lc) if and only if $(X,B+ { \boldsymbol{\beta}})$ is gklt (resp. gdlt, glc).
\end{proposition}
\begin{proof}
The `if part' has been proved in \cite[Lemma 2.13]{DHY23}. The converse direction follows by 
\cite[Remark 2.6.(a)]{DHY23} combined with the Definition \ref{logdiskrep}.
\end{proof}

When $ { \boldsymbol{\beta}}$ does not descend to $X$, it will contribute to the singularities (cf. \cite[Remark 4.2]{BZ16}).
% \begin{example}[{}]
% Assume $X=\mathbb{P}^2$ and that $$\phi: X' = \text{Bl}_x \mathbb{P}^2 \to X = \mathbb{P}^2$$is the blow up of a point $x$. Let $E^{\prime}$ be the exceptional divisor, $L$ a line passing through $x$ and $L^{\prime}$ the birational transform of $L$. We then consider the following two cases.\\
% (1) Let $M^{\prime}=2 L^{\prime}$ be a divisor on $X'$, and $B= 0$ $M = \phi_* M' = 2L$ be divisors on $X$. By \cite[Example 1.4.6]{Laza04}, $L'^2 = (\phi^*(L)-E')^2 = \deg L + E'^2 = 1-1 = 0 \ge 0$ so that it is a nef divisor, hence $\mathbf M=\overline{M'}$.


% and $$\phi:(X', E' + 2L') \to (X, 0+ 2L).$$B Since $$\phi^*(K_{X}+ 0 + 2L) = K_{X'} -E'+ 2L' + 2E'= K_{X'} + E' + M',$$ then by Definition \ref{logdiskrep},  the generalized pair $(X, B+ \mathbf{M})$ is glc but not generalized klt while the ordinary log pair $(X,B)$ is klt.\\
% (2) If $B=L$ and $M^{\prime}=2 L^{\prime}$ with $M = \phi_*(M') = 2L$,  and$$\phi: (X',(L'+2E')+ 2L')\to (X,L + 2L),$$where $M' = 2L'$ is a nef divisor on $X'$ and $\mathbf M=\overline{M'}$. Since $$\phi^*(K_X+L + 2L) = K_{X'} - E' +3L' +3E' = K_{X'} + (L'+2E') +M'.$$By definition $(X, B+\mathbf{M})$ is not glc, while the ordinary pair $(X,B)$ is lc. Note that in both cases, the b-divisor $\mathbf{M}$ does not descend to $X$, since $\phi^*(M) = 2L + 2E \ne M'$.
% \end{example}

 
\subsection{Nef, big and pseudo-effective classes} In this subsection, we will introduce nef, big and pseudo-effective Bott--Chern cohomology classes on a normal compact complex analytic variety and their basic properties. 

\begin{definition}[Metrically nef, quasi-nef and algebraically nef classes]
Let $X$ be a normal compact complex analytic variety in the Fujiki class $\mathcal{C}$, and let $\alpha = [\theta] \in H_{\mathrm{BC}}^{1,1}(X,\mathbb{R})$ be a class represented by a form $\theta$ with local potentials. Then \\
\emph{(1)} We say that the class $\alpha$ is \emph{(metrically) nef} if for some positive smooth $(1,1)$-form $\omega$ on $X$ and for every $\varepsilon>0$ there exists $f_\varepsilon \in \mathcal{A}^0(X)$ such that $$ \theta+i \partial \overline{\partial} f_\varepsilon \ge-\varepsilon \omega .$$
\emph{(2)} We say that the class $\alpha$ is \emph{quasi-nef} if there exists a K\"ahler modification $\nu \colon \widehat{X} \to X$ such that $\nu^*\alpha$ is a metrically nef class on $\widehat{X}$.\\
\emph{(3)} We say that the class $\alpha$ is \emph{algebraically nef} if for all irreducible curves $C \subset X$, we have  $\alpha\cdot C  \ge 0$.\\
\emph{(4)} Let $f:X\to S$ be a proper surjective morphism of normal relatively compact complex analytic varieties, and  $\alpha \in H^{1,1}_{\rm{BC}}(X,\mathbb{R})$ a Bott--Chern (1,1)-class on $X$. We say that $\alpha$ is \emph{$f$-metrically nef (resp. $f$-algebraically nef)} if $\alpha|_{X_s}$ is metrically nef (resp. algebraically nef) for any $s\in S$.
\end{definition}

Using the cone theorem for K\"ahler varieties \cite{CH20, HP24}, we first establish the equivalence of several notions of nefness for adjoint classes on K\"ahler varieties (cf. \cite[\S\ 2]{Naka86}).


\begin{proposition}\label{nef}
Let $X$ be a normal compact K\"ahler variety, such that $(X, B + \boldsymbol{\beta}_X)$ is a generalized klt pair, and $\alpha \in H^{1,1}_{\mathrm{BC}}(X, \mathbb{R})$. Then the class $\alpha$ is metrically nef if and only if it is quasi-nef. Moreover, if we further assume that $\alpha = K_X + B + \boldsymbol{\beta}_X$, then $\alpha$ is algebraically nef if and only if it is metrically nef.
\end{proposition}
    
\begin{proof}
By definition, $\alpha$ is quasi-nef if and only if $f^* \alpha$ is metrically nef on a resolution $f: \widehat{X} \to X$. On the other hand, by \cite[Theorem 1]{Paunnef},  $\alpha$ is metrically nef if and only if $f^* \alpha$ is metrically nef. Hence clearly $\alpha$ is metrically nef if and only if it is quasi-nef. 

Assume now that $\alpha = K_X + B + \bbeta_X$. By \cite[Theorem 0.5]{HP24} (or alternatively \cite[Theorem 1.3]{CH20}), together with Theorem~\ref{Theorem Ou25}, if $K_X + B + \bbeta_X$ is not metrically nef, then there exists a $(K_X + B + \bbeta_X)$-negative rational curve $C$ such that $
(K_X + B + \bbeta_X)\cdot C < 0$. 
This implies that $K_X + B + \bbeta_X$ is not algebraically nef. 
\end{proof}

\begin{definition}[Pseudo-effective and big classes]
On a normal compact complex analytic variety $X$, we say that:\\  
\emph{(1)} A cohomology class $\alpha \in H^{1,1}_{\rm{BC}}(X,\mathbb{R})$ is \emph{pseudo-effective} if there exists a closed positive (1,1)-current $T\ge 0$, such that $T$ represents the class $\alpha$.\\
\emph{(2)} A cohomology class $\alpha \in H^{1,1}_{\rm{BC}}(X,\mathbb{R})$ is \emph{big} if there exists a closed positive (1,1)-current $T$ to represent the class $\alpha$,  such that $T\ge \delta \omega$ for some smooth Hermitian metric $\omega$ and some number $\delta >0$. 
\end{definition}

If we perturb a pseudo-effective class in the K\"ahler direction, we obtain a big class. 
\begin{proposition}\label{bigandpseff}
Let $X$ be a normal compact K\"ahler variety. A class $\alpha \in H^{1,1}_{\rm{BC}}(X,\mathbb{R})$ is pseudo-effective if and only if there exists a K\"ahler class $\omega$ such that $\alpha + \varepsilon \omega$ is big for any $\varepsilon >0$. 
\end{proposition}
% \begin{proof}
% We first take a resolution $\pi:\widetilde{X} \to X$ with $\widetilde{X}$ be a compact K\"ahler manifold with a K\"ahler class $\omega_{\widetilde{X}}$. Suppose that $\alpha\in H^{1,1}_{\rm{BC}}(X,\mathbb{R})$ is pseudo-effective, then there exists a closed positive $(1,1)$-current $T \in \alpha$. For any $\varepsilon>0$ the current
% $$
% T_{\varepsilon}:=T+\varepsilon \omega
% $$
% is closed and satisfies $T_{\varepsilon} \ge \varepsilon \omega$. Hence $T_{\varepsilon}$ is a K\"ahler current in the class $\alpha+\varepsilon[\omega]$. Therefore $\alpha+ \varepsilon[\omega]$ is big. 

% Conversely, assume that there exists a smooth Hermitian form $\omega$ such that $\alpha + \varepsilon [\omega]$ is big. By Proposition \ref{PushPSEF}, the pull back class $\pi^*(\alpha + \varepsilon [\omega])$ is still big. That is for every $\varepsilon>0$ there exists a K\"ahler current $T_{\varepsilon} \in  \pi^*(\alpha+\varepsilon[\omega])$;  In particular each $T_{\varepsilon}$ is a closed positive current. Choose a sequence $\varepsilon_j \searrow 0$ and corresponding K\"ahler currents $T_j:=T_{\varepsilon_j}$. The family $\left\{T_j\right\}$ is a family of closed positive $(1,1)$-currents with uniformly bounded mass on $\widetilde{X}$.


% Hence, by the usual compactness theorem for positive currents, there exists a subsequence (still denoted $T_j$) that converges weakly to a closed positive $(1,1)$-current $T$. (Weak limits of closed currents are closed, and weak limits of positive currents are positive.) Moreover, the cohomology classes satisfy
% $$
% \left[T_j\right]=\pi^*(\alpha+\varepsilon_j[\omega]) \longrightarrow \pi^*\alpha
% $$
% in $H_{\mathrm{BC}}^{1,1}(\widetilde{X},\mathbb{R})$ so the limit current $T$ lies in the class $\pi^*\alpha$. Therefore $T$ is a closed positive current in $\pi^*\alpha$. So that by Proposition \ref{PushPSEF}, $\alpha$ is pseudo-effective. 
% \end{proof}

\subsection{Lelong numbers and restricted non-K\"ahler locus}
Let us introduce the definitions of Lelong numbers and restricted non-K\"ahler locus on normal complex analytic varieties. 
\begin{definition}[{Lelong numbers, \cite[Définition 2]{Demailly82}}]
Let X be a normal complex $n$-dimensional analytic variety, and ${T} \geq 0$ a closed positive $(1,1)$-current on X. Let ${x} \in {X}$ be an arbitrary point. Consider an embedding
$$
({X}, {x}) \hookrightarrow\left(\mathbb{C}^{{N}}, 0\right).
$$
The coordinate functions $\left({z}_{{i}}\right)_{{i}=1,\cdots, {N}}$ on $\mathbb{C}^{{N}}$ restricted to X induce a generating system $\left({g}_{{i}}\right)_{{i}=1,\cdots, {N}}$ of the maximal ideal of the ring $\mathcal{O}_{{X}, {x}}$. The function $\varphi:=\sum_{{i}}\left|{g}_{{i}}\right|$ is defined on some small open subset $U$ containing $x$. We then define the \emph{Lelong number} of $T$ at $x$ to be
$$
\nu({T},{x}):=\lim _{{r} \rightarrow 0} \frac{1}{(2 \pi {r})^{2 {n}-2}} \int_{\varphi<{r}} {~T} \wedge(\sqrt{-1} \partial \bar{\partial} \varphi)^{{n}-1}.$$
\end{definition}

\begin{definition}[{Restricted non-K\"ahler locus, \cite[Definition 4.18]{HP24}}]
Let $X$ be a normal complex analytic variety and $T$ a closed positive $(1,1)$-current on $X$ with local potentials. For each $c>0$ we define the \emph{Lelong level set} as 
$$
E_c:=\{x \in X: \nu(T, x)\ge c\},
$$
where $\nu(T, x)$ denotes the Lelong number of $T$ at $x \in X$.
This is an analytic subset of $X$. We define
$$
E_{+}(T):=\bigcup_{c \in \mathbb{R}^{+}} E_c .
$$
The \emph{restricted non-K\"ahler locus} of a class $\alpha\in H^{1,1}_{\rm{BC}}(X,\R)$ is the set
$$
E_{n K}^{a s}(\alpha):=\bigcap_{T \in \alpha} E_{+}(T),
$$
where $T$ above is assumed to be a K\"ahler current with weak analytic singularities, and $E_{+}(T) \subset X$ is the (analytic) subset of $X$ for which the Lelong numbers of $T$ are strictly positive. 
\end{definition}

\begin{lemma}\label{nonKahlerlocus}
Let $\alpha \in H^{1,1}_{\rm{BC}}(X,\mathbb{R})$ on a normal compact complex analytic variety $X$. Then $E_{nK}^{as}(\alpha) = \emptyset $ if and only if $\alpha$ is K\"ahler class; $E_{nK}^{as}(\alpha) = X$ if and only if $\alpha$ is not a big class.
\end{lemma}

% Restricted non-K\"ahler locus satisfies the following restriction property. 
% \begin{lemma}

% \end{lemma}
% \begin{proof}

% \end{proof}

\subsection{Volume of cohomology classes}
Compared with the projective setting, a complex analytic variety may admit only a few divisors but many cohomology classes in $H_{\rm{BC}}^{1,1}(X,\mathbb{R})$. While the notion of a global section does not apply to these classes, we can still define their volume. In this subsection, we briefly recall some basic facts about the volume of a cohomology class. 

\begin{definition}[Volume of a cohomology class]\label{Def:Vol}
For $\alpha \in H^{1,1}_{\rm{BC}}(X, \mathbb{R})$ on a normal compact K\"ahler variety $X$, if $\alpha$ is a pseudo-effective cohomology class, we define the \emph{volume} of $\alpha$ by pulling back via a resolution $\pi: X'\to X$ 
$$
\mathrm{vol}  (\alpha):=\sup _{T\in \pi^*(\alpha)} \int_{X'} T_{\rm{ac}}^n =  \sup _{T\in \pi^*(\alpha) }\int_{X' \setminus \operatorname{Sing}(T)}T^n,
$$
for $T$ ranging over all the closed positive (1,1)-currents in $\pi^*(\alpha)$; otherwise, we set $\mathrm{vol}(\alpha)=0$. Here $T_{\mathrm{ac}}$ denotes the absolutely continuous part of the current $T$, and $\operatorname{Sing}(T)$ denotes the singular locus of $T$.
\end{definition}

\begin{remark}
Since the volume is a birational invariant for smooth varieties (cf. \cite[Example 2.2.49]{Laza04} and \cite{Boucksom02}), the volume of a cohomology class on a normal compact K\"ahler variety is independent of the choice of resolution.
\end{remark}
\begin{remark}
Given a pseudo-effective line bundle $L$ on a compact K\"ahler manifold $X$, \cite[Theorem 1.2]{Boucksom02} showed that the usual definition of the volume of a line bundle,
$$
\vol(L) := \limsup_{k \to +\infty} \frac{n!}{k^n} h^0(X, kL),
$$
coincides with Definition~\ref{Def:Vol}.
\end{remark}
The volume function is continuous on the Bott–Chern cohomology space of a normal compact K\"ahler variety.
\begin{proposition}[{\cite[Corollary 4.11]{Boucksom02}}]\label{ContVol}
Let $X$ be a normal compact K\"ahler variety. Then the volume function $$\vol : H^{1,1}_{\rm{BC}}(X,\mathbb{R})\to \mathbb{R}$$is continuous. 
\end{proposition}

For a smooth family of K\"ahler manifolds, the volume function is upper semi-continuous.
\begin{proposition}[{\cite[Proposition 4.13]{Boucksom02}}]\label{uscvolume}
Let ${X} \rightarrow S$ be a smooth family of K\"ahler manifolds. If $\alpha \in H^{1,1}_{\rm{BC}}(X_0, \mathbb{R})$ is a pseudo-effective class on the central fiber $X_0$, then one has
$$
\mathrm{vol}(\alpha) \ge \limsup _{k \rightarrow+\infty} \mathrm{vol}\left(\alpha_k\right)
$$
for every sequence $\alpha_k \in H^{1,1}_{\rm{BC}}\left(X_{t_k}, \mathbb{R}\right)$ converging to $\alpha$ with $(t_k \rightarrow 0)$.
\end{proposition}
\begin{remark}
Recently, \cite{Jiao} proved the upper semi-continuity of volume of divisors for a projective morphism when fibers are irreducible and reduced, by use of a Fujita approximation technique. 
\end{remark}
Volume function satisfies some special properties for big classes and nef classes.

\begin{lemma}[{\cite[Theorems 4.1, 4.7]{Boucksom02}}]\label{BigVol}
For a normal compact K\"ahler variety $X$, one has: \\
(a) A class $\alpha \in H^{1,1}_{\mathrm{BC}}(X, \mathbb{R})$ is big if and only if $\mathrm{vol}(\alpha)>0$;\\
(b) For a nef class $\beta\in H^{1,1}_{\mathrm{BC}}(X,\mathbb{R})$, $\mathrm{vol}(\beta) =\beta^n.$
\end{lemma}

Given a smooth family, the top intersection number of a cohomology class is constant under deformations.
\begin{lemma}\label{constIntersection}
Let $f: X \to S$ be a smooth family of compact complex manifolds over a smooth connected curve $S$, with $n = \dim_{\mathbb{C}}(X/S)$. Suppose that $\alpha \in H^{1,1}_{\mathrm{BC}}(X, \mathbb{R})$. Then the top intersection number
$$
s \longmapsto \int_{X_s} (\alpha|_{X_s})^n
$$
is constant in the family.
\end{lemma}
\begin{proof}
Since $S$ is connected, given any two points $s_1, s_2 \in S$, there exists a smooth real curve $\gamma : [0,1] \to S$ with $\gamma(0) = s_1$, $\gamma(1) = s_2$. Since $f:X\to S$ is a smooth morphism, the constant rank theorem implies that the inverse image $\overline{M} = f^{-1}(\gamma([0,1]))$ is a smooth manifold with boundary $\partial \overline{M} = X_{s_1} \cup X_{s_2}$. Thus, Stokes' theorem gives rise to  
$$
\int_{X_{s_1}}  (\alpha|_{X_{s_1}})^n - \int_{X_{s_2}} (\alpha|_{X_{s_2}})^n = \int_{\overline{M}} d(\alpha^n) = 0.
$$
That is, $s \longmapsto \int_{X_s} (\alpha|_{X_s})^n$ is constant in the family.
\end{proof}

%When the family is non-smooth, a singular Stokes' formula in Proposition \ref{stokes} will be useful in Theorem \ref{InvarVol}.

Given a pseudo-effective class on a compact complex manifold or more generally a normal compact complex analytic variety, we can take its Boucksom--Zariski decomposition. 
% \begin{definition}[{Lelong numbers on normal complex spaces,  and \cite[Definition 4.4]{HP24}}]
% Let $X$ be a normal complex space, and $T \geq 0$ a closed positive $(1,1)$-current on $X$. Let $x \in X$ be an arbitrary point. Consider an embedding
% $$
% (X, x) \hookrightarrow\left(\mathbb{C}^N, 0\right)
% $$
% and then the coordinate functions $\left(z_i\right)_{i=1, \ldots, N}$ on $\mathbb{C}^N$ restricted to $X$ induce a generating system $\left(g_i\right)_{i=1, \ldots, N}$ of the maximal ideal of the ring $\mathcal{O}_{X, x}$. The function
% $$
% \varphi_x:=\sum_i\left|g_i\right|
% $$
% is defined on some small open subset $U$ containing $x$. We then define the \emph{Lelong number} of $T$ at $x$ to be $$\nu(T, x):=\lim _{r \rightarrow 0} \frac{1}{(2 \pi r)^{2 n-2}} \int_{\varphi<r} T \wedge(\sqrt{-1} \partial \bar{\partial} \varphi)^{n-1}.$$
% \end{definition}
\begin{definition}[Minimal multiplicity]
Let $X$ be a compact complex manifold and $\alpha \in H^{1,1}_{\mathrm{BC}}(X, \mathbb{R})$ a pseudo-effective class. The \emph{minimal multiplicity} of $\alpha$ at a point $x \in X$ is defined by
\[
\nu(\alpha, x) := \sup_{\varepsilon > 0} \nu\!\left(T_{\min,\varepsilon}, x\right),
\]
where $T_{\min,\varepsilon}$ denotes a current with minimal singularities in the set
\[
\alpha[-\varepsilon \omega] := \{\text{closed almost positive } (1,1)\text{-currents } T \mid T \text{ represents } \alpha - \varepsilon \omega,\ T \ge -\varepsilon \omega \,\},
\]and $\nu(T_{\min, \varepsilon},x)$ is the Lelong number of $T_{\min,\varepsilon}$ at $x$. 
If $D$ is a prime divisor on $X$, the \emph{generic minimal multiplicity} of $\alpha$ along $D$ is defined as
\[
\nu(\alpha, D) := \inf \{\, \nu(\alpha, x) \mid x \in D\}.
\]
\end{definition}

\begin{definition}[{Boucksom--Zariski decomposition, \cite[Definition 3.7]{Boucksom04}}]
Let $X$ be a compact complex manifold and $\alpha\in H^{1,1}_{\rm{BC}}(X,\mathbb{R})$ a pseudo-effective class. Then we define the \emph{negative part} of $\alpha$ as $$N(\alpha) = \sum  \nu (\alpha , D)\cdot D,$$
where $D$ runs through all the divisors on $X$,  
and $Z(\alpha) = \alpha  - [N(\alpha)].$ 
Then we say that the decomposition 
$$\alpha = Z(\alpha)+ [N(\alpha)]$$ 
is the \emph{Boucksom--Zariski decomposition}. 
\end{definition}
\begin{remark}\label{NegZariski}
Given a normal compact complex analytic variety $X$, let $f \colon Y \to X$ be a resolution of singularities of $X$. For a pseudo-effective class $\alpha\in H^{1,1}_{\rm{BC}}(X,\mathbb{R})$, we define the \emph{negative part} of $\alpha$ by
$$
N(\alpha) := f_*\bigl(N(f^*\alpha)\bigr).
$$
By \cite[Lemma A.7]{DHY23}, this definition is independent of the choice of resolution, and hence $N(\alpha)$ on a normal complex analytic variety is well defined.
\end{remark}
We then describe the behavior of the volume of a big class under the Boucksom--Zariski decomposition:
\begin{proposition}[{\cite[Proposition 3.20]{Boucksom04}}]\label{VolZariski}
Let $X$ be a compact K\"ahler manifold. 
If $\alpha \in H^{1,1}_{\mathrm{BC}}(X, \mathbb{R})$ is big and $E \ge 0$ is a divisor on $X$ such that 
$\mathrm{supp}(E) \subset \mathrm{supp}(N(\alpha))$, then
\[
\mathrm{vol}(\alpha) = \mathrm{vol}(Z(\alpha) + [E]) = \mathrm{vol}(Z(\alpha)),
\]
where $N(\alpha)$ is the negative part of the Boucksom--Zariski decomposition of $\alpha$, 
and $Z(\alpha) := \alpha - [N(\alpha)]$ is the positive part.
\end{proposition}
% \begin{proof}
% Note that Boucksom's result is for compact K\"ahler manifold. However, this can be generalize as follows. Let $X$ be the normal compact K\"ahler variety, and $\pi:X'\to X$ be a resolution. By \cite[Proposition 3.20]{Boucksom04}, we have $$\vol(\pi^*(\alpha)) = \vol (\pi^*(\alpha) -[N(\pi^*(\alpha))]).$$On the other hand, by definition of volume and Remark \ref{NegZariski}, we then have $$\vol(\alpha): = \vol(\pi^*(\alpha)) = \vol (\pi^*(\alpha) -[N(\pi^*(\alpha))]) = \vol(\pi^*(\alpha - [N(\alpha)])) = \vol(\alpha-[N(\alpha)]).$$
% \end{proof}
Let us end this subsection with some volume inequalities.
\begin{lemma}\label{vollemma}
Let $X$ be a compact K\"ahler manifold of dimension $n$ and $\omega \in H^{1,1}_{\mathrm{BC}}(X, \mathbb{R})$ a K\"ahler class on $X$. If $\alpha \in H^{1,1}_{\mathrm{BC}}(X, \mathbb{R})$ is a pseudo-effective class, then for any $\varepsilon >0$, the volume satisfies the inequality $$\vol (\alpha + \varepsilon \omega) \ge \varepsilon^n  \int_X\omega^n.$$
\end{lemma}
\begin{proof}
Let $T \in \alpha$ be a closed positive $(1,1)$-current on $X$. 
By Demailly's regularization theorem (\cite[Proposition (18.7)]{Demailly12}), 
there exists a sequence $\{T_k\}_{k \in \mathbb{N}}$ of closed $(1,1)$-currents 
with analytic singularities in the same cohomology class $\alpha$, weakly converging to $T$, such that 
\[
T_k + \varepsilon_k \omega \ge 0, \qquad \varepsilon_k \downarrow 0.
\]


In particular, for large $k$ we have $\varepsilon_k<\varepsilon$, so $T_k + \varepsilon\omega \ge (\varepsilon-\varepsilon_k)\omega$. Set $\delta_k=\varepsilon-\varepsilon_k>0$ and then $\delta_k \to \varepsilon$. By the definition of volume, we have $$\vol(\alpha + \varepsilon \omega) = \sup_{T_{\varepsilon} \in \alpha + \varepsilon \omega}\int _X (T_{\varepsilon})_{ac}^n \ge \limsup_{k\to \infty} \int_X (T_k + \varepsilon \omega)^n_{ac} \ge  (\limsup_{k\to \infty} \delta_k^n) \int_X\omega^n = \varepsilon^n \int_X\omega^n.$$
\end{proof}

The volume function satisfies certain log concavity properties.
\begin{proposition}\label{r-logc}
Let $X$ be a compact K\"ahler manifold of complex dimension $n$. If $\alpha, \alpha ' \in H^{1,1}_{\rm{BC}}(X, \mathbb{R})$ are big classes, then the volume satisfies the log concavity property $$\vol(\alpha + \alpha ')^{1/n} \ge \vol (\alpha)^{1/n} + \vol (\alpha ')^{1/n}.$$
\end{proposition}
\begin{proof}
Note that the log concavity of the volume is well known for line bundles and $\R$-divisors (cf. \cite[Remark 2.2.50]{Laza04} and also \cite[Corollary 4.12]{LM09}) and also for nef classes by the Khovanskii–-Teissier inequalities proved in \cite{Dem93}. The general case then follows easily from the analytic version of Fujita’s
approximation theorem  \cite[Theorem 1.4]{Boucksom02}. Given big classes $\alpha,\alpha' \in H^{1,1}_{\rm{BC}}(X, \mathbb{R})$ and $\varepsilon >0$, there exists a resolution $\nu : X_\varepsilon \to X$ and K\"ahler classes $\omega ,\omega ' \in H^{1,1}_{\rm{BC}}(X_{\varepsilon}, \mathbb{R})$ on $X_\varepsilon$ such that
$${\rm vol}(\omega )\leq {\rm vol}(\alpha)\leq {\rm vol}(\omega )+\varepsilon ,$$ $${\rm vol}(\omega') \leq {\rm vol}(\alpha ')\leq {\rm vol}(\omega ') +\varepsilon,$$ $${\rm vol}(\omega +\omega')\leq{\rm vol}(\alpha+\alpha ').$$
But then
$$({\rm vol}(\alpha )-\varepsilon)^{1/n}+({\rm vol}(\alpha')-\varepsilon)^{1/n}\le {\rm vol}(\omega )^{1/n} +{\rm vol}(\omega')^{1/n}\le {\rm vol}(\omega +\omega')^{1/n}\le {\rm vol}(\alpha+\alpha')^{1/n},$$
and taking the limit as $\varepsilon \to 0$ yields ${\rm vol}(\alpha)^{1/n}+{\rm vol}(\alpha')^{1/n} \le {\rm vol}(\alpha+\alpha')^{1/n}$.
\end{proof}

\subsection{Bigness and pseudo-effectivity in families} 
We start this subsection with: 
\begin{definition}[Relatively big and pseudo-effective classes]\label{pseffclass}
Let $f:X\to S$ be a proper surjective morphism from a normal complex analytic variety $X$ onto a relatively compact base $S$. Then we say that:\\
\emph{(1)} A cohomology class $\alpha\in H^{1,1}_{\rm{BC}}(X,\mathbb{R})$ is \emph{$f$-pseudo-effective over $S$} or \emph{relatively pseudo-effective over $S$} if there is a complement $V$ of a countable union of Zariski closed proper subsets of $S$ such that $\alpha_s:=\alpha|_{X_s}$ is pseudo-effective for any $s\in V$.\\
\emph{(2)} A cohomology class $\alpha\in H^{1,1}_{\rm{BC}}(X,\mathbb{R})$ is \emph{$f$-big over $S$} or \emph{relatively big over $S$} if there is a non-empty Zariski open subset $U$ of $S$ such that $\alpha_s$ is big for any $s\in U$.  
\end{definition}

The following proposition shows that a relatively pseudo-effective class is pseudo-effective on general fibers.
\begin{proposition}\label{genericpseff}
Let $f:X\to S$ be a family from a normal K\"ahler variety $X$ with a K\"ahler class $\omega$ onto a relatively compact base $S$. A class $\alpha \in H^{1,1}_{\rm{BC}}(X, \mathbb{R})$ is $f$-pseudo-effective over $S$ if and only if there exists a non-empty Zariski open subset $U \subset S$ such that $\alpha_s$ is pseudo-effective for any $s\in U$. 
\end{proposition}
\begin{proof}
After shrinking $S$ and replacing $X$ by a resolution,  we may assume that $X\to S$ is smooth. By assumption $\alpha _s$ is pseudo-effective for any very general $s\in S$.
By Lemma \ref{vollemma}, $$\vol(\alpha _s+\varepsilon \omega _s)  \ge \varepsilon ^d \cdot v \text{ for very general } s\in S,$$ where $d=\dim (X/S)$ and $v=\int _{X_s}\omega _s^d$. Note that $v= \int_{X_s} \omega_s^d>0$ is a constant independent of $s\in S$ by Lemma \ref{constIntersection}. 

Let $X_{s_0}$ be a special fiber, and define the pseudo-effective threshold
\[
\tau_{s_0} = \inf \{ t \ge 0 \mid \alpha_{s_0} + t \omega_{s_0} \text{ is pseudo-effective} \}.
\]
Then $\alpha_{s_0} + \tau_{s_0} \omega_{s_0}$ is pseudo-effective but not big; hence $
\vol(\alpha_{s_0} + \tau_{s_0} \omega_{s_0}) = 0$. 
By Boucksom's upper semicontinuity Proposition~\ref{uscvolume},
\[
0 = \vol(\alpha_{s_0} + \tau_{s_0} \omega_{s_0})
\ge \limsup_{k \to \infty} \vol(\alpha_{s_k} + \tau_{s_0} \omega_{s_k})
\ge \tau_{s_0}^d \, v.
\]
Therefore $\tau_{s_0} = 0$, and in particular $\alpha_{s_0}$ is pseudo-effective.
\end{proof}

Before proceeding, let us first recall the definition of the pullback of a closed positive $(1,1)$-current with local potentials. 

\begin{definition}[{Pull back of closed positive $(1,1)$-current with local potential, \cite[Section 3.1]{Meo96}}]\label{Def: pullcurrent}
Let $f: X \rightarrow Y$ be a  surjective morphism between complex analytic varieties. For a closed positive $(1,1)$-current $T$ with local potentials on $Y$, we define the pull back $f^* T$ as follows:

Let $\Omega$ be an open set in $Y$ on which $T_{\Omega}=i \partial \bar{\partial} u$ with $u$ a plurisubharmonic potential function in $\Omega$. If the plurisubharmonic function $u \circ f$ is not identically $-\infty$ on each connected component of $f^{-1}(\Omega)$, we then define $$\left.\left(f^* T\right)\right|_{f^{-1}(\Omega)}=i \partial \bar{\partial}(u \circ f).$$If $\Omega^{\prime}$ is another open set in $Y$ with a similar representation $T_{\Omega^{\prime}}=i \partial \bar{\partial} u^{\prime}$, the function $u-u^{\prime}$ is pluriharmonic on $\Omega \cap \Omega^{\prime}$, and the same holds for $u \circ f-u^{\prime} \circ f$ on $f^{-1}(\Omega) \cap f^{-1}\left(\Omega^{\prime}\right)$, so the above definition is independent of the choice of $\Omega$ and the potential $u$.
\end{definition}
\begin{remark}\label{Rmk: Globalpull}
Note that the closed positive current $T$ admits a global representative 
$T = \theta + i \partial \bar{\partial} \varphi$, where $\theta$ is a smooth $(1,1)$-form and $\varphi$ is an almost plurisubharmonic function on $Y$, obtained via a partition of unity argument (cf.~\cite[Proposition 3.1.1]{Meo96}). Hence the pullback $f^*T$ can be written as 
$f^*T = f^*\theta + i \partial \bar{\partial} f^*\varphi$, and therefore $[f^*T] = f^*[T]$.
\end{remark}

We can now describe the behavior of pseudo-effective and big classes under bimeromorphic change.

\begin{proposition}\label{Prop: PullPSEf}
Let $f: Y \rightarrow X/S$ be a proper bimeromorphic morphism of normal complex analytic varieties over a relatively compact base $S$. If $\alpha \in H^{1,1}_{\rm BC}(X,\mathbb{R})$ and ${{\beta}} \in H^{1,1}_{\rm BC}(Y,\mathbb{R})$ are relatively pseudo-effective over $S$, then so are the classes ${{\beta}}' = f_*{{\beta}}$ and $\alpha' = f^*\alpha$.
\end{proposition}
\begin{proof}
First, by Definition \ref{pseffclass}, to check relative pseudo-effectivity, it is sufficient to check it on the general fiber. By \cite[Corollary 3.2]{CRC25}, $f_s : Y_s \to X_s$ is bimeromorphic on the general fiber. Hence, in what follows, without loss of generality, we can reduce the problem to the absolute case and assume that $S$ is a point. 

We then show that if $\alpha \in H^{1,1}_{\rm BC}(X, \mathbb{R})$ is pseudo-effective, then so is $f^*\alpha$. Let $T_\alpha$ be a closed positive $(1,1)$-current such that $[T_\alpha] = \alpha$. By Definition \ref{Def: pullcurrent}, we can find an open cover $\{\Omega_i\}$ of $X$ such that
\[
T_\alpha|_{\Omega_i} = i \partial \bar{\partial} u_i,
\]
for some plurisubharmonic function $u_i$, and
\[
f^*T_\alpha \big|_{f^{-1}(\Omega_i)} = i \partial \bar{\partial} (u_i \circ f).
\]
By \cite[Proposition~1.1.8]{Demailly12}, the pullback of a plurisubharmonic function is again plurisubharmonic, so $f^*T_\alpha$ is a positive current. By Remark \ref{Rmk: Globalpull}, $[f^* T_{\alpha}] = f^* \alpha$. Hence $f^* \alpha$ is a pseudo-effective class.

% $
% T_\alpha = \theta + i \partial \overline{\partial}\varphi$ for some closed smooth form $\theta$ of bidegree (1,1) such that $[T_\alpha ] = [\theta] =\alpha$, with the pull back 
% \[
% f^*T_\alpha = f^* \theta + i \partial \overline{\partial}(f^* \varphi).
% \]
% And therefore $[f^* T_\alpha]  = [f^* \theta] = f^*\alpha$. 

We end our proof by showing that if $\beta \in H^{1,1}_{\rm BC}(X,\mathbb{R})$ is pseudo-effective, then so is $f_* \beta$. First, there exists a closed positive $(1,1)$-current $T_{\beta}$ such that $[T_{\beta}] = \beta$. By definition, a closed $(1,1)$-current is positive if and only if it is positive on every smooth positive $(n-1,n-1)$-form $\varphi$. Moreover, since $\varphi$ is positive, the pullback $f^*\varphi$ is semi-positive. Hence
\[
\langle f_* T_{\beta}, \varphi \rangle = \langle T_{\beta}, f^*\varphi \rangle \ge 0,
\]
for every smooth positive $(n-1,n-1)$-form $\varphi$. Thus, the pushforward of a pseudo-effective class is pseudo-effective.   
\end{proof}
\begin{proposition}\label{Prop: PullBig}
Let $f: Y\rightarrow X/S$ be a proper bimeromorphic morphism of normal K\"ahler varieties over a relatively compact base $S$. If $\alpha \in H^{1,1}_{\rm BC}(X,\mathbb{R})$ and ${{\beta}} \in H^{1,1}_{\rm BC}(Y,\mathbb{R})$ are relatively big over $S$, then so are the classes ${{\beta}}' = f_*{{\beta}}$ and $\alpha' = f^*\alpha$.
\end{proposition}
\begin{proof}
Similar to Proposition~\ref{Prop: PullPSEf}, we may reduce to the absolute setting and assume that $S$ is a point. In particular, both $X$ and $Y$ are compact.

We first show that if $\alpha \in H^{1,1}_{\mathrm{BC}}(X, \mathbb{R})$ is big, then $f^{*}\alpha$ is big. Since
$$
\mathrm{vol}(f^{*}\alpha) = \mathrm{vol}(\alpha) > 0,
$$
Proposition~\ref{BigVol} implies that $f^{*}\alpha$ is big.

We next prove that if $\beta \in H^{1,1}_{\mathrm{BC}}(Y, \mathbb{R})$ is big, then $\beta' = f_{*}\beta$ is big. Let $T_{\beta}$ be a K\"ahler current representing $\beta$. By definition, there exists a smooth Hermitian form $\omega_{Y}$ with
$$
T_{\beta} \ge \omega_{Y}.
$$
Let $\omega_{X}$ be a smooth Hermitian form on $X$. Since $Y$ is compact, there exists a constant $C > 0$ such that
$$
f^{*}\omega_{X} \leq C\, \omega_{Y}.
$$
It follows that
$$
f_{*}T_{\beta} \ge f_{*}\omega_{Y} \ge C^{-1} f_{*}f^{*}\omega_{X}
  = C^{-1}\, \omega_{X}.
$$
Thus $f_{*}T_{\beta}$ is again a K\"ahler current. Because $[f_{*}T_{\beta}] = f_{*}\beta$, the class $f_{*}\beta$ is big.
\end{proof}

We say that a bimeromorphic contraction $ \phi : X \dashrightarrow Y $ is a \emph{$(K_X + B +  { {\bbeta}})$-non-positive contraction} if under a common resolution $ p : W \to X $ and $ q : W \to Y $, the divisor $ E $ in
\[
p^*(K_X + B +  { {\bbeta}_X}) = q^*(K_Y + B_Y +  { {\bbeta}}_Y) + E
\]
is an effective $ q $-exceptional divisor (that is, $ q_* E = 0 $).

A $(K_X+B+ \bbeta_X)$-non-positive contraction preserves pseudo-effectivity (resp. bigness) of adjoint classes. 
\begin{proposition}\label{PushPSEF}
Let $(X,B + \bbeta_X)$ be a generalized K\"ahler pair. Given a $(K_X+B+ \bbeta_X)$-non-positive contraction (e.g. a sequence of steps of the $(K_X+B + \bbeta_X)$-MMP) over the relatively compact $S$
\[
\phi : (X, B +  { \boldsymbol{\beta}_X}) 
\dashrightarrow 
(X', B' +  { \boldsymbol{\beta}}_{X'})/S,
\]
the adjoint class $ K_{X} + B +  { {\bbeta}_X} $ is relatively pseudo-effective (resp.\ big) over $ S $ if and only if $ K_{X'} + B' +  { {\bbeta}}_{X'}$ is relatively pseudo-effective (resp.\ big) over $ S $.
\end{proposition}

\begin{proof}
%Let $(X,B+\bbeta_X)$ be the generalized K\"ahler pair as above, and let $$\phi: (X,B + \bbeta_X) \dashrightarrow (X', B'+ \bbeta_{X'})/S$$be $(K_X+B+ \bbeta_X)$-non-positive contraction. 
Assume that $K_{X}+ B+ {\bbeta}_X$ is relatively pseudo-effective (resp.\ big) over $S$. 
Let $p:Y \to X$ and $q:Y \to X'$ be a common resolution of $\phi$. 
Since $\phi$ is $(K_{X}+B+{\bbeta}_X)$-non-positive over $S$,  $$p^*(K_{X}+ B+{\bbeta}_X) \equiv_S 
q^*(K_{X'}+B'+ {\bbeta}_{X'}) + E,$$
with $E \ge 0$ effective and $q$-exceptional. 
Since we already showed in Propositions \ref{Prop: PullPSEf} and \ref{Prop: PullBig} that pullback and pushforward under bimeromorphic morphisms preserve pseudo-effectivity (resp.\ bigness), it follows that
\[
q_*\bigl(p^*(K_{X}+B + {\bbeta}_X)\bigr) 
\equiv_S 
q_*\bigl(q^*(K_{X'}+B'+ {\bbeta}_{X'}) + E\bigr) 
= K_{X'}+ B'+ {\bbeta}_{X'}
\]
is relatively pseudo-effective (resp.\ big) over $S$. 

Conversely, if $K_{X'}+B'+ {\bbeta}_{X'}$ is relatively pseudo-effective (resp.\ big) over $S$, then by the effectiveness of $E$ and $p^*(K_{X}+ B+{\bbeta}_X) \equiv_S 
q^*(K_{X'}+B'+ {\bbeta}_{X'}) + E$, it is easy to see that 
$p^*(K_{X}+ B+ {\bbeta}_X)$ is relatively pseudo-effective (resp.\ big) over $S$, and therefore so is 
$K_{X}+ B+ {\bbeta}_X$.
\end{proof}

\begin{remark}
Note that for a bimeromorphic contraction morphism, the pseudo-effectivity of $ f_*  \alpha $  does not necessarily imply that of $  { {\alpha}} $. Consider the blow-up $ f : X \to Y $ of a smooth projective variety $ Y $. Let $ D \ge 0 $ be an effective divisor on $ X $ and $ E \ge 0 $ an exceptional divisor. For $ m \gg 1 $, the divisor $ D - mE $ is not pseudo-effective, but
\[
f_*(D - mE) = f_* D \ge 0
\]
is pseudo-effective.
\end{remark}


\begin{proposition}\label{AbsRelativepseff}
Let $f:X\to S$ be a family from a normal K\"ahler variety $X$ to a complex analytic base $S$, and $\alpha \in H^{1,1}_{\rm{BC}}(X,\mathbb{R})$. Let $W$ be a compact subset of $S$ and $X_W = f^{-1}(W)$. If $\alpha|_{X_W} \in H^{1,1}_{\rm{BC}}(X_W,\mathbb{R})$ is absolute big (resp. pseudo-effective), then it is relatively big (resp. pseudo-effective) over $W$. 
\end{proposition}

\begin{proof}
We first prove that if $\alpha|_{X_W}$ is absolute big then it is also relatively big over $W$. By Lemma \ref{nonKahlerlocus}, the restricted non-K\"ahler locus of $\alpha$ is a proper subvariety of $X_W$. 
In particular, there exists a non-empty Zariski open subset $U\subset W$ such that the fibers $X_s$ for $s\in U$ are not contained in $E_{nK}^{as}(\alpha|_{X_W})$. Since $$E_{nK}^{as}(\alpha|_{X_s}) \subset E_{nK}^{as}(\alpha|_{X_W})\cap X_s,$$ $E_{nK}^{as}(\alpha|_{X_s})$ does not contain the fiber $X_s$ for $s\in U$. Hence, by Lemma \ref{nonKahlerlocus}, the class $\alpha|_{X_s}$ is big for $s\in U$, i.e., $\alpha|_{X_W}$ is $f$-big over $W$. 

Given a pseudo-effective class $\alpha|_{X_W}$ on $X_W$, Proposition \ref{bigandpseff} shows that there exists a K\"ahler class $\omega$ on $X_W$ such that $\alpha|_{X_W}  + \varepsilon \omega$ is (absolute) big for any $\varepsilon >0$. Hence, by what we just proved above, $\alpha|_{X_W} + \varepsilon \omega$ is also $f$-big over $W$ for any $\varepsilon >0$. Thus, $\alpha$ is $f$-pseudo-effective over $W$. 
\end{proof}

\begin{proposition}[{\cite[Lemma 5.1]{DH24}}]\label{KVVanishing}
Let $g: X \rightarrow S$ be a proper surjective morphism from a generalized K\"ahler pair $(X, B+ {\boldsymbol{\beta}})$ with generalized klt singularities onto a relatively compact complex analytic variety $S$. If $D$ is a $\mathbb{Q}$-Cartier $\mathbb{Z}$-divisor on $X$ such that $D-\left(K_X+B+ { {\bbeta}}_X\right)$ is relatively nef and big over $S$, then $R^i g_* \mathcal{O}_X(D)=0$ for all $i>0$.
\end{proposition}
\begin{proof}
The proof is based on \cite[Lemma 5.1]{DH24} with a minor change. Let $\nu:(X',B'+{\boldsymbol{\beta}}_{X'})\to (X,B+{\boldsymbol{\beta}}_X)$ be a log resolution such that ${\boldsymbol{\beta}}_{X'}$ is nef over $S$, with 
$$
K_{X'}+B'+\boldsymbol{\beta}_{X'}=\nu^*(K_X+B+\boldsymbol{\beta}_X).
$$
By Propositions \ref{Prop: PullBig} and \ref{nef}, the divisor
$$
 \nu^*(D) - (K_{X'}+B')=\boldsymbol{\beta}_{X'} + \nu^*(D - (K_X+B+\boldsymbol{\beta}_X)) 
$$
is a big and nef over $S$. Since $X'$ is smooth, there exists some $m\in \mathbb{N}_{+}$ such that $m(\nu^*(D) - (K_{X'}+B'))$ is a relatively big line bundle over $S$. By \cite[Lemma 2.18]{DH24}, there exists a projective modification $\mu:X'' \to X'$ such that $g'':X'' \to S$ is projective.

Then the proof follows along the same lines as \cite[Lemma~5.1]{DH24}. 
Let $\psi = \mu \circ \nu$. We have the following commutative diagram:
\begin{center}
\begin{tikzcd}[ampersand replacement=\&]
	{X''} \& {X'} \& X \\
	\&\& S
	\arrow["\mu", from=1-1, to=1-2]
	\arrow["{g''}"', from=1-1, to=2-3]
	\arrow["\nu", from=1-2, to=1-3]
	\arrow["{g'}"{description}, from=1-2, to=2-3]
	\arrow["g", from=1-3, to=2-3]
\end{tikzcd}
\end{center}
such that $\boldsymbol{\beta}_{X''} = \mu^*\boldsymbol{\beta}_{X'}$ is nef over $S$, and
\[
K_{X''} + B'' + \boldsymbol{\beta}_{X''}
= \psi^*\bigl(K_X + B + \boldsymbol{\beta}_X\bigr).
\]
We then set 
$$
E := \left\lfloor B'' + \{-\psi^*D\} \right\rfloor, \quad B^* := B'' + \{-\psi^*D\} - E = \{B'' + \{-\psi^*D\}\},
$$ 
where $\{\cdot\},\left\lfloor\cdot\right\rfloor,\left\lceil\cdot\right\rceil $ denote the fractional, round down, and round up part of a divisor, respectively. Hence, 
$$
\left\lceil\psi^*D\right\rceil - E - (K_{X''}+B^*) \equiv \boldsymbol{\beta}_{X''} + \psi^*\left(D - (K_X+B+\boldsymbol{\beta}_X)\right).
$$
By assumption, $D-(K_X +B + \boldsymbol{\beta}_X)$ is nef and big over $S$, hence $\left\lceil\psi^*D\right\rceil - E - (K_{X''}+B^*)$ is also nef and big over $S$. 

Since $g'':X'' \to S$ is projective, we may apply the standard version of the Kawamata--Viehweg vanishing theorem, obtaining
$$
R^p(g'')_* \mathcal{O}_{X''}(\left\lceil\psi^*D\right\rceil - E) = 0, \quad 
R^p \psi_* \mathcal{O}_{X''}(\left\lceil\psi^*D\right\rceil - E) = 0 \quad \text{for any } p > 0.
$$
By the same argument as in \cite[Lemma 5.1]{DH24}, we can prove that $$
\psi_* \mathcal{O}_{X''}(\left\lceil\psi^*D\right\rceil - E) = \mathcal{O}_X(D).
$$ Therefore, a standard Leray spectral sequence argument yields
$$
R^p g_* \mathcal{O}_X(D) 
= R^p g_*(\psi_* \mathcal{O}_{X''}(\left\lceil\psi^*D\right\rceil - E)) 
= R^p g''_* \mathcal{O}_{X''}(\left\lceil\psi^*D\right\rceil - E) 
= 0
$$
for all $p > 0$.
\end{proof}


\subsection{Uniruled varieties}We next introduce some basic results about uniruled varieties.
\begin{definition}[Uniruled varieties]\label{Def: Uniruled}
A complex analytic variety $X$ is called \emph{uniruled} if for every general point $x\in X$ there is a non-constant morphism $f:  \mathbb{P}^1\to X$ such that $f(0) = x$.
\end{definition}

\begin{remark}\label{charaterization-uniruled}
With a slight modification of the proof of \cite[Proposition II.5.4]{MP97}, if $X$ is a compact complex manifold, then the following are equivalent.\\
(1) $X$ is uniruled (in the sense of Definition \ref{Def: Uniruled});\\
(2) There exists a complex analytic variety $Y$ and a dominant meromorphic map $$\Phi: \mathbb{P}^1 \times Y  \dashrightarrow X,$$that is not constant on $\mathbb{P}^1$;\\
(3) There exists a free rational curve $f: \mathbb{P}^1 \to X$;\\
(4) For a general point $x\in X$, there exists a free rational curve $f: \mathbb{P}^1 \to X$ with $f(\infty) = x$. 

Obviously, uniruledness is a bimeromorphically invariant property. If, moreover, $X$ is K\"ahler, then uniruledness is also equivalent to $K_X$ being non–pseudo-effective by Theorem~\ref{Theorem Ou25}.
\end{remark}


Any uniruled variety with canonical singularities is covered by family of $K_X$-negative curves.
Recall that a normal complex variety $M$ is said to have
\emph{(only) canonical singularities} if the canonical divisor $K_M$
is a {$\mathbb{Q}$-Cartier divisor}, i.e., some integer multiple of
$K_M$ is a Cartier divisor, and if the discrepancy divisor
$K_{\tilde{M}}-\mu^*K_M$ is effective for a (or every) resolution of singularities
$\mu:\tilde{M}\rightarrow M$. 
\begin{lemma}\label{Lemma: NegCurve}
Let $X$ be a uniruled variety with canonical singularities. Then $X$ is covered by a family of rational curves $\{{C}_\lambda\}_{\lambda \in \Lambda}$ such that $K_X\cdot {C}_\lambda <0$. 
\end{lemma}
\begin{proof}
Since $X$ is uniruled, by Definition \ref{Def: Uniruled}, $X$ is covered by rational curves $\{{C}_\lambda\}_{\lambda \in \Lambda}$. Since $X$ admits canonical singularities, let $f:X'\to X$ be a resolution then $$ K_{X'}=f^*K_X+E,$$for some effective, $f$-exceptional divisor $E\ge 0$. Let ${C}_\lambda' = f^{-1}_* ({C}_\lambda)$ be the strict transform of a general member in the family of curves $\{{C}_\lambda\}_{\lambda \in \Lambda}$, then $$K_X\cdot {C}_\lambda=(K_{X'}-E)\cdot {C}'_\lambda\leq K_{X'}\cdot {C}'_\lambda$$ and so it suffices to show that $K_{X'}\cdot {C}'_\lambda<0$; in particular, we may assume that $X$ is smooth.

Then the rest of the proof follows from \cite{Kollar96}. Since $X$ is a smooth uniruled variety, there is a family of rational curves $\mathcal C\to T$ and a dominant holomorphic map $g:\mathcal C \to X$. Replacing $T$ by an open subset, we may assume that $T$ is isomorphic to a disc. Cutting by hyperplanes, we may assume that $\dim \mathcal C=\dim X$ and hence $g$ is a generically finite holomorphic map of complex manifolds. We may then write $K_{\mathcal C}=g^*K_X+R$ where $R$ is an effective divisor. We then have \[-2={\rm deg}(K_{\mathbb P ^1})={\rm deg}(K_{\mathcal C}|_{ C _\lambda})=K_X\cdot  C _\lambda+R\cdot  C _\lambda\ge K_X\cdot C _\lambda\] and hence $K_X\cdot  C _\lambda<0$.
\end{proof}
% \begin{proof}
% Note that $(4) \implies (1)$ and $(4) \implies (3)$ are clear. 
% $(1) \implies (2)$.
% Given $x\in X$, let $C_x$ be a rational curve passing through $x\in X$. Let $\{C_x\}_{x\in X}$ be a uncountable family of rational curves. Then there exists a component of Barlet cyclc space $\Lambda$ that contains them such that $$X = \overline{\bigcup_{C_x \in \Lambda} C_x}.$$Consider the normalization of the universal family over $Z$. Then $$\widetilde{U}\to \widetilde{ Z}$$bimeromorphic to a $\mathbb{P}^1$-bundle and $Y\to \widetilde{Z}$ is bimeromorphic. 

% $(2) \implies (1)$. Clearly $\{C_t\}_{t\in Y}$ for $C_t = \Phi(\mathbb{P}^1,t)$ will be a dominant family of rational curves on $X$.  

% $(3) \implies (4)$. 
% \end{proof}


We can easily generalize the remarkable Theorem \ref{Theorem Ou25} of Ou to the singular setting. 
\begin{lemma}\label{BDPPCanonical}
For a normal compact $\mathbb{Q}$-Gorenstein K\"ahler variety $X$, we have:\\
\emph{(1)} If $X$ is not uniruled, then $K_X$ is pseudo-effective;\\
\emph{(2)} If $X$ admits canonical singularities and $K_X$ is pseudo-effective, then $X$ is not uniruled. 
\end{lemma}
\begin{proof}
(1) Let $\pi: Y \to X$ be a resolution. If $X$ is not uniruled, then neither is $Y$. Since $Y$ is smooth and K\"ahler, Theorem \ref{Theorem Ou25} implies that $K_Y$ is pseudo-effective. By Proposition \ref{Prop: PullPSEf}, $K_X$ is pseudo-effective. 

(2) If $X$ has canonical singularities and $\pi: Y \to X$ is a resolution, then $K_Y=\pi^*K_X  + E$ with the exceptional divisor $E \ge 0$. Thus, if $K_X$ is pseudo-effective, then $K_Y$ is also pseudo-effective. By \cite[Theorem 1.1]{Ou25}, $Y$ is not uniruled and thus, $X$ is not uniruled by bimeromorphic invariance of uniruledness, Remark \ref{charaterization-uniruled}.
\end{proof}

\begin{remark}
If $X$ is log terminal, the pseudo-effectivity of $K_X$ does not necessarily imply that $X$ is not uniruled as shown in the following two counterexamples. As the first one, in \cite[Section 5]{k-bmy}, Koll\'ar constructed a hypersurface $$H\left(a_1, \ldots, a_n\right):=\left(x_1^{a_1} x_2+x_2^{a_2} x_3+\cdots+x_{n-1}^{a_{n-1}} x_n+x_n^{a_n} x_1=0\right) \subset \mathbb{P}\left(w_1, \ldots, w_n\right)$$ with quotient singularities (hence log terminal), which is rational and $K_{H(a_1,\ldots ,a_n)}$ is ample for a certain choice of $(a_1, \ldots , a_n)$. The second one was given by Totaro in \cite[Example on p. 245]{Totaro10}, and is a rational Calabi--Yau quotient surface as $Y  = E \times E / (\mathbb{Z}/3)$ where $\mathbb{Z}/3$ acts by $(\zeta ,\zeta )$ on $E\times E $ with $\zeta$ a third root of the unity. Note that $Y$ is rational, since it can be obtained by blowing up 12 points on $\mathbb{P}^2$ in the Hesse configuration and then blowing down 9 lines.
\end{remark}


\subsection{Pseudo-effective threshold}
We end this section by introducing the pseudo-effective threshold and several basic facts about it, which will play an important role in the proof of our main Theorem \ref{Deformation invariance of pseudo-effectivity}.

\begin{definition}[Pseudo-effective threshold]
Let $X$ be a normal compact K\"ahler variety, and $\omega \in H^{1,1}_{\rm{BC}}(X, \mathbb{R})$ a K\"ahler class on $X$. We define the \emph{pseudo-effective threshold of} $\alpha \in H^{1,1}_{\rm{BC}}(X,\mathbb{R})$ to be $$\tau(X, \alpha; \omega):=\inf \left\{t \in \mathbb{R}_{\ge 0} \mid \alpha+t \omega \text { is pseudo-effective}\right\}.$$
\end{definition}
Similarly, we can define the relative pseudo-effective threshold. 
\begin{definition}[Relative pseudo-effective threshold]
Let $f:X\to S$ be a proper surjective morphism with connected fibers from a normal K\"ahler variety onto a relatively compact base $S$,  %$X$ onto a smooth connected curve $S$ 
and $\omega \in H^{1,1}_{\rm{BC}}(X, \mathbb{R})$ a K\"ahler class on $X$. 
We define the \emph{relative pseudo-effective threshold} of $\alpha\in H^{1,1}_{\rm{BC}}(X,\mathbb{R})$ over $S$ to be $$\tau(X/S,\alpha ;\omega) := \inf\{t\in \mathbb{R}_{\ge 0}\mid \alpha + t \omega \text{ is }f\text{-pseudo-effective over }S\}.$$
For a line bundle $L$ on $X$, set $\tau(X/S, L ; \omega) :=  \tau(X/S, c_1(L) ; \omega)$. 

Moreover, we define
\[
\tau'(X/S, \alpha; \omega) := \inf \{ t \in \mathbb{R}_{\ge 0} \mid \alpha + t \omega \text{ is } f\text{-big over } S \}.
\] 
\end{definition}

\begin{lemma}
Let $f:X\to S$ be a proper surjective morphism with connected fibers from a normal K\"ahler variety onto a relatively compact base $S$, and $\omega\in H^{1,1}_{\rm{BC}}(X, \R)$ a K\"ahler class on $X$. Then  $$\tau(X/S, \alpha; \omega)=\tau'(X/S, \alpha; \omega).$$
\end{lemma}
\begin{proof} Set $\tau=\tau(X/S, \alpha; \omega)$ and $\tau'= \tau'(X/S, \alpha; \omega)$. We only need to prove $\tau' \le \tau$. By definition of $\tau$, there exists a non increasing sequence $t_i$ with $\lim_{i \to \infty} t_i =  \tau$ in $\mathbb{R}_{\ge 0}$ and a countable union $Z_i$ of proper Zariski closed subsets of $S$ such that $\alpha_s + t_i \omega_s$ is pseudo-effective for every $s\in V_i : = S \setminus Z_i$. Therefore $\alpha_s + \tau \omega_s$ is pseudo-effective for every $s \in V_\infty = \bigcap_{i =1}^\infty V_i$ (note that $V_\infty$ is also a complement of countable union of proper Zariski closed subsets of $S$). By Proposition \ref{genericpseff}, there exists a dense Zariski open subset $V$ such that $\alpha_s + \tau \omega_s$ is pseudo-effective for $s\in V$. 

Passing to a resolution and further shrinking $V$ we may assume that $f$ is smooth over $V$. Thus, Lemma \ref{vollemma} implies that for any $s\in V$ and any $t> \tau$, $$\vol(\alpha_s + t\omega_s) \ge (t-\tau)^n v>0,$$ with $n = \dim_{\mathbb{C}} (X/S)$ and $v =\int _{X_s} \omega_s^n$. Therefore, by Proposition \ref{BigVol}, for any $t> \tau$, $\alpha + t\omega$ is $f$-big, and thus $\tau' \le  \tau$.
\end{proof}

\begin{lemma}\label{nonpseffcriterion}
Let $f:X\to S$ be a proper surjective morphism with connected fibers from a normal K\"ahler variety onto a relatively compact base $S$, and $\omega \in H^{1,1}_{\rm{BC}}(X,\mathbb{R})$ a K\"ahler class on $X$. Then $\alpha$ is not $f$-pseudo-effective over $S$ if and only if the pseudo-effective threshold $\tau(X/S,\alpha ; \omega)> 0$.
\end{lemma}
\begin{proof}
Assume that $\alpha$ is $f$-pseudo-effective over $S$. Then we have $\tau(X/S, \alpha; \omega) = 0$ by definition. 
Conversely, if $\tau(X/S, \alpha; \omega) = 0$, then for any $\varepsilon > 0$, the class $\alpha + \varepsilon \omega$ is $f$-big. By Proposition \ref{bigandpseff}, $\alpha$ is $f$-pseudo-effective over $S$. 
\end{proof}


\begin{remark}\label{r-pseft} If $S$ is relatively compact, then the relative pseudo-effective threshold
$$\tau(X,\alpha,S;\omega) = \inf\{ t\in \mathbb{R}_{\ge 0}\mid  \alpha + t \omega \text{ is }f\text{-pseudo-effective over some neighborhood } U\subset S \}$$ 
is well defined, i.e.,\ $\tau (X, \alpha ,S ;\omega)<+\infty$. This follows since for $t\gg 0$, 
$\alpha+t\omega$ is K\"ahler (and hence pseudo-effective) over $S$. Thus, $\tau (X,\alpha ,S;\omega)<+\infty$.
\end{remark}

\begin{lemma}\label{Lemma: relativebig}
Let $f:X\to S$ be a proper surjective morphism from a normal K\"ahler variety onto a relatively compact base $S$. If $\alpha \in H^{1,1}_{\rm{BC}}(X, \R)$ is relatively big over $S$, then $-\alpha$ is not relatively big over $S$. 
\end{lemma}
\begin{proof}
To see this, assume by contradiction that  $-\alpha$ is big over a general fiber $X_s$. 
Then there exist K\"ahler currents
\[
T_+ \ge \varepsilon_+ \omega',\quad T_- \ge \varepsilon_- \omega'
\]
for some K\"ahler metric $\omega'$ on $X_s$ and $\varepsilon_+, \varepsilon_->0$, such that $[T_+]=\alpha|_{X_s},\ [T_-]=-\alpha|_{X_s}$.

Since $[T_+ + T_-] = 0$ and $T_+  +T_- \ge 0$, we have $T_+ + T_-= 0$ (because a plurisubharmonic function on a compact complex variety is constant). This is impossible since then $$0 =T_+ + T_- \ge (\varepsilon_+ + \varepsilon_-) \omega'.$$
\end{proof}

\begin{proposition}\label{NonPSEFFMFS}
Let $f \colon X \to S$ be a proper surjective morphism from a generalized K\"ahler pair $(X, B+\boldsymbol{\beta})$ to a relatively compact base $S$.  
If $-(K_X+B+\boldsymbol{\beta}_X)$ is relatively K\"ahler over $S$, then the adjoint class
$K_X+B+\boldsymbol{\beta}_X$ is not relatively pseudo-effective over $S$.  
Moreover, if $S$ is compact, then $K_X+B+\boldsymbol{\beta}_X$ is not (absolutely) pseudo-effective on $X$.
\end{proposition}
\begin{proof}
Let $\omega$ be a K\"ahler class on $X$. Since $S$ is relatively compact, 
$-(K_X+B+{{\bbeta}_X}+\varepsilon\omega)$ is still relatively K\"ahler over $S$ for sufficiently small $0<\varepsilon\ll 1$. By Lemma \ref{Lemma: relativebig}, 
$K_X+B+{{\bbeta}_X}+\varepsilon\omega$ is not relatively big over $S$.

Therefore the pseudo-effective threshold is $\tau(X,B+ { \boldsymbol{\beta}_X},S;A)>0$. By Lemma \ref{nonpseffcriterion}, $K_X+B +  { {\bbeta}_X}$ is not relatively pseudo-effective over $S$. When $S$ is compact, Proposition \ref{AbsRelativepseff} implies that $K_X+B+\bbeta_X$ is not absolutely pseudo-effective on $X$ either. 
\end{proof} 
\begin{lemma}\label{Lemma: PSEFandMFS}
Let $X$ be a compact K\"ahler threefold with the K\"ahler class $\omega$, and $(X,B+\bbeta)$ a generalized klt pair. Assume that $K_X+B+\bbeta_X$ is not pseudo-effective and that the $(K_X+B+\bbeta_X)$-MMP with scaling of $\omega$, $$
(X,B+\bbeta_X) \dashrightarrow (X_1,B_1+\bbeta_{X_{1}}) \dashrightarrow  \cdots 
\dashrightarrow (X_N,B_N+\bbeta_{X_{N}}) = (X',B'+\bbeta_{X'})
$$ terminates with a Mori fibre space $f:X' \to Z'$.
Set 
$$
\tau := \inf \bigl\{ t \ge 0 \,\big|\, K_X + B + \bbeta_X + t\omega \text{ is pseudo-effective}\bigr\}(>0).
$$
Then the following hold:\\
(a) The Mori fibre space $X'\to Z'$ is $(K_{X'}+B'+ \bbeta_{X'} + \tau \omega')$-trivial;\\
(b) Conversely, if the Mori fibre space $X'\to Z'$ is $(K_{X'}+B'+ \bbeta_{X'} + \delta \omega')$-trivial for some $\delta>0$, then $\delta = \tau$ (i.e, $\delta$ is the pseudo-effective threshold $\tau$). 
\end{lemma}
\begin{proof}[Proof of (a)]
Running the $(K_X+B+\bbeta_X)$-MMP with scaling of $\omega$, we obtain a sequence
$$
(X,B+\bbeta_X) \dashrightarrow (X_1,B_1+\bbeta_{X_{1}}) \dashrightarrow  \cdots 
\dashrightarrow (X_N,B_N+\bbeta_{X_{N}}) = (X',B'+\bbeta_{X'}).
$$
Here we denote by $\omega_i$ the push forward of $\omega$ on $X_i$ and  
$$
\lambda_i := \inf \left\{ t \ge 0 \,\middle|\, K_{X_i}+B_i+\bbeta_{X_i}+t\omega_i \text{ is nef} \right\},
$$
where $\{\lambda_i\}$ is a non increasing sequence bounded below by $\tau$
$$
\lambda_1 \ge \lambda_2 \ge \cdots \ge \lambda_i \ge \cdots \ge \lambda_N \ge \tau>0.
$$


Next, we prove that $\lambda_N = \tau$. Since $K_{X_N}+B_N+\bbeta _{X_N}+\lambda _N\omega _N$ is nef, then $\lambda _N\geq \tau$, and it suffices to show that $\lambda_N \le \tau$.  
Since the MMP terminates with a $(K_{X'}+B'+\bbeta_{X'}+\lambda_N\omega')$-trivial Mori fibre space which is $(K_{X'}+B'+\bbeta_{X'})$-negative.  
If $\tau < \lambda_N$, then for the extremal ray $R$ contracted by the Mori fibre space
$X' \to Z'$, we have
$$
\bigl(K_{X'}+B'+\bbeta_{X'}+\tau\omega'\bigr)\cdot R < 0.
$$
This contradicts the pseudo-effectivity of
$K_{X'}+B'+\bbeta_{X'}+\tau\omega'$.  
Hence $\lambda_N = \tau$. 
\end{proof}

\begin{proof}[Proof of (b)]
Conversely, assume that the MMP ends with a $(K_{X'}+ B' + \bbeta_{X'}+ \delta \omega')$-trivial Mori fibre space.

We already showed in (a) that the Mori fibre space is $(K_{X'}+ B' + \bbeta_{X'} + \tau \omega')$-trivial and $(K_{X'}+ B' + \bbeta_{X'})$-negative. Thus
$$
\omega' \cdot R > 0, \quad (\delta - \tau)\,\omega' \cdot R = 0,
$$
and the only possible case is $\delta = \tau$.
\end{proof}
The following proposition was inspired by \cite[Lemma 2.35]{BAB}.

\begin{proposition}\label{pseffthreshold}
Let $X \to S$ be a family from a  K\"ahler variety $X$ onto a smooth, connected, and relatively compact curve $S$ such that $X_t$ admits canonical singularities for $t\in S$, and $\omega$ a K\"ahler class on $X$. 

If $K_X$ is not relatively pseudo-effective over an  open neighborhood $U \subset S$ of $0$, then $(K_X  + \delta \omega)$ is not relatively pseudo-effective over $U$ for some constant $\delta >0$ independent of the choice of $U$.
\end{proposition}

\begin{proof}
Since $K_{X}$ is not relatively pseudo-effective over $U$, there is an uncountable set $B\subset U$, such that $K_{X_t}$ are not pseudo-effective for $t\in B \subset U$. By Lemmata \ref{Lemma: NegCurve} and \ref{BDPPCanonical}, the fibers $X_t$ over $t\in B$ are covered by $K_{X_t}$-negative curves $\{C_{t,s}\}_{s\in S_t}$, where $S_t$ denotes the covering family. 

Since the number of irreducible components of the relative Douady space is countable, there exists a component $D$ of the relative Douady space containing $\{C_{t,s}\}$ for $t\in B' \subset B$ where $B'$ is an uncountable subset of $S$. Therefore there is a covering family of curves $\{C'_{t,s}\}$ (parameterized by the component $D$) that covers the fiber for each $t\in B'$, such that for the K\"ahler class $\omega$ on $X$, $C'_{t,s}\cdot \omega =K>0$ is fixed. Since $X_t$ admits canonical singularities for $t\in S$, $K_X$ is $\Q$-Cartier. Let $R$ be the Cartier index of $K_X$, then $K_{X_t}\cdot C'_{t,s}\leq - 1/R$. So $$(K_{X_t}+\frac{1}{RK+1}\omega_t)\cdot C'_{t,s}<0$$ and thus $K_{X_t}+\frac{1}{RK+1}\omega_t$ is not pseudo-effective for fibers $X_t$ at $t\in B'\subset U$. In particular, $K_X+\frac{1}{RK+1} \omega $ is not relatively pseudo-effective over $U$. 
\end{proof}



%%%%%%%%%%%%%%%%%%%%%%%%%%%%%%%%%%%%%%%%%%%%%%%%%%%%%%%%%%%%%%%%%%%%%%%%%%%%%%%%%%%%%%%%%%%%%%%%%%%%%%%%%%%%%%%%%%%%%%%%%%%%%%%%%
%%%%%%%%%%%%%%%%%%%%%%%%%%%%%%%%%%%%%%%%%%%%%%%%%%%% Subsection Positivity in Kahler family %%%%%%%%%%%%%%%%%%%%%%%%%%%%%%%%%%%%%
%%%%%%%%%%%%%%%%%%%%%%%%%%%%%%%%%%%%%%%%%%%%%%%%%%%%%%%%%%%%%%%%%%%%%%%%%%%%%%%%%%%%%%%%%%%%%%%%%%%%%%%%%%%%%%%%%%%%%%%%%%%%%%%%%

\section{Singularities and positivities of line bundles in  families}\label{Section: SingandPositivty}
In this section, we study the deformation behavior of various kinds of singularities and positivities in K\"ahler families. 
\subsection{Deformation of singularities}
We first study the behavior of singularities under deformation, generalizing several results from \cite{FH11}.

\begin{proposition}[{\cite[Proposition 3.5]{FH11}}]\label{DefCanSing}
Let $f: X \to S$ be a flat proper morphism from a generalized K\"ahler pair $(X, B+\boldsymbol{\beta})$ onto a smooth, connected curve $S$. Fix a point $0 \in S$ and assume that the boundary divisor $B$ does not contain the fiber $X_0$, and that the pair $(X_0, B_0+\boldsymbol{\beta}_0)$ has canonical singularities with $\lfloor B_0 \rfloor = 0$, where $(K_X + B +X_0 + \boldsymbol{\beta}_X)|_{X_0} = K_{X_0} + B_0 + \boldsymbol{\beta}_0$. Then, after shrinking $S$, the pair $(X, B+\boldsymbol{\beta})$ has canonical singularities, and $(X_t, B_t+\boldsymbol{\beta}_t)$ has canonical singularities for any $t \in S$.
\end{proposition}
\begin{proof}
By assumption, the pair $(X_0, B_0+  { \boldsymbol{\beta}}_0)$ is canonical and in particular, $\lfloor B_0\rfloor = 0$. By \cite[Theorem 0.1]{HP24}, the pair $(X,B+X_0+  { \boldsymbol{\beta}})$ is generalized plt near $X_0$.  %It follows easily from the arguments of \cite[Corollary 4.10]{KollarSing} that $(X,B+X_t+  { \boldsymbol{\beta}})$ is glc for $t$ in a neighborhood of $0\in S$. 
By the arguments of \cite[Corollary 1.4.3]{BirkarCHM}, it follows that $(X, B+ { \boldsymbol{\beta}})$ has canonical singularities, and $\left(X_t, B_t+ { \boldsymbol{\beta}}_t\right)$ has canonical singularities for any $t \in S$. 
%We will now show that $(X,B +  { \boldsymbol{\beta}})$ admits canonical singularities near $X_0$. Since $(X,B+\boldsymbol{\beta})$ is generalized klt, by the proof of \cite[Theorem 1.5]{HP24}, there is a proper birational morphism $\nu :X'\to X$ such that $X'$ is strongly $\mathbb Q$-factorial and $a(X,B,\boldsymbol{\beta}; E)<0$ if and only if $E$ is a $\nu$-exceptional divisor. If we write \[K_{X'}+\nu ^{-1}_*B+F+\boldsymbol{\beta}_{X'}=\nu ^*(K_X+B+\boldsymbol{\beta}),\] then $(X', \nu ^{-1}_*B+F+\boldsymbol{\beta})$ is canonical and hence so is $(X', \nu ^{-1}_*B+F)$ and $F\ne 0$ if and only if $(X,B+\boldsymbol{\beta})$ is not canonical. Let $\phi :X'\dasharrow X^m$ be the output of the relative $(K_{X'}+\nu ^{-1}_*B)$-MMP over $X$ \cite{FujinoBCHM}. We let $\mu:X^m\to X$ be the induced map and $F^m=\phi _*F$. We cf. that $K_{X'}+\mu ^{-1}_*B$ is nef over $X$ and since this mmp does not contract components of $\nu ^{-1}B$, then $(X^m, \mu ^{-1}_*B)$ is canonical.  We also have that $\mu(\mathrm{Supp } F^m)\cap X_0\ne 0$ if and only if $(X,B+\boldsymbol{\beta})$ is not canonical on any neighborhood of $X_0$. Notice however that $-F^m$ is nef over $X$ and so by the negativity lemma, for every $x \in \mu(\operatorname{Supp}(F^m))$, we have $\mu^{-1}(x) \subseteq \operatorname{Supp}(F^m)$. It follows that if $\mu (\mathrm{Supp} F^m)\cap X_0\ne 0$, then $\mathrm{Supp}F^m\cap \mu ^{-1}_*X_0\ne 0$. If $X^m_0=\mu ^{-1}_*X_0$ and $\mu_0=\mu |_{X^m_0}$ then  \[ K_{X^m_0}+B_0^m+\boldsymbol{\beta}_0^m=(K_{X^m}+\mu ^{-1}_*B+F^m+\boldsymbol{\beta}_{X^m})|_{X^m_0}=\mu_0^*((K_X+B+X_0+\boldsymbol{\beta}_X)|_{X_0})=\mu_0^*(K_{X_0}+B_0+\boldsymbol{\beta}_0).\]   It follows that $B^m_0-(\mu _0^{-1})_*B_0\ge F|_{X^m_0}\ne 0$ and hence $(X_0,B_0+\boldsymbol{\beta}_0)$ is not canonical.
\end{proof}


\begin{proposition}\label{DefKLT}
Let $f: X \rightarrow S$ be a flat proper morphism from %a normal $\mathbb{Q}$-factorial 
generalized K\"ahler pair $(X, B+ { \boldsymbol{\beta}})$ onto a smooth, connected curve $S$. Fix a point $0\in S$ and assume that the boundary divisor $B$ does not contain the fiber $X_0$, and that the pair $(X_0, B_0+ { \boldsymbol{\beta}}_0)$ is generalized klt. Then, after shrinking $S$, the pair $(X, B+ { \boldsymbol{\beta}})$ is generalized klt and  $(X_t, B_t+\left. { \boldsymbol{\beta}}_t\right)$ has generalized klt singularities for any $t\in S$.
\end{proposition}

\begin{proof}
The proof is similar to the canonical singularities case. Since $(X_0,B_0 +  { \boldsymbol{\beta}}_0)$ is generalized klt, and $(X,B+X_0+ { \boldsymbol{\beta}})$ is a generalized pair. By inversion of adjunction \cite[Theorem 0.1]{HP24}, we know that $(X,B+X_0+  { \boldsymbol{\beta}})$ is generalized plt. Arguing as above, $(X,B+X_t+  { \boldsymbol{\beta}})$ is glc and since $X_t$ is Cartier, the pair $(X_t, B_t +  { \boldsymbol{\beta}}_t)$ is generalized klt for $t\in S$.
\end{proof}

\begin{remark}
By \cite{Kawamata99} it can happen that $(X_0, S_0)$ admits klt singularities but $(X,S)$ does not have klt singularities as $K_X+S$ is not $\mathbb R$-Cartier.  In other words, the $\mathbb R$-Cartier property is not well behaved for families of klt pairs. 
\end{remark}
% \begin{example}
% Let $X\to S$ be a family of surfaces whose central fiber $X_0$ is a $\mathbb{Q}$-factorial klt surface. If $X\to Y/S$ be a flipping contraction whose exceptional locus is contained in the central fiber $X_0$, then $Y$ is not $\mathbb{Q}$-factorial, but the central fiber $Y_0$ is $\mathbb{Q}$-factorial and klt. 
% \end{example}
% Under additional assumptions, we can prove the deformation openness of $\mathbb{Q}$-factorial singularities.
% \begin{proposition}[{\cite[Corollary 3.2]{FH11}}]
% Let $S$ be a normal irreducible Cartier divisor on a complex analytic variety $X$. If $S$ is $\mathbb{Q}$-factorial with terminal singularities, then $X$ is $\mathbb{Q}$-factorial (and normal) in a neighborhood of $S$.
% \end{proposition}

\subsection{Relative ampleness between families}
Given a proper family $f : X \to S$ and a line bundle $ L \in \mathrm{Pic}(X) $, 
if $ L_0 = L|_{X_0} $ is ample on $ X_0 $, then there exists a Zariski open subset of $ S $ 
such that $ L_t = L|_{X_t} $ is ample on $ X_t $ for any $t$ in this subset (cf. \cite[Theorem 1.2.17]{Laza04}). 
This result easily generalizes to the relative setting (Proposition \ref{DefAmple}).
 
\begin{lemma}\label{ResEmbed}
Let $p:X \to S$ and $g:Z \to  Y$ be proper morphisms of complex analytic varieties where $p$ is flat, and let $\Phi: X\to Z$ be a morphism such that the diagram commutes 
\begin{center}
\begin{tikzcd}[ampersand replacement=\&]
	X \&\& Z \\
	\& Y \\
	\& S.
	\arrow["\Phi", from=1-1, to=1-3]
	\arrow["f", from=1-1, to=2-2]
	\arrow["p"', from=1-1, to=3-2]
	\arrow["g"', from=1-3, to=2-2]
	\arrow["q", from=1-3, to=3-2]
	\arrow["r"{description}, from=2-2, to=3-2]
\end{tikzcd}
\end{center}
Fix a point $0\in S$ and assume that the restriction of the diagram $\Phi_0: X_0 \rightarrow Z_0/Y_0$ over $0\in S$ is a closed immersion (resp. an isomorphism). Then $\Phi:X\to Z /Y$ is a closed immersion (resp. an isomorphism) over some open neighborhood $U\subset S$ of $0$. 
\end{lemma}


\begin{proposition}\label{DefAmple}
Let $f: X \to Y/S$ be a proper surjective contraction morphism of complex analytic varieties over a smooth connected curve $S$. For each $t \in S$, denote by
\[
f_t : X_t \to Y_t
\]
the induced proper contraction on the fiber over $t$. 

(1) If $L$ is a $\mathbb{Q}$-line bundle on $X$ such that $L|_{X_0}$ is $f_0$-ample, then there exists a Zariski open subset $U \subset S$ around $0$ such that $L|_{X_t}$ is $f_t$-ample for any $t \in U$.

(2) If $\alpha\in H^{1,1}_{\rm{BC}}(X, \R)$ is a class on $X$ such that $\alpha|_{X_0}$ is $f_0$-K\"ahler, then there exists a (Euclidean) open subset $U\subset S$ around $0$ such that $\alpha|_{X_t}$ is $f_t$-K\"ahler for any $t\in U$. 
\end{proposition}
\begin{proof}
Since $L_0 = L|_{X_0}$ is $f_0$-ample,  there exists some $m\ge 0$ such that
$$\rho_m : f^* f_* (L^{\otimes m})\to L^{\otimes m}$$is surjective near $X_0$. To see this, consider the exact sequence 
\[
0 \to L^{\otimes m} \otimes \mathcal{I}_{X_0} \to L^{\otimes m} \to L^{\otimes m}|_{X_0} \to 0.
\]
Since $L|_{X_0}$ is $f_0$-ample, Serre’s vanishing theorem gives 
\[
R^i (f_0)_* \bigl(L^{\otimes m}|_{X_0}\bigr) = 0 \quad \text{for } i>0 \text{ and } m \gg 0.
\]
Therefore,
\[
\psi: R^1 f_* (L^{\otimes m}) \otimes \mathcal{I}_{\{0\}} 
\cong R^1 f_*\bigl(L^{\otimes m} \otimes \mathcal{I}_{X_0}\bigr) 
\longrightarrow R^1 f_* (L^{\otimes m})
\]
is surjective for $m \gg 0$, where the first isomorphism follows from the projection formula together with the fact that $\mathcal{I}_{X_0}=f^*\mathcal{I}_{\{0\}}$ and $\mathcal{I}_{\{0\}}$ are invertible on the smooth connected curve $S$. Taking the stalk at $0$ and tensor with $\kappa(0)\cong \mathbb{C}$, we obtain a surjective map
\[
\psi(0): \;  R^1 f_* (L^{\otimes m})_0 \otimes \bigl(\mathfrak{m}_0 / \mathfrak{m}_0^2\bigr)
\longrightarrow R^1 f_* (L^{\otimes m})_0 \otimes \kappa(0).
\]
Note that this is the zero map, since the composition in the short exact sequence 
\[
0\to \mathfrak{m}_{0} /\mathfrak{m}_{0}^2 \to \mathcal{O}_{T,0} /\mathfrak{m}_{0}^2 \to \mathcal{O}_{T,0} /\mathfrak{m}_{0} = \kappa (0) \to 0
\] 
is zero. Therefore, $R^1 f_* (L^{\otimes m})_0 \otimes \kappa(0) = 0$ as well. By Nakayama’s lemma, we conclude that 
\[
R^1 f_* (L^{\otimes m})_0 = 0.
\]
Since $R^1 f_* (L^{\otimes m})$ is coherent, its support is a closed analytic subset. Therefore, $R^1 f_* (L^{\otimes m}) = 0$ in a neighborhood of $0$. Consequently,
\[
f_* L^{\otimes m} \longrightarrow (f_0)_* \bigl(L^{\otimes m}|_{X_0}\bigr)
\]
is surjective. Since $L^{\otimes m}|_{X_0}$ is $f_0$-globally generated for $m\gg 0$, i.e., the map
\[
f_0^* (f_0)_* (L^{\otimes m} |_{X_0} ) \to L^{\otimes m}|_{X_0}
\]
is surjective for $m\gg 0$, the morphism
\[
\rho_m: f^* f_* L^{\otimes m} \longrightarrow L^{\otimes m}, \quad \text{for } m\gg 0
\]
is surjective in a neighborhood of $X_0$.


Therefore, after shrinking $S$, it will induce the commutative diagram  
\begin{center}
\begin{tikzcd}[ampersand replacement=\&]
	X \&\& {\mathbb{P}_Y(f_* L^{\otimes m})} \\
	\& Y \\
	\& S.
	\arrow["\Phi", from=1-1, to=1-3]
	\arrow["f"{description}, from=1-1, to=2-2]
	\arrow["g"', from=1-1, to=3-2]
	\arrow[from=1-3, to=2-2]
	\arrow[from=1-3, to=3-2]
	\arrow["h"{description}, from=2-2, to=3-2]
\end{tikzcd}
\end{center}
Since $L|_{X_0}$ is $f_0$-ample, the restriction to the fiber over $0 \in S$ induces an embedding
$$\Phi_0 = \Phi|_{X_0}: X_0 \hookrightarrow \mathbb{P}_{Y_0}((f_0)_* (L^{\otimes m}|_{X_0}))/Y_0.$$
By Lemma \ref{ResEmbed}, after further shrinking the base $S$, $\Phi:X \hookrightarrow \mathbb{P}_{Y}(f_* L^{\otimes m})/Y$ is still a closed embedding.

The second part of the proposition is clear, since positive form is still positive when restricted to the nearby fibers. 
\end{proof}
\subsection{Relative pseudo-effectivity of line bundles in projective families}
Using a recent result of \cite{Jiao}, we prove that the global stability of pseudo-effectivity of a line bundle over a projective family with irreducible and reduced fibers, if this line bundle is pseudo-effective on a Zariski dense set of fibers. 

\begin{proposition}\label{p-psefrest}
Let $f:X\to S$ be a projective flat morphism of complex analytic spaces with reduced and irreducible fibers. If $L$ is a line bundle on $X$ whose restriction to a Zariski  dense set of fibers is pseudo-effective, then $L$ is $f$-pseudo-effective over $S$, and moreover $L|_{X_s}$ is pseudo-effective for every $s \in S$. 
\end{proposition}
\begin{proof}
Let $ A $ be a relatively ample line bundle on $ X $. Assume that there exists a dense subset $ B \subset S $ such that $ L|_{X_s} $ is pseudo-effective for any $ s \in B $. Then, for any $ \varepsilon > 0 $, 
\[
\mathrm{vol}\big((L + \varepsilon A)|_{X_s}\big) \ge  \varepsilon ^nv \quad \text{for all } s \in B,
\] where $n=\dim_{\mathbb{C}} (X/S)$ and $v=(A|_{X_s})^n$ which is independent of $s\in S$.
By Lemma \ref{Hilbertscheme}, the family $f:(X,L,A)\to S$ is the pull back of some projective family $(X',L',A') \to S'$ via a dominant morphism $\phi:S\to S'$ such that $S'$ is a projective scheme. Let $\eta$ be the generic point of $S'$ and denote the generic fiber of $X' \to S'$ by $X'_{\eta}$. Since $\phi:S\to S'$ is a dominant morphism, the image $B' = \phi(B)$ is dense in $S'$. By \cite[Theorem 1.1]{Jiao},
\[
\mathrm{vol}\big((L'+\varepsilon A')|_{X'_\eta}\big) = \inf_{s\in B'} \mathrm{vol}\big((L'+\varepsilon A')|_{X'_s}\big) \ge \varepsilon ^nv >0.
\]
In particular, for every $\varepsilon > 0$, the divisor $L' + \varepsilon A'$ is big on the generic fiber. It follows that $L'$ is relatively pseudo-effective over $S'$, and hence $L$ is relatively pseudo-effective over $S$.

For any $\varepsilon > 0$, the line bundle $L + \varepsilon A$ is relatively big over $S$. That is, there exists a dense Zariski open subset $U_{\varepsilon} \subseteq S$ such that $(L+\varepsilon A)|_{X_s}$ is big for every $s \in U_{\varepsilon}$. By \cite[Proposition~4.1]{RT1}, the restrictions $(L+\varepsilon A)|_{X_s}$ are big for all $s \in S$. Since $\varepsilon > 0$ is arbitrary, it follows that $L|_{X_s}$ is pseudo-effective for every $s \in S$.
\end{proof}


% \begin{lemma}
% Let $f:X \rightarrow Y$ be a morphism of proper algebraic varieties, with analytification $f^{an}: X^{an} \to Y^{an}$, and $Y$ is irreducible with generic point $\eta$. Let $(P)$ be a property of varieties. If the generic fiber $X_\eta$ satisfies $(P)$, and the property is constructible, then a general fiber satisfies  $(P)$ as well, meaning that there is an (analytic) Zariski open subset $U^{an} \subset Y^{an}$ such that $X_s^{an}$ satisfies Property $(P)$ for every $s \in U^{an}$.
% \end{lemma}

% \begin{proof}
% Let $$S = \{t\in Y \mid X_t \text{ satisfies property }(P)\}.$$By assumption, $S$ is a constructible set, so that it can be written in the form $S= \bigcup_{i =1}^r  (U_i \cap Z_i)$ for some Zariski closed subsets $Z_i$ and Zariski open subsets $U_i$. By assumption $\eta \in S = \bigcup_{i=1}^r (U_i \cap Z_i)$. In particular, there exists some $i$, such that $\eta \in U_i \cap Z_i$. Since $\eta$ is the generic point of $Y$, if a Zariski closed subset $Z_i$ contains that generic point $\eta$, then $Y = \bar{\eta} \subset \bar{Z_i}  = Z_i$. In particular, we find some Zariski open subset $\eta\in U_i \subset S$, such that over that Zariski open subset $U_i$, all the fibers satisfies property $(P)$. Therefore, there exists a Zariski open subset $U_i^{an} \subset Y^{an}$ such $X^{an}_s$ satisfies property $(P)$ for every $s\in U_i^{an}$.
% \end{proof}

\begin{lemma}\label{Hilbertscheme}
If $ f : X \to S $ is a projective flat morphism of complex analytic spaces with reduced and irreducible fibers, together with a line bundle $ L $ on $ X $, then there exists a projective family $f' : X' \to S'$ over some projective scheme $ S' $, equipped with a line bundle $ L' $ on $ X' $ and a dominant morphism $ \phi : S \to S' $, such that $f : (X, L) \to S$ is the pullback of
\[
f' : (X', L') \to S'
\]
via $ \phi : S \to S' $.
\end{lemma}
\begin{proof}
Since $f: X \to S$ is a projective morphism, there exists a closed embedding 
$X \hookrightarrow \mathbb{P}_S^N / S$, where $N$ is some fixed positive integer. 

Let $A$ be a sufficiently $f$-ample line bundle on $f \colon X \to S$. 
If $m \gg 0$, then $M := L \otimes A^{\otimes m}$
is sufficiently $f$-ample, relatively globally generated, and $f_*(M)$ is locally free. 
In particular, a general section $\sigma \in H^0(S, f_*M) = H^0(X, M)$ cuts out a relatively effective divisor $D$ such that 
\[
M|_{X_s} \cong \mathcal{O}_{X_s}(D|_{X_s}).
\]
Since $f$ is flat, the fibers of $f$ define a flat family $
\{X_s\}_{s \in S}$ of projective varieties in $\mathbb{P}^N$. 
Note that the Hilbert polynomial 
\[
\Phi_{X_s}(k) = \chi\!\left(\mathcal{O}_{X_s} \otimes A^{\otimes k}\right)
\]
with respect to the ample line bundle $A$ is constant under flat deformations, denoted by $\Phi(k)$.

Therefore, by the representability of the flag Hilbert scheme (cf. \cite[Theorem 4.5.1]{Sernesi}), there exists a projective scheme $H$ and a classification map 
\[
\phi: S \to H, \quad s \mapsto [X_s].
\]
The family $(X,D) \to S$ is the pullback of the projective universal family $f': (X', D')\to H$ via $\phi:S\to H$, and the canonical line bundle $\mathcal{O}_{X'/H}(1)$ on $X'$ pull back to the ample line bundle $A$ on $X$. 

If we denote by $S' = \phi(S)$ the scheme theoretical image, then $\phi:S \to S'$ is a dominant morphism. 
We now consider the line bundle
\[
L' = \mathcal{O}_{X'}({D}') \otimes \mathcal{O}_{X'/S'}(-m).
\]
Clearly, the pair $(X', L') \to H$ pulls back to the pair $(X, L) \to S$ via the morphism $\phi : S \to S'.$ 
\end{proof}

\subsection{Global generation and semi-ampleness of line bundles in families}
Next, we discuss the global generation of line bundles in families.
\begin{lemma}\label{ResBasechange}
Let $f: X \to S$ be a proper surjective morphism of normal complex analytic spaces, and $L$ a line bundle on $X$. Assume that $S$ is Stein.
Then for any $s\in S$, the base change map $$\varphi^0(s): (f_* L)_s \otimes_{\mathcal{O}_{S,s}} \kappa(s) \;\longrightarrow\; H^0(X_s, L|_{X_s})$$ is surjective if and only if the restriction map $$r: H^0(X,L) \to H^0(X_s, L|_{X_s})$$ is surjective. 

\end{lemma}
\begin{proof}
By Cartan's Theorem A, the direct image sheaf $f_*L$ is globally generated on $S$. Therefore the evaluation map $\alpha_s$ in the diagram is surjective
\[
\begin{tikzcd}[ampersand replacement=\&]
	{H^0(S,f_* L)} \& {(f_*L)(s)} \\
	{H^0(X,L)} \& {H^0(X_s,L|_{X_s})}.
	\arrow["{\alpha_s}", from=1-1, to=1-2]
	\arrow["{=}"', from=1-1, to=2-1]
	\arrow["{\varphi^0(s)}", from=1-2, to=2-2]
	\arrow["r", from=2-1, to=2-2]
\end{tikzcd}
\]
Thus, $\varphi^0(s)$ is surjective if and only if $r$ is surjective.
%Moreover, if we assume that  $\varphi^0(s)$ is surjective, then the natural restriction $$r: H^0(X,L) \to H^0(X_s, L|_{X_s})$$is surjective.
\end{proof}
\begin{lemma}\label{BsoverStein}
With the same setting as in Lemma \ref{ResBasechange}, we have
\[
\operatorname{coker}\!\bigl(f^* f_* L \to L\bigr) \;=\; \operatorname{coker}\!\bigl(H^0(X,L)\otimes \mathcal{O}_X \to L\bigr).
\]
In particular,
\[
\operatorname{Bs}(L/S) \;=\; \operatorname{Bs}(L) \;=\; \bigcap_{\sigma \in H^0(X,L)} V(\sigma).
\]
Here the \emph{base loci} $\operatorname{Bs}(L/S)$ and $\operatorname{Bs}(L)$ of $L/S$ and $L$ are defined as $$\operatorname{Bs}(L/S) : = \mathrm{supp}\operatorname{coker}\!\bigl(f^* f_* L \to L\bigr)$$ and $$\operatorname{Bs}(L):=\mathrm{supp}\operatorname{coker}\!\bigl(H^0(X,L)\otimes \mathcal{O}_X \to L\bigr),$$ respectively. 
\end{lemma}
\begin{proof}
By Cartan's Theorem~A, the coherent sheaf $f_*L$ is globally generated on $S$. That is, the natural map 
\[
H^0(X,L)\otimes \mathcal{O}_S \;=\; H^0(S,f_*L)\otimes \mathcal{O}_S \;\longrightarrow\; f_*L
\]
is surjective. Since pullback via $f^*$ preserves surjectivity, we obtain the surjective morphism
\[
H^0(X,L)\otimes \mathcal{O}_X \;\longrightarrow\; f^* f_* L.
\]
So $
\operatorname{coker}\!\bigl(f^* f_* L \to L\bigr) \;=\; 
\operatorname{coker}\!\bigl(H^0(X,L)\otimes \mathcal{O}_X \to L\bigr)$. 
\end{proof}

Based on Lemmata \ref{ResBasechange} and \ref{BsoverStein}, we obtain the equivalence between relative global generation and fiberwise global generation of line bundles under the base change assumption.
\begin{proposition}\label{Bslocus}
%Let $f: X \to Y$ be a proper surjective morphism of complex analytic spaces, and $L$ a line bundle on $X$. Assume that for every $y \in Y$, the base change map \[ (f_* L)_s \otimes_{\mathcal{O}_{Y,y}} \kappa(y) \to H^0(X_s, L|_{X_s}) \] is an isomorphism, where $X_s = f^{-1}(y)$ is the fiber over $y$ and $L_s = L|_{X_s}$. Let $\mathrm{Bs}(L/Y) = \mathrm{supp}(\mathrm{coker}(f^*f_* L \to L))$ and $\mathrm{Bs}(L|_{X_s}) = \mathrm{supp}(\mathrm{coker}(H^0(X_s, L|_{X_s})\otimes_{\kappa(y)}\mathcal{O}_{X_s}\to L|_{X_s}))$. Then $$\mathrm{Bs}(L/Y)|_{X_s} = \mathrm{Bs}(L|_{X_s}) ,\quad \text{for any } y \in Y.$$ In particular, the natural map \[ f^* f_* L \to L \]is surjective if and only if $L|_{X_s}$ is globally generated for every $y \in Y$.

Let $f: X \to S$ be a proper surjective morphism of normal complex analytic spaces, and $L$ a line bundle on $X$. 
Assume that for every $s \in S$, the base change map
\[
\varphi^0(s): (f_* L)_s \otimes_{\mathcal{O}_{S,s}} \kappa(s) \;\longrightarrow\; H^0(X_s, L|_{X_s})
\]
is an isomorphism, where $X_s = f^{-1}(s)$ denotes the fiber over $s$.

Then the natural morphism
$f^{*} f_{*} L \;\longrightarrow\; L$ is surjective if and only if $L|_{X_s}$ is globally generated for every $s \in S$.
\end{proposition}

\begin{proof}
Since surjectivity of a sheaf homomorphism can be checked on an open cover, we may without loss of generality, assume that the base $S$ is Stein. Let $x\in X_s$ be a fixed point on the fiber $X_s$ for any $s\in S$. 
Denote by $L(x)$ and $(L|_{X_s})(x)$ the fibers of $L$ and the restriction $L|_{X_s}$ at $x \in X$, respectively. Clearly, both fibers are isomorphic to $\mathbb{C}$. 
Consider the commutative diagram (in the category of $\mathbb{C}$-vector spaces)
\[
\begin{tikzcd}
H^0(X, L)  \arrow[r, "\widetilde{\alpha}_x"] \arrow[d,"r"] & L(x) \cong \mathbb{C} \arrow[d,"\cong"] \\
H^0(X_s, L|_{X_s})  \arrow[r, "{\alpha}_x"] & {(L|_{X_s})}(x) \cong \mathbb{C},
\end{tikzcd}
\]
where $\widetilde{\alpha}_x$ (resp. ${\alpha}_x$) denotes the evaluation map 
\[
\widetilde{\alpha}_x (\widetilde{\sigma}) = \widetilde{\sigma}(x), \qquad \text{(resp. ${\alpha}_x (\sigma) = \sigma(x)$)},
\]
for $\widetilde{\sigma} \in H^0(X,L)$ (resp. $\sigma\in H^0(X_s,L|_{X_s})$), while $r$ is the natural restriction. 
By assumption, the base-change map 
\[
\varphi^0(y): (f_*L)\otimes \kappa(s) \longrightarrow H^0(X_s, L|_{X_s})
\]
is an isomorphism for every $s \in S$. Thus, by Lemma~\ref{ResBasechange}, the restriction map $r$ is surjective, and hence the cokernels of $\widetilde{\alpha}_x$ and ${\alpha}_x$ coincide.  

%Since a $\mathbb{C}$-linear homomorphism into $\mathbb{C}$ is surjective if and only if its kernel is a proper linear subspace, it follows 
Note that $\widetilde{\alpha}_x$ is surjective precisely when there exists a section $\widetilde{\sigma} \in H^0(X,L)$ such that $\widetilde{\alpha}_x(\widetilde{\sigma}) \neq 0$. Similarly, ${\alpha}_x$ is surjective if and only if there exists a section $\sigma \in H^0(X_s, L|_{X_s})$ with ${\alpha}_x(\sigma) \neq 0$.
Moreover, 
$\operatorname{Bs}(L|_{X_s}) \;=\; \bigcap_{\sigma\in H^0(X_s, L|_{X_s})} V(\sigma)$ and by Lemma~\ref{BsoverStein},  
$$\operatorname{Bs}(L/S) \;=\; \bigcap_{\widetilde{\sigma}\in H^0(X,L)} V(\tilde{\sigma}).
$$
Therefore, ${\alpha}_x$ is surjective if and only if $x \notin \operatorname{Bs}(L|_{X_s})$, and $\widetilde{\alpha}_x$ is surjective if and only if $x \notin \operatorname{Bs}(L/S)|_{X_s}$. It follows that the restriction $L|_{X_s}$ is globally generated at $x \in X_s$ if and only if the map 
\[
f^*f_*L \longrightarrow L
\]
is surjective at $x \in X_s$.
\end{proof}
\begin{remark}\label{SteinBslocus}
Let $f: X \to S$ be a proper surjective morphism of normal complex analytic spaces with Stein $S$ and $L$ a line bundle on $X$. If we assume that at some $s \in S$, the base change map $$\varphi^0(s): (f_* L)_s \otimes_{\mathcal{O}_{S,s}} \kappa(s) \;\longrightarrow\; H^0(X_s, L|_{X_s})$$ is surjective, then
Proposition \ref{Bslocus} shows that for any $x\in X_s$, the following are equivalent:\\
(a) The natural morphism $(f^* f_* L)_x \to L_x$ is surjective;\\
(b) The natural morphism $H^0(X_s,L|_{X_s}) \otimes \mathcal{O}_{X_s ,x} \to (L|_{X_s})_x$ is surjective;\\
(c) The evaluation map $H^0(X_s, L|_{X_s}) \to (L|_{X_s})(x)$ is surjective. \\
Moreover, we have $x \in \mathrm{Bs}(L|_{X_s})$ if and only if $x \in \mathrm{Bs}(L/S)\cap X_s$.
%(c) $x\in \mathrm{Bs}(L|_{X_0}) \iff  x\in \mathrm{Bs}(L/S)|_{X_0}$.\\

\end{remark}
\begin{remark}\label{Grauertbasechange}
Let $f:X\to S$ be a flat proper morphism of normal complex analytic spaces, and $L$ a line bundle on $X$. Then \cite[Corollary~III.3.7]{BS} gives the equivalent conditions for the base change maps to be surjective: \\
(1) The line bundle $L$ satisfies the \emph{base change property in degree $0$}, i.e., $$ \varphi^0(s):  f_*(L) \otimes_{\mathcal{O}_{S,s}} \kappa(s) \rightarrow H^0\left(X_s, L|_{X_s}\right)$$is surjective for any $s\in S$.\\
(2) The $0$-cohomological dimension function $$y \mapsto h^0(X_s, L|_{X_s})$$ is locally constant over $S$.\\
(3) The line bundle $L$ is \emph{cohomologically flat in dimension 0}, i.e., the functor $$F: \text{Coh}_X \to \text{Coh}_S,\quad  \mathcal{M} \mapsto f_*(L \otimes_{\mathcal{O}_S}\mathcal{M})$$is exact. In what follows, we will use the equivalent characterization  implicitly.
\end{remark}

As direct corollaries of Proposition \ref{Bslocus}, we first give several equivalent characterizations of relative global generation of line bundles. 
\begin{corollary}\label{equi-rgg}
Let $f:X\to \Delta$ be a flat proper morphism from a normal complex analytic variety onto the unit disk in $\mathbb{C}$, and $L$ a line bundle on $X$. Then the following are equivalent: \\
\emph{(a)} The natural morphism $f^* f_* L \to L$ is surjective over $f^{-1}(U)$, for some non-empty Zariski open subset $U \subset \Delta$;\\
\emph{(b)} $L|_{X_t}$ is globally generated for any $t \in V$, where $V \subset \Delta$ is a non-empty Zariski open subset;\\
\emph{(c)} There exists a subset $E\subset \Delta$ with an accumulation point contained in $\Delta$, such that $L|_{X_t}$ is globally generated for $t\in E$;\\
\emph{(d)} There exists a point $s_0\in \Delta$, such that the base change map $$\varphi^0(s_0): (f_* L)(s_0)\to H^0(X_s, L|_{X_{s_0}})$$is surjective, and $L|_{X_{s_0}}$ is globally generated. 
\end{corollary}

\begin{proof}
$(a) \Rightarrow (b)$. Since Grauert's base change theorem holds over some Zariski open subset $W\subset \Delta$, $W\cap U = V $ is a non-empty Zariski open subset of $\Delta$ over which $f^*f_* L \to L$ is surjective and the base change property holds. Thus, by Proposition \ref{Bslocus}, $L|_{X_t}$ is globally generated for any $t\in V$. \\
$(b) \Rightarrow (c)$ is clear, since given any point $t_0\in V$ we can always find a countable sequence of points $E = \{t_i\}\subset V$ that converge to $t_0$. \\
$(c) \Rightarrow (d)$. There exists some non-empty Zariski open subset $W\subset \Delta$, such that Grauert's base change holds. Since $E\cap W \ne \emptyset$ (by the identity principle), there exists a point $s\in E\cap W$ such that the base change map $\varphi^0 (s)$ is surjective and $L|_{X_s}$ is globally generated.\\
$(d) \Rightarrow (a)$. By \cite[Theorem III.12.11]{Hartshorne},  
$$W = \{s\in \Delta \mid \varphi^0(s): (f_* L)(s)\to H^0(X_s, L|_{X_s}) \text{ is surjective}\}$$ 
is a non-empty Zariski open subset of $\Delta$, which is Stein. By assumption (d), at the point $s_0\in W\subset \Delta$, $\varphi^0(s_0)$ is surjective and $L|_{X_{s_0}}$ is globally generated. Therefore, by Remark \ref{SteinBslocus} 
$$\mathrm{Bs}(L/\Delta)|_{X_{s_0}} = \mathrm{Bs}(L|_{X_{s_0}}) = \emptyset.$$

In particular, $U = W \setminus f(\mathrm{Bs}(L/\Delta))$ is a non-empty Zariski open subset of $\Delta$, such that $f^*f_* L \to L$ is surjective on $U$.
\end{proof}

Furthermore, as a more special corollary of Proposition \ref{Bslocus}, one establishes the Zariski openness of global generation for line bundles under the assumption that the dimension of their global sections is deformation invariant.
\begin{corollary}\label{opensemiample}
Let $ f: X \to S $ be a flat proper morphism from a normal complex analytic variety $ X $ onto a smooth, connected, and relatively compact curve $S$, and  $ L $ a line bundle on $ X $. Assume that $ L|_{X_0} $ is globally generated and that the dimension $ h^0(X_s, L|_{X_s}) $ is constant for every $ s \in S $. Then there exists a dense Zariski open subset $ U \subset S $ such that $ L|_{X_s} $ is globally generated for any $ s \in U $.
\end{corollary} 
\begin{proof} Although the proof is implicit in that of Corollary \ref{equi-rgg}, we still include a direct proof here. 

Without loss of generality, we may shrink the base around $0$ and assume that $S$ is Stein. Since $L|_{X_0}$ is globally generated, $\text{Bs}(L|_{X_0}) = \emptyset$. By Remark \ref{SteinBslocus}  $$\mathrm{Bs}(L/S)|_{X_0} = \mathrm{Bs}(L|_{X_0}) = \emptyset.$$

Since $\mathrm{Bs}(L/S)$ is an analytic closed subset of $X$, $f(\mathrm{Bs}(L/S)) = W\subset S$ is an analytic closed subset that does not contain $0$. Hence, there exists a Zariski open subset $U = S \setminus W$ such that $$\mathrm{Bs}(L /S)|_{X_s} = \mathrm{Bs}(L|_{X_s}) = \emptyset ,\quad \text{for any } s \in U.$$In particular, $L|_{X_s}$ is globally generated for any $s\in U$.
\end{proof}
% \begin{example}
% A divisor on a smooth, connected, and relatively compact complex analytic curve can be an infinite sum, so that in the proposition above the complement of some Zariski open subset can have infinitely many points inside the unit disk $\Delta$. Consider the Blaschke product $$B(z)=\prod_{n=1}^{\infty} \frac{\left|a_n\right|}{a_n} \frac{a_n-z}{1-\overline{a_n} z},$$with $a_n  = 1- \frac{1}{2^n}$. Since $\sum\left(1-\left|a_n\right|\right)<\infty$, by Morera's theorem, this defines a holomorphic function on the smooth, connected, and relatively compact curve. Thus the vanishing set $D= V(B(z)) = \sum_{n=1}^{\infty} [a_n]$ defines a divisor which is an infinite sum, with the accumulation point $a_n \to 1$ which does not lie in the smooth, connected, and relatively compact curve $S$. The above arguments do not apply directly to this case since the accumulation point does not lie inside the smooth, connected, and relatively compact curve. 
% \end{example}
It is worth noting that the global generation criterion for line bundles developed above (Proposition~\ref{Bslocus}) also applies to relatively flat coherent analytic sheaves of rank $1$.
At the end of this subsection, we apply the Proposition~\ref{Bslocus} to the semi-ampleness of canonical sheaves.
%, \cite{Tosati}
\begin{definition}[Semi-ample coherent sheaves, semi-ample (1,1)-classes]\label{DefinitionSemiample}
Let $X$ be a normal complex analytic variety. Then:\\
\emph{(1)} A coherent analytic sheaf $\mathcal{S}$ on $X$ is \emph{semi-ample} if there exists some sufficiently large $m \gg 0$ such that $\mathcal{S}^{\otimes m}$ is globally generated. \\
\emph{(2)} Let $f: X \rightarrow S$ be a proper surjective morphism of normal complex analytic spaces, and $\mathcal{S}$ a coherent analytic sheaf on $X$. We say that $\mathcal{S}$ is \emph{$f$-semi-ample} (or \emph{relatively semi-ample}) if there exists an integer $m>0$ such that the natural evaluation map
$$
f^* f_*\left(\mathcal{S}^{\otimes m}\right) \longrightarrow \mathcal{S}^{\otimes m}
$$
is surjective as a morphism of sheaves on $X$.\\
\emph{(3)} A class $\alpha \in H^{1,1}_{\mathrm{BC}}(X,\mathbb{R})$ is a \emph{semi-ample $(1,1)$-class} if there exists a contraction morphism  $f : X \to Z$ 
together with a K\"ahler form $\omega_Z$ on the K\"ahler variety  $Z$ such that  
\[
f^*([\omega_Z]) = \alpha.
\] 
\end{definition}

Now let us strengthen \cite[Corollaries 1.2, 1.12]{LRW25} via a different approach. 
\begin{proposition}\label{defsemiample}
Let $f:X\to S$ be a smooth family of projective manifolds (or a smooth K\"ahler family of threefolds, a flat family of varieties with only canonical singularities and uncountably many fibers of general type). Assume that $0\in S$ such that $K_{X_{0}}$ is semi-ample. Then $K_{X_t}$ is semi-ample on the
 general fiber $X_t$ of $f$. 

In particular, if $f$ is a smooth projective family or a smooth K\"ahler family of threefolds and the canonical divisor of one fiber of $f$ is semi-ample, then the canonical divisor $K_X$ is relatively semi-ample and the canonical divisor of any fiber of $f$ is semi-ample. 
\end{proposition}
\begin{proof}
By \cite[Theorem 1.2]{RT2} and \cite[Corollary 1.12]{LRW25} (or Theorem \ref{DefVol3fold} here for the K\"ahler case), the deformation invariance of plurigenera holds for each family in the assumption. Then Corollary \ref{opensemiample} implies that the canonical sheaf of a general fiber of $f$ is semi-ample (and hence nef).

In particular, when $f$ is a smooth projective family or a smooth K\"ahler family of threefolds, \cite[Theorems 1.1, 1.10]{LRW25} gives the deformation closedness of nefness. Thus, the canonical divisors of all fibers of $f$ are nef. Then Kawamata's relative freeness theorem \cite{Kawa85} (also \cite[Theorem 5.8]{Naka85} and \cite[Theorem 1]{FujinoKawa}) yields that $K_X$ is $f$-semi-ample. By Proposition \ref{Bslocus}, this is equivalent to say that the canonical divisor of any fiber of $f$ is semi-ample.  
\end{proof}




%%%%%%%%%%%%%%%%%%%%%%%%%%%%%%%%%%%%%%%%%%%%%%%%%%%%%%%%%%%%%%%%%%%%%%%%%%%%%%%%%%%%%%%%%%%%%%%%%%%%%%%%%%%%%%%%%%%%%%%%%%%%%%%%%
%%%%%%%%%%%%%%%%%%%%%%%%%%%%%%%%%%%%%%%%%%%%%%%%%%%% Subsection Extension of MMP  %%%%%%%%%%%%%%%%%%%%%%%%%%%%%%%%%%%%%%%%%%%%%%%
%%%%%%%%%%%%%%%%%%%%%%%%%%%%%%%%%%%%%%%%%%%%%%%%%%%%%%%%%%%%%%%%%%%%%%%%%%%%%%%%%%%%%%%%%%%%%%%%%%%%%%%%%%%%%%%%%%%%%%%%%%%%%%%%%
\section{Extension of the analytic MMP and stability of nefness and bigness}\label{Section: ExtMMP}
The relative minimal model program is a powerful tool in the study of deformation problems. The purpose of this section is to study the behavior of the minimal model program under deformations. The definitions below follow \cite{DHY23}. 
\begin{definition}[Divisorial contractions]\label{DefinitionDivisorial}
Let $(X,B+\boldsymbol{\beta})$ be a generalized pair.
A $(K_X+B+\boldsymbol{\beta}_X)$-\emph{divisorial contraction} $f:X\to Z$ is a bimeromorphic morphism with the relative Picard number $\rho(X/Z) = 1$ and exceptional locus being an (irreducible) divisor, such that $-(K_X+B+ \boldsymbol{\beta}_X)$ is K\"ahler over $Z$.
\end{definition}

\begin{definition}[Small bimeromorphic maps]
Let $X$ and $Y$ be complex analytic varieties. A bimeromorphic map $\phi:X \dashrightarrow Y$ is a \emph{bimeromorphic contraction} if there exists a resolution $p: W \longrightarrow X$ and $q: W \longrightarrow Y$ of $\phi$ such that $p$ and $q$ are contraction morphisms and every $p$-exceptional divisor is $q$-exceptional. We say that $\phi$ is a \emph{small bimeromorphic map} if both $\phi$ and $\phi^{-1}$ are bimeromorphic contractions. 
\end{definition}

\begin{definition}[Flipping contractions, flips]\label{DefinitionFlip} Let $(X,B+\boldsymbol{\beta})$ be a generalized pair.
A $(K_X+B+\boldsymbol{\beta}_X)$-\emph{flipping contraction} $f: X \rightarrow Z$ is a small bimeromorphic morphism such that $\rho(X / Z)=1$ and $-\left(K_X+B+ { \boldsymbol{\beta}}_X\right)$ is K\"ahler over $Z$. The corresponding \emph{flip} (if it exists) is a small bimeromorphic morphism $f^{+}: X^{+} \rightarrow$ $Z$ such that $\rho\left(X^{+} / Z\right)=1$ and $K_{X^{+}}+B^{+}+ \boldsymbol{\beta}_{X^{+}}$ is  K\"ahler over $Z$, where $B^{+}$ is the strict transform of $B$.
\end{definition}
\begin{remark}
By \cite[Theorem 5.12]{DH24}, the flipping contraction is a locally projective morphism. When the transcendental class satisfies $ { \boldsymbol{\beta}}_X = 0 $ in Definitions \ref{DefinitionDivisorial} and \ref{DefinitionFlip}, the divisorial and flipping contractions are projective morphisms.
\end{remark}
\begin{definition}[{Mori fiber spaces of a generalized K\"ahler pair}]\label{MFSKahler}
If $(X / S, B+ { \boldsymbol{\beta}})$ is a generalized dlt pair over a relatively compact analytic variety $S$, then we say that $f:(X/S,B + { \boldsymbol{\beta}})\to Z/S$ is a \emph{Mori fiber space of $(X,B+  { \boldsymbol{\beta}})$} over $S$ if $f$ is a contraction morphism associated to a $(K_X+B+ { {\bbeta}_X})$-negative extremal face such that $-(K_X+B+ { {\bbeta}_X})$ is $f$-K\"ahler and $\dim Z < \dim X$.
\end{definition}

\begin{remark}
Some references require the additional assumption that relative Picard number satisfies $ \rho(X / S) = \rho(Z / S) + 1 $ (e.g., \cite{DHY23}). For simplicity and without causing any problem, we omit this condition in our discussion.
\end{remark}
\subsection{Extension of the analytic MMP to nearby fibers} By \cite{KM92}, we have the crucial extension result: 
\begin{theorem}[{Extension of a contraction morphism, \cite[Proposition (11.4)]{KM92}}]\label{ExtendCont}
Let $\phi_0: X_0 \to X_{0}'$ be a proper surjective morphism of normal compact complex analytic spaces. Assume that $(\phi_0)_* \mathcal{O}_{X_0} = \mathcal{O}_{X_0'}$ and $R^1(\phi_0)_* \mathcal{O}_{X_0} = 0$. If $f:X\to S$ is a small deformation of $X_0$, then $\phi_0$ can be extended to a contraction morphism $\phi: X\to X'$ over $S$:
% \begin{center}
% \begin{tikzcd}[ampersand replacement=\&]
% 	{X_0} \&\& {X} \\
% 	\& {X_0'} \&\& {X'} \\
% 	{\{0\}} \&\& S.
% 	\arrow[hook, from=1-1, to=1-3]
% 	\arrow["{\phi_0}", from=1-1, to=2-2]
% 	\arrow[from=1-1, to=3-1]
% 	\arrow["{\phi}", from=1-3, to=2-4]
% 	\arrow[from=1-3, to=3-3]
% 	\arrow[hook, from=2-2, to=2-4]
% 	\arrow[from=2-2, to=3-1]
% 	\arrow[from=2-4, to=3-3]
% 	\arrow[hook, from=3-1, to=3-3]
% \end{tikzcd}
% \end{center}
\[\begin{tikzcd}
	{X_0} && X \\
	& {X'_0} & {} & {X'} \\
	{\{0\}} && {S.}
	\arrow[hook, from=1-1, to=1-3]
	\arrow["{\phi_0}", from=1-1, to=2-2]
	\arrow[from=1-1, to=3-1]
	\arrow["\phi", from=1-3, to=2-4]
	\arrow[from=1-3, to=3-3]
	\arrow[hook, no head, from=2-2, to=2-3]
	\arrow[from=2-2, to=3-1]
	\arrow[from=2-3, to=2-4]
	\arrow[from=2-4, to=3-3]
	\arrow[hook, from=3-1, to=3-3]
\end{tikzcd}\]
\end{theorem}
\begin{remark}
The result does not hold without the condition  `$R^1 (\phi_0)_* \mathcal{O}_{X_0} = 0$'. Consider the product of abelian varieties $ \pi_2 : A_1 \times A_2 \to A_2 $; note that
\[
R^1 (\pi_2)_* \mathcal{O}_{A_1 \times A_2} \neq 0,
\]
since $\dim H^1(A_1,\mathcal{O}_{A_1}) = \dim A_1  >0$. In this case, a general deformation of $ A_1 \times A_2 $ is a simple abelian variety $ \mathcal{A} $, and there cannot exist any non-trivial morphism to a lower-dimensional abelian variety $ \mathcal{A}\to \mathcal{A}' $.
\end{remark}

% \begin{remark}[Remark about Theorem \ref{ExtendCont}]
% Let us compare Theorem \ref{ExtendCont} with the blow-down theorem of Fujiki: 
% \begin{theorem}[{\cite[Theorem 2]{Fujiki75}}]
% Let $X$ be a reduced complex space, $S$ an effective Cartier divisor on $X$ and $f: S \rightarrow S^{\prime}$ a proper surjective morphism of $S$ onto another complex space $S^{\prime}$. Assume that $\mathcal{O}_S(-S)$ is $f$-ample and $R^1 f_* \mathcal{O}_S(-k S)=0$ for all $k>0$. Then there exists a modification $\psi: X \rightarrow Y$ with $\psi |_ S=f$.  
% \end{theorem}
% Fujiki's blow-down theorem requires checking the ampleness and vanishing conditions on the conormal bundle of a (not necessarily reduced) Cartier divisor (cf. \cite[Section 4 and Theorem 5.8]{DH24}). On the other hand, the positivity condition on the conormal bundle is not needed in Theorem \ref{ExtendCont}. Both Fujiki's blow-down theorem and Theorem \ref{ExtendCont} require the vanishing of the first direct image, as this sheaf contains the obstruction to extending the structure sheaf.
% \end{remark}


We now apply Theorem \ref{ExtendCont} to the study of the minimal model program under deformations. Our first result concerns one of its building blocks: Mori fiber spaces under deformations.

\begin{proposition}[Extension of Mori fiber spaces]\label{extFano}
Let $f:X\to S$ be a flat proper morphism from a generalized pair $(X,B+ { \boldsymbol{\beta}})$  onto a smooth, connected, and relatively compact curve $S$. Fix a point $0\in S$ and assume that the boundary divisor $B$ does not contain the fiber $X_0$, and that $(X_0,B_0 + { \boldsymbol{\beta}}_0)$ be a generalized klt pair, where $(K_X+X_0+B+\bbeta _X)|_{X_0}=K_{X_0}+B_0+\bbeta _0$. 

If there is a $(K_{X_0}+B_0+ \bbeta_0)$-Mori fiber space $\phi_0:X_0 \to Z_0$ (cf. Definition \ref{MFSKahler}), then the family admits a relative Mori fiber space $\phi: X\to Z /U$ over a neighborhood $U \subset S$ of $0$.
\end{proposition}

\begin{proof}
Since $\phi_0 : X_0\to Z_0$ is a $(K_{X_0} +B_0+ \bbeta_0)$-Mori fiber space, $-(K_{X_0}+ B_0 + { {\bbeta}}_0)$ is $\phi_0$-K\"ahler. 
Therefore, by Kawamata--Viehweg vanishing theorem (Proposition \ref{KVVanishing}), $R^1(\phi_0)_*(\mathcal{O}_{X_0}) = 0$. 
Then, Theorem \ref{ExtendCont} implies that the contraction morphism extends to $\phi:X\to Z/U$ for some $U \subset S$. 
We claim that,  for $t$ near 0, $\phi_t:X_t \to Z_t$ are also  contractions of fiber type. Indeed, by Chevalley's upper semi-continuouity theorem,  $\dim Z_t \le \dim Z_0$. 
Moreover, since $f:X\to S$ is flat, $\dim X_0 = \dim X_t$. Therefore, $$ \dim X_t  - \dim Z_t \ge \dim X_0 - \dim Z_0 >0.$$
By Proposition \ref{DefAmple}, $-(K_{X_t}+B_t +  { {\bbeta}}_t )$ is $\phi_t$-K\"ahler for $t$ near $0$. 
\end{proof}

\begin{remark}
Note that in Proposition \ref{extFano}, the existence of the generalized pair $(X,B+\bbeta )$ is crucial; without it, the deformation stability of Fano type structures remains an open problem (cf. \cite[Question 1.1]{CLZ25}).
\end{remark}

The next result concerns two more building blocks of the MMP: divisorial and flipping contractions under deformations.

\begin{proposition}[Extension of divisorial and flipping contractions]\label{ContInMMP}
Let $f:X\to S$ be a flat proper morphism from a generalized pair $(X,B+ { \boldsymbol{\beta}})$ onto a smooth, connected, and relatively compact curve $S$. Fix a point $0\in S$ and assume that the boundary divisor $B$ does not contain the fiber $X_0$, and that $(X_0,B_0 + { \boldsymbol{\beta}}_0)$ is a generalized klt pair, where $(K_X+X_0+B+\bbeta _X)|_{X_0}=K_{X_0}+B_0+\bbeta _0$. 

If there exists a divisorial/flipping contraction $\phi_0:X_0 \to Z_0$ (cf. Definitions \ref{DefinitionDivisorial} and \ref{DefinitionFlip}), then the divisorial/flipping contraction can be extended to a bimeromorphic contraction $\phi:X\to Z/U$ over a neighborhood $U \subset S$ of $0$. If $\phi _0$ is a flipping contraction, then $\phi$ is small.
\end{proposition}

\begin{proof}
Since $\phi_0 :X_0 \to Z_0$ is a divisorial or flipping contraction, 
$- (K_{X_0} + B_0 +  {{\bbeta}}_0)$ is $\phi_0$-K\"ahler. It follows from Kawamata--Viehweg vanishing (Proposition \ref{KVVanishing}) that 
$R^1 (\phi_0)_*(\mathcal{O}_{X_0}) = 0$. By Theorem \ref{ExtendCont}, after shrinking $S$ to some open neighborhood $U\subset S$, the flipping or divisorial contraction extends to a contraction 
$\phi : X \to Z$ over $U$.

Next, we show that if the central fiber admits a flipping contraction $\phi_0:X_0 \to Z_0$, then it extends to small bimeromorphic morphisms for the nearby fibers. Set
\[
\mathrm{Ex}(\phi_t) := \{ x \in X_t \mid \dim \phi_t^{-1}(\phi_t(x)) > 0 \}.
\]
Since $\phi_0$ is an isomorphism in codimension $2$, 
$\mathrm{codim}_{X_0}(\mathrm{Ex}(\phi_0)) \ge 2$. Consider the restriction of the morphism 
$\mathrm{Ex}(\phi) \to S$. By Chevalley’s upper semi-continuity theorem,
\[
\dim \mathrm{Ex}(\phi)|_{X_t} \le \dim \mathrm{Ex}(\phi)|_{X_0}.
\]
Thus,
\[
\dim \mathrm{Ex}(\phi_t) \le \dim \mathrm{Ex}(\phi_0) \le \dim X_0 - 2 = \dim X_t - 2.
\]
Therefore on the fibers near $0$, $\phi_t : X_t \to Z_t$ are small morphisms and $-(K_{X_t}+B_t+\boldsymbol{\beta} _t)$ is relatively K\"ahler. 

By similar arguments, we can show that if the central fiber is a divisorial contraction, then the nearby fibers are also bimeromorphic contractions (which need not be divisorial contractions, as the example below indicates).
\end{proof}
\begin{example}
In \cite[Remark 4.5]{FH11}, de Fernex--Hacon construct a family 
$ f: X \to \mathbb{A}^1 $ of quadric surfaces degenerating to the Hirzebruch surface 
$ X_0 = \mathbb{F}_2 $. A line in one of the two rulings on the general fiber sweeps out a divisor 
$ R $ on $ X $, which restricts on the central fiber as  
\[
R|_{X_0} = E + F,
\]
where $ F $ is a fiber of $ \mathbb{F}_2 \to \mathbb{P}^1 $ and $ E $ is the $(-2)$-curve on $ \mathbb{F}_2 $. They then construct a $(K_X + \varepsilon R)$-negative contraction $ X \to Z $, which restricts to a divisorial contraction $ X_0 \to Z_0 $ but is an isomorphism 
$ X_t \to Z_t $ for $ t \neq 0 $. This example shows that, in some situations, the deformation of a divisorial contraction need not remain a divisorial contraction on nearby fibers. However, under some reasonable assumptions, \cite[Theorem 5.7]{FH11} shows that a divisorial contraction deforms to a family of divisorial contractions. 
\end{example}


The divisors contracted by a negative contraction are contained in the support of the negative part of the Boucksom--Zariski decomposition.
\begin{lemma}[{\cite[Theorem A.11]{DHY23}}]\label{ZariskiCont}
Let $\phi: X \dashrightarrow X^{\prime}$ be a bimeromorphic contraction of normal compact K\"ahler varieties. Let $\left(X, B+\boldsymbol{\beta}\right)$ and $\left(X^{\prime}, B^{\prime}+\boldsymbol{\beta}'\right)$ be generalized dlt pairs such that $K_X+B+\boldsymbol{\beta}$ is pseudo-effective and $B^{\prime}+\boldsymbol{\beta}'=\phi_*\left(B+\boldsymbol{\beta}\right)$.

If $\phi$ is $(K_X+B+\boldsymbol{\beta})$-negative (which holds for steps of the MMP), then the divisors contracted by $\phi$ are contained in the support of negative part of the Boucksom--Zariski decomposition $N\left(K_X+B+\boldsymbol{\beta}\right)$, that is $\mathrm{Ex}(\phi)\subset \operatorname{Supp} N(K_X+B+ \boldsymbol{\beta})$, and moreover, $$\phi_*(N(K_X+B+ \bbeta)) = N(K_{X'}+ B' + \bbeta ').$$
\end{lemma}
Before proceeding, we recall that if $D=\sum_i d_iD_i$ and $D'=\sum_i d'_iD_i$, for distinct prime divisors $D_i$, then 
$$D \wedge D^{\prime}:=\sum_i \min \left\{d_i, d_i^{\prime}\right\} D_i.$$
\begin{lemma}\label{singular}
In the same setting as Lemma~\ref{ZariskiCont}, if $(X,B+\bbeta)$ is a generalized pair with canonical singularities and 
$B \wedge N(K_X+B+\bbeta)=0$, then $(X',B'+\bbeta')$ also admits canonical singularities and $B' \wedge N(K_{X'}+B'+\bbeta')=0$. 
\end{lemma}

\begin{proof}
Take a common resolution $p:Y \to X, \quad q:Y \to X'$
of the birational map $\phi: X \dashrightarrow X'$. %Since $(X,B+\bbeta)$ has canonical singularities, we have   \[ p^*(K_X+B+\bbeta) = K_s + B_s + \bbeta_s, \] where for any $p$-exceptional divisor $E'$ appearing in $B_s$, the discrepancy satisfies $ a(X,B+\bbeta,E') \ge 0$. 
Since the contraction $\phi$ is $(K_X+B+\bbeta)$-negative, \[
p^*(K_X+B+\bbeta) - q^*(K_{X'}+B'+\bbeta') = E \ge 0
\]
for some effective $q$-exceptional divisor $E$. Since $\phi$ only contracts divisors and does not extract any divisor, every $p$-exceptional divisor is also $q$-exceptional. 

Therefore, by $E \ge 0$, for any  $p$-exceptional divisor $E'$ we obtain  
\[
a(X',B'+\bbeta',E') \;\ge\; a(X,B+\bbeta,E') \;\ge\; 0.
\]
Moreover, since $B \wedge N(K_X+B+\bbeta)=0$, 
Lemma~\ref{ZariskiCont} implies that $\phi$ does not contract any component of $B$. 

For the remaining $q$-exceptional divisors $E''$, we have that $p_*(E'') \neq 0$ is a divisor on $X$. Since $p_*(E'')$ is not contained in $\mathrm{supp}(B)$, it follows that  
\[
a(X',B'+\bbeta',E'') > a(X,B+\bbeta,E'') = 0.
\]
Thus $(X',B'+\bbeta')$ is generalized canonical.  Finally, since  
\[
\phi_* N(K_X+B+\bbeta) = N(K_{X'}+B'+\bbeta'),
\]
and $B' = \phi_*B$, it follows immediately that $B' \wedge N(K_{X'}+B'+\bbeta')=0$. 
\end{proof}


As for the minimal model program in the complex analytic setting, the following Property $(P)$ is useful.
\begin{definition}[{\cite{FujinoBCHM}}]
Let $\pi: X \rightarrow Y$ be a projective morphism of complex analytic spaces and $W$ a compact subset of $Y$ with the properties:\\
\emph{(P1)} $X$ is a normal complex variety,\\
\emph{(P2)} $Y$ is a Stein space,\\
\emph{(P3)} $W$ is a Stein compact subset of $Y$, and\\
\emph{(P4)} $W \cap Z$ has only finitely many connected components for any analytic subset $Z$ which is defined over an open neighborhood of $W$.\\
We say that $\pi : X \to Y$ and $W$ satisfy \emph{Property $(P)$} if the four conditions above hold.
\end{definition}
% \begin{definition}[Canonical models of a generalized K\"ahler pair]\label{goodMM}
% If $(X / S, B+ { \boldsymbol{\beta}})$ is a generalized dlt pair over a  relatively compact analytic variety $S$, then we say that a bimeromorphic map $\phi: X \dashrightarrow X^m$ (proper over $S$) is a \emph{log canonical model} (resp. a \emph{weak log canonical model}) over $S$ if the following conditions hold.\\
% \emph{(1)} $\left(X^m, B^m+ { \boldsymbol{\beta}^m}\right)$ is a  generalized lc pair, where $B^m=\phi_* B+E$, and $E$ is the reduced sum of all $\phi^{-1}$-exceptional divisors,\\
% \emph{(2)} $K_{X^m}+B^m+ { \boldsymbol{\beta}}_{X^m}$ is K\"ahler over $S$ (resp. nef over $S$), and\\
% \emph{(3)} $a(P, X, B,  { \boldsymbol{\beta}}) \leq a\left(P, X^m, B^m,  { \boldsymbol{\beta}^m}\right)$ for every $\phi$-exceptional divisor $P$. 
% \end{definition}

\begin{proposition}[{Extension of MMP to nearby fibers, \cite[Theorem 2]{DefGeneralType}}]\label{ExtendMMP}
Let $g: X \rightarrow S$ be a flat proper morphism with connected fibers from a generalized pair $(X,B+ \bbeta)$ to a smooth, connected, and relatively compact curve. Fix a point $0\in S$ and assume that the support of the boundary divisor $B$ does not contain the fiber $X_0$. Suppose that $(X_0,B_0 + { \boldsymbol{\beta}}_0)$ is projective, with canonical singularities, and $\lfloor B_0 \rfloor = 0$, where $$(K_X+X_0+B+\bbeta _X)|_{X_0}=K_{X_0}+B_0+\bbeta _0,$$ and the negative part of Boucksom--Zariski decomposition satisfies the relation $$N(K_{X_0}+ B_0 + \bbeta_0) \wedge B_0 = 0.$$ Then every sequence of transcendental MMP-steps \[(X_0, B_0  +\bbeta_0) \dashrightarrow (X_0^{(1)}, B_0^{(1) }+  \bbeta_0^{(1)}) \dashrightarrow (X_0^{(2)},B_0^{(2)}+ \bbeta_0^{(2)}) \dashrightarrow \cdots\]  extends to a sequence of $(K_X+B+\bbeta _X)$-negative proper meromorphic maps 
$$ (X, B+ \bbeta)/U \dashrightarrow (X^{(1)},B^{(1)}+ \bbeta^{(1)})/U \dashrightarrow (X^{(2)},B^{(2)}+ \bbeta^{(2)}) /U\dashrightarrow \cdots,$$over some open neighborhood $U\subset S$ of $0$.
\end{proposition}
\begin{proof}
The proof is based on \cite[Theorem 2]{DefGeneralType} and \cite[Theorem 5.12]{DH24}. Note that $({X_0},B_0)$ is also canonical and so $K_X+B$ is $\R$-Cartier and canonical on a neighborhood of $X_0$ by \cite[Proposition 3.5]{FH11}. Assume that $\phi_0:X_0 \to Z_0$ is a $(K_{X_0}+ B_0 + \bbeta_0)$-negative extremal contraction. Then, by Kawamata--Viehweg vanishing (Proposition \ref{KVVanishing}), $R^1 (\phi_0)_* \mathcal{O}_{X_0} = 0$. By Theorem \ref{ExtendCont}, $\phi_0$ extends to a contraction $\phi: X \to Z$. Note that by \cite{DH24MMP}, there exists a klt pair $(X_0,\Delta _0)$ such that $K_{X_0}+\Delta  _0\equiv K_{X_0}+B_0+\bbeta _0$ and an ample divisor $H_0$ such that $K_{X_0}+\Delta  _0+H_0$ is semi-ample and defines the contraction $X_0\to Z_0$. It follows that $$K_{Z_0}+\Delta  _{Z_0}+H_{Z_0}:=\phi _{0,*}(K_{X_0}+\Delta  _0+H_0)$$ is klt and hence $Z_0$ has rational singularities. But then $Z$ also has rational singularities on a neighborhood of $Z_0$ by \cite{Elkikj}. 

Suppose that $\phi _0$ is of flipping type. %If $K_{X_0}+B_0\not \equiv _{Z_0} 0$, then $K_{X_0}+B_0$ is relatively ample or relatively anti-ample and hence so is $K_X+B$ (after possibly shrinking $U$). In particular we may assume that $\phi$ is projective. Since $\rho (X_0/Z_0)=1$, we have $\bbeta _0\equiv _{Z_0} \lambda (K_{X_0}+Z_0)$. 
We will construct  $X^+\to Z$ locally over $Z$. Pick a point $z\in Z_0$ and fix a neighborhood $z\in W\subset Z$ such that $W$ is relatively compact, Stein and satisfies Property (P). %Working locally over $W$
Let $\nu :X'_W\to X_W$ be a resolution, projective over $W$ such that $X'_z:=X'_W\times _W\{z\}$ is a divisor with simple normal crossings and $\bbeta _{X'_W}$ is nef over $W$. Arguing as in the proof of \cite[Theorem 5.12]{DH24} we may assume that there is a divisor $D'_W\equiv_W \bbeta _{X'_W}$ and hence $D'_W$ is nef and big over $W$. We may then assume that $(X'_W,B'_W+D'_W)$ is sub-klt and hence $(X_W, B_W+D_W\equiv \nu _*( B'_W+D'_W))$ is klt and in fact canonical.
By \cite[Theorem C]{FujinoBCHM},  $({X_W},B_W+D_W)$ has a log canonical model over $W$ say $X^+_W$. Note that these log canonical models are defined locally and uniquely determined by the class of the $\R$-divisor $$K_{X_W}+B_W+D_W\equiv _W K_{X_W}+B_W+\bbeta _{X_W}.$$ It follows that the $X^+_W$ glue together to give $\phi ^+:X^+\to Z$.
    
Finally, we observe that $\phi ^+_0=X^+_0\to Z_0$ is the flip of $\phi _0:X_0\to Z_0$. This can be checked locally over $Z_0$. Let $W_0=Z_0\cap W$, then \[K_{X_{W_0}}+B_{X_{W_0}}+\bbeta |_{X_{W_0}}\equiv K_{X_{W_0}}+B_{X_{W_0}}+D_{X_{W_0}}\equiv  (K_{X_{W}}+B_{X_{W}}+X_{W}+D_{X_W})|_{X_{W_0}}.\] 
Note that since $({X_{W_0}},B_{X_{W_0}}+D_{X_{W_0}})$ is canonical, then $X_{W_0}\dasharrow X^+_{W_0}$ is a small birational map since otherwise there is a divisor $F$ on $X_{W_0}^+$ that is exceptional over $X_{W_0}$. But then, letting $$K_{X_{W_0}^+}+B_{X^+_{W_0}}+D_{X^+_{W_0}}=(K_{X_{W}^+}+B_{X^+_{W}}+D_{X^+_{W}})|_{X^+_{W_0}}$$ we have the following inequalities between discrepancies \[0\ge a(X_{W_0}^+,B_{X^+_{W_0}}+D_{X^+_{W_0}};F)>a(X_{W_0},B_{X_{W_0}}+D_{X_{W_0}};F)\ge 0\] which is impossible. 
Note that the first inequality is an equality unless $F$ is contained in the support of $B_{X^+_{W_0}}+D_{X^+_{W_0}}$, the second strict inequality follows by the monotonicity lemma as in \cite[Lemma 3.38]{KM98}, and the last inequality is immediate as $(X_{W_0},B_{X_{W_0}}+D_{X_{W_0}})$ is canonical. Thus $X_{W_0}\dasharrow X_{W_0}^+$ is a small birational map and $$K_{X_{W_0}^+}+B_{X^+_{W_0}}+D_{X^+_{W_0}}=(K_{X_{W}^+}+B_{X^+_{W}}+X^+_{W}+D_{X^+_{W}})|_{X^+_{W_0}}$$ is ample over $W_0$ so that $X_{W_0}\dasharrow X_{W_0}^+$ is the $(K_{X_{W_0}}+B_{X_{W_0}}+D_{X_{W_0}})$-flip.
    
One easily checks that $(X_0^+, B_0^+ + \bbeta _{X^+_0})$ is projective with  canonical singularities, and the negative part of the divisorial Zariski decomposition satisfies the relation $$N (K_{X_0^+} + B_0^+ + \bbeta _{X^+_0}) \wedge B_0^+  = 0.$$ 


Suppose now that $\phi _0$ is a divisorial contraction and so $K_{Z_0}+B_{Z_0}+\bbeta _{Z_0}$ is canonical and $K_{Z_0}+B_{Z_0}$ is $\R$-Cartier. As observed above $Z_0$ has rational singularities and hence we may assume that $Z$ has rational singularities (after shrinking in a neighborhood of $Z_0$). Note that it also follows that $(Z_0,B_{Z_0})$ has canonical singularities and hence $(Z,B_Z)$ has canonical singularities \cite[Proposition 3.5]{FH11}.
Arguing as above we may construct $X\dasharrow Z^c$ the log canonical model for $(X,B+\bbeta _X)$ over $Z$.

It follows that $$K_{Z_0^c}+B_{Z^c_0}+\bbeta _{Z^c_0}=(K_{Z^c}+B_{Z^c}+Z^c_0+\bbeta _{Z^c})|_{Z^c_0}$$ is K\"ahler over $Z$ and hence $Z^c_0\to Z_0$ is an isomorphism. Therefore, $Z^c\to Z$ is also an isomorphism in a neighborhood of $Z_0$, and so $(Z,B_Z+\bbeta)$ is a generalized canonical pair and $-(K_X+B+\bbeta _X)$ is K\"ahler over $Z$. 

If we let $X^+:=Z$, then one easily checks that $(X_0^+, B_0^+ + \bbeta _{X^+_0})$ is projective with 
canonical singularities, and the negative part of divisorial Zariski
decomposition satisfies the relation $N (K_{X_0^+} + B_0^+ + \bbeta _{X^+_0}) \wedge B_0^+  = 0$.  

Replacing $(X,B+\bbeta)$ by $(X^+,B^++\bbeta^+)$ and proceeding by induction, we conclude the proof. Note that $X^+$ is K\"ahler by Lemma \ref{ContKahler}.
\end{proof}


\begin{remark}\label{Remark:ExtMMP}
Thanks to the contraction theorem for K\"ahler threefolds (cf. \cite[Theorem 1.2]{DH24}), the same argument shows that Proposition~\ref{ExtendMMP} also holds for families of K\"ahler threefolds, without assuming that the central fiber is projective.
\end{remark}

\begin{lemma}[{Contractions in the K\"ahler MMP preserve the K\"ahler condition, \cite[Lemma 2.24]{DHY23}}]\label{ContKahler}
Let $\pi: X \rightarrow S$ be a proper surjective morphism of compact complex analytic varieties such that $X$ is K\"ahler. If $(X, B+ { \boldsymbol{\beta}})$ is a generalized dlt pair and $\phi$ : $X \dashrightarrow X^{\prime}/S$ is a $(K_X+B+{ \boldsymbol{\beta}})$-flip, flipping contraction or divisorial contraction over $S$, then $X^{\prime}$ is K\"ahler over $S$. 
\end{lemma}


For a projective generalized klt pair $(X, B + \bbeta)$, \cite{DH24MMP} proved the following transcendental version of \cite{BirkarCHM}, which will be repeatedly used in what follows. 
\begin{theorem}[{\cite[Theorem 1.2]{DH24MMP}}]\label{Theorem: TransBCHM}
Let $(X, B+\boldsymbol{\beta})$ be a compact generalized klt pair, and $f:X\to S$ a morphism of projective varieties such that $B+\boldsymbol{\beta}_X$ is big over $S$. 
%Assume that $X$ is a %$\mathbb{Q}$-factorialprojective variety over $\mathbb{C}$. 
Then the following hold:\\
(1) if $K_X+B+\boldsymbol{\beta}_X$ is pseudo-effective over $S$, then $(X, B+\boldsymbol{\beta})$ has a good minimal model over $S$, and\\
(2) if $K_X+B+\boldsymbol{\beta}_X$ is not pseudo-effective over $S$, then $(X, B+\boldsymbol{\beta})$ has a Mori fiber space over $S$.
\end{theorem}
\begin{proof}
The above result holds when $\dim S=0$ and $X$ is $\Q$-factorial by \cite[Theorem 1.2]{DH24MMP}. Assume now that $X$ is $\Q$-factorial and $\dim S>0$. Let $H$ be an ample divisor on $S$ and $A\in |(2\dim X+1)H|$ a general member. Replacing $B$ by $B+f^*A$, one sees that every $(K_X+B+\bbeta _X)$-negative extremal ray is a $(K_X+B+\bbeta _X)$-negative extremal ray over $S$ (cf. \cite[Section 3]{HP24}). It follows that $K_X+B+\bbeta _X$ is nef if and only if it is nef over $S$ and every step of the $(K_X+B+\bbeta _X)$-MMP is a step of the $(K_X+B+\bbeta _X)$-MMP over $S$.
Thus the claim holds in this case.
    
 If $X$ is not $\Q$-factorial, then the $\Q$-factorialization $\nu:X'\to X$ is a small birational morphism such that $X'$ is $\Q$-factorial.  (The existence of the $\Q$-factorialization $\nu$ follows by standard arguments using the $\Q$-factorial case and following the proof of \cite[Theorem 1.5]{HP24}).  By {\cite[Theorem 1.2]{DH24MMP}}, (1) and (2) hold for $K_{X'}+B'+\bbeta _{X'}:=\nu ^*(K_X+B+\bbeta _X)$ and hence (1) and (2) also hold for $K_X+B+\bbeta _X$.
\end{proof}

\subsection{Stability of nefness and bigness of adjoint classes}
Das--Hacon (cf. \cite[Conjecture 1.6]{DH24} and also \cite{TosatiKAWA}) made the following transcendental base point free conjecture. 
\begin{conjecture}[{Transcendental base point free conjecture}]\label{Transbpf}
Let $(X, B+\bbeta)$ be a generalized klt K\"ahler pair, where $X$ is a normal compact K\"ahler variety, and $\alpha \in H_{\mathrm{BC}}^{1,1}(X,\mathbb{R})$ a nef class on $X$. If $\alpha-\left(K_X+B+ \bbeta_X\right)$ is nef and big, then there exists a proper surjective morphism with connected fibers $f: X \rightarrow Y$ to a normal compact K\"ahler variety $Y$ with rational singularities and a K\"ahler class $\omega_Y \in H_{\mathrm{BC}}^{1,1}(Y,\mathbb{R})$ such that $\alpha=f^* \omega_Y$. 
\end{conjecture}

Thanks to \cite{DH24MMP}, the transcendental base point free conjecture holds when $X$ is projective.

\begin{theorem}[{\cite[Corollary~3.7]{DH24MMP}}]\label{DHTransbpf}
Let $(X, B+ \bbeta)$ be a projective generalized klt pair and $\alpha \in H_{\mathrm{BC}}^{1,1}(X,\mathbb{R})$ a nef class. If $\alpha-\left(K_X+B+ \bbeta_X\right)$ is nef and big, then there is a proper surjective morphism $f: X \rightarrow Z$ with connected fibers to a normal projective variety $Z$ with rational singularities such that $\alpha=f^* \omega_Z$ for some K\"ahler class $\omega_Z$ on $Z$.
\end{theorem}

Theorem \ref{DHTransbpf} allows us to prove the deformation stability of bigness-and-nefness of adjoint classes in K\"ahler families with projective central fiber (Proposition \ref{Extbigandnef}). Before that, let us present several lemmas.
\begin{lemma}\label{Lemma: cohomologicaltrivial}
Let $X$ be a proper Moishezon variety with rational singularities, with $L$ a holomorphic $\mathbb{R}$-line bundle on $X$, and a class $\alpha \in H^{1,1}_{\rm{BC}}(X, \mathbb{R})$. If $L\equiv 0$ (resp. $\alpha \equiv 0$), then $c_1(L) = 0 \in H^2(X ,\mathbb{R})$ (resp. $\alpha  =0 \in H^2(X, \mathbb{R})$) . 
\end{lemma}
\begin{proof}
First we treat the case of a line bundle. Since $X$ is Moishezon, there exists a bimeromorphic modification $\mu : X^{p} \to X$ with $X^{p}$ a projective manifold. The pull-back line bundle $\mu^{*}L$ remains numerically trivial, i.e., $\mu^{*}L \equiv 0$. By Poincar\'e duality and Hodge--Riemann bilinear relationship the pairing $${\langle -,-\rangle}:H^{1,1}(X^p,\R)\times H^{n-1,n-1}(X^p,\R) \to H^{2 n}(X^p, \mathbb{R})\cong\R,\quad (\alpha, \beta)\mapsto \int_{X^p} \alpha \cup \beta$$is a perfect pairing. Since $\mu^* L \equiv 0$, we have $$\mu^* L \cdot C = {\langle c_1(\mu^* L), [ C]\rangle}= 0$$for any irreducible curve $C$, where $[C]\in H^{n-1,n-1}(X^p, \R)$ represents its fundamental class. By the Lefschetz (1,1)-theorem, $H^{n-1,n-1}(X^p,\R)$ is generated by these classes $[C]$, therefore using the perfect pairing condition, we have $\mu^* c_1(L) = c_1(\mu^* L) = 0$. Since $X$ has rational singularities, the pullback map
\[
\mu^{*} : H^{2}(X,\mathbb{R}) \to H^{2}(X^{p},\mathbb{R})
\]
is injective by \cite[Lemma~2.6]{DH24MMP}. Hence $c_{1}(L)=0 \in H^{2}(X,\mathbb{R})$ as well.

Next we handle the case of a Bott--Chern class $\alpha \in H^{1,1}_{\mathrm{BC}}(X,\mathbb{R})$. Since $\mu^{*}\alpha \in H^{1,1}_{\mathrm{BC}}(X^{p},\mathbb{R})$, the Lefschetz $(1,1)$-theorem implies that $\mu^{*}\alpha$ is the first Chern class of some $\mathbb R$-divisor on $X^{p}$. Hence the same argument as above applies, and the conclusion follows.
\end{proof}
\begin{lemma}\label{l-bir}
Let $f:U\to V$ be a proper Moishezon morphism of analytic varieties, and $\alpha$ a locally exact $(1,1)$-form on $U$. If $\alpha |_{f^{-1}(v)}\equiv 0$ for some point $v\in V$, then there exists a neighborhood $V'\subset V$ around $v$, such that $\alpha |_{f^{-1}(w)}\equiv 0$ for all $w\in V'$.

If furthermore $f:U\to V$ is bimeromorphic, and $\alpha \equiv K_U+B+\bbeta _U$ for some generalized klt pair $(U,B+\bbeta )$, then replacing $V$ by $V'$ and $U$ by $f^{-1}(V')$, we have that $
K_V+B_V+\bbeta _V:=f_*(K_U+B+\bbeta _U)$
is generalized klt and $K_U+B+\bbeta _U=f^*(K_V+B_V+\bbeta _V)$.
\end{lemma}
\begin{proof}
Replacing $U$ by a resolution, we may assume that $U$ is smooth and $f$ is projective. Since there are countably many families of $f$-vertical curves, $\alpha |_{f^{-1}(w)}\equiv 0$ for very general $w\in W\subset V$ (here $W$ is a complement of countably many closed subsets not containing $v$). Let $V_j$ be a stratification of $V$ by locally closed subsets such that $f$ is smooth over each $V_j$ and $U_j=f^{-1}(V_j)$, then $R^2f_*\mathbb R _{U_j}$ is a locally constant sheaf whose fibers are $(R^2f_*\mathbb R _{U_j})_p=H^2(f^{-1} (p),\mathbb R )$ for $p\in V_j$. Suppose that $v\in \bar V_j$. Since each $V_j$ is locally closed, we have a non-empty intersection $W\cap V_j \ne \emptyset$. For any $w\in W \cap V_j$ we have $\alpha|_{f^{-1}(w)} \equiv 0$. Since $R^2 f_* \mathbb{R}_{U_j}$ is locally constant, this means $\alpha|_{f^{-1}(p)} =0 $ for all $p\in V_j$ by Lemma \ref{Lemma: cohomologicaltrivial}. 
Let $V'=\cup _{v\in \overline V_j}V_j$, then $\alpha |_{f^{-1}(w)}\equiv 0$ for all $w\in V'$.

Suppose now that  $\alpha \equiv K_U+B+\bbeta _U$ for some generalized klt pair $(U,B+\bbeta )$. Replacing $V$ by $V'$ and $U$ by $f^{-1}(V')$, we may assume that $\alpha |_{f^{-1}(w)}\equiv 0$ for all $w\in V$.
Let $\nu :U'\to U$ be a resolution such that $f':U'\to V$ is projective and $\bbeta _{U'}$ is nef. 
We write $K_{U'}+B_{U'}+\bbeta _{U'}=\nu ^*(K_U+B+\bbeta _U)$. 
Working locally over $V$, we may assume that $V$ is Stein and $\bbeta _{U'}\equiv -(K_{U'}+B_{U'})$ is nef and big over $V$. But then \[B_{U'}^{\ge 0}+\varepsilon {\rm Ex}(\nu )+\bbeta _{U'}\equiv \Delta _{U'},\] where $K_{U'}+\Delta _{U'}\equiv _V B_{U'}^{\leq 0}+\varepsilon {\rm Ex}(\nu )=:F$ is klt. By the mmp for projective morphisms (\cite{FujinoBCHM} and \cite{DHP24}), $K_{U'}+\Delta _{U'}$ has a good minimal model $\phi:U'\dasharrow U''$ over $U$, as in the commutative diagram:
\begin{center}
\begin{tikzcd}
	{U'} && {U''} \\
	& U \\
	& V.
	\arrow["\phi", dashed, from=1-1, to=1-3]
	\arrow["\nu", from=1-1, to=2-2]
	\arrow["{f'}"', from=1-1, to=3-2]
	\arrow["\mu"', from=1-3, to=2-2]
	\arrow["{f''}", from=1-3, to=3-2]
	\arrow["f"{description}, from=2-2, to=3-2]
\end{tikzcd}
\end{center}
Since $K_{U'}+\Delta  _{U'}\equiv _{U}F$ where $F$ is an effective $\mathbb R$-divisor with support equal to ${\rm Ex}(\nu)$, the negativity lemma gives $\phi _*F=0$, and so $\mu:U''\to U$ is a small birational morphism. Thus $K_{U''}+\Delta _{U''}$ is a klt pair, numerically trivial over $V$. By the base point free theorem, $K_{U''}+\Delta _{U''}$ is semi-ample over $V$, and hence $K_{U''}+\Delta _{U''}=(f'')^*(K_V+\Delta _V)$ where $(V,\Delta _V)$ is klt and in particular $V$ has rational singularities. By \cite[Lemma 3.3]{HP16} or \cite[Lemma 2.6]{DH24MMP}, $\alpha =f^*\alpha _V$ for some locally exact $(1,1)$-form $\alpha _V$ on $V$. But then $K_V+B_V+\bbeta _V:=f_*(K_U+B+\bbeta _U)$ is generalized klt. 
\end{proof}
Now we can prove:
\begin{proposition}\label{Extbigandnef}
Let $f:X\to S$ be a proper surjective morphism from a %normal $\mathbb{Q}$-factorial 
generalized K\"ahler pair $(X,B +  { \boldsymbol{\beta}})$ onto a smooth, connected, relatively compact curve $S$. Fix a point $0\in S$ and assume that the support of the boundary divisor $B$ does not contain the fiber $X_0$. 

Assume that the restriction to the central fiber $(X_0,B_0 +  { \boldsymbol{\beta}}_0)$ is a projective  generalized klt pair, such that $K_{X_0}+B_0 +  { {\bbeta}}_0$ is nef and big. Then $K_{X}  + B +  { {\bbeta}_X}$ is nef and big over $U$ for some Euclidean open neighborhood $U\subset S$ of $0$. 
\end{proposition}
\begin{proof}
Since $(X_0, B_0 + \bbeta_0)$ is a generalized klt pair, by Proposition \ref{DefKLT}, after shrinking the base, $(X,B + \bbeta)$ also has generalized klt singularities. 
By the transcendental base point free Theorem \ref{DHTransbpf} and semiample fibration, we know that $K_{X_0} + B_0 +  { {\bbeta}}_0$ is a semi-ample class and there exists a bimeromorphic contraction morphism  $$g_0: X_0 \to Y_0,$$ with a generalized klt pair $(Y_0, B_{Y_0} +  { \boldsymbol{\beta}}_0)$, such that $$g_0^* (\omega_0) = K_{X_0 } +B_0+  { {\bbeta}}_0,$$ where $\omega_0$ is a K\"ahler class on $Y_0$. As in the proof of Lemma \ref{l-bir}, there exists a klt pair $(Y_0, \Delta_{Y_0})$ such that $K_{Y_0} + \Delta_{Y_0}$ is ample, and some class $\delta_{Y_0} \in H^{1,1}_{\rm{BC}}(Y_0 , \mathbb{R})$ such that $(K_{Y_0}, \Delta_{Y_0} + {\delta}_{Y_0})$ is a generalized klt pair with $\omega_0 = K_{Y_0}+ \Delta_{Y_0} +\delta_{Y_0}$. 

By Kawamata--Viehweg vanishing (Proposition \ref{KVVanishing}), $R^i (g_0)_* \mathcal{O}_{X_0} = 0$, for $i>0$. Then, by Theorem \ref{ExtendCont}, we know that the contraction extends to a relative contraction $$g:X\to Y/S.$$
We denote $B_Y = g_* B$ and $ { {\bbeta}}_Y = g_*  { {\bbeta}_X}$. By a similar upper semi-continuity argument as in Proposition \ref{ContInMMP}, we know that the extension $g:X\to Y/S$ is a bimeromorphic contraction, which is fiberwise bimeromorphic by \cite[Corollary 3.2]{CRC25}. 
Since $g_0^* \omega_0 = K_{X_0}+ B_0 +  { {\bbeta}}_0$, $g_0$ is a $(K_{X_0}+ B_0+  { {\bbeta}}_0)$-trivial contraction morphism. By Lemma \ref{l-bir}, after shrinking the base,  $g$ is a $(K_X+ B+ { {\bbeta}})$-numerically trivial contraction morphism. 
By \cite[Lemma 2.6]{DH24MMP}, it follows that there exists a class $\omega\in H^{1,1}_{\rm{BC}}(Y,\R)$ such that  $$(K_X+B + { {\bbeta}_X} ) = g^* \omega.$$
Since $\omega_{0}=\omega|_{Y_0}$ is K\"ahler, $\omega$ is K\"ahler near $Y_0$. So $K_X+B+ \bbeta_X$ is nef and big over $U$ for some Euclidean open neighborhood $U\subset S$ of $0$. 
\end{proof}

We propose the following more general statement. 
\begin{conjecture}\label{conj: bignef}
Let $f:X\to S$ be a proper surjective morphism from %a normal $\mathbb{Q}$-factorial 
generalized K\"ahler pair $(X,B +  { \boldsymbol{\beta}})$ onto a smooth, connected, relatively compact curve $S$. Fix a point $0\in S$ and assume that the support of the boundary divisor $B$ does not contain the fiber $X_0$. 

Assume that the restriction to the central fiber $(X_0,B_0 +  { \boldsymbol{\beta}}_0)$ is a generalized klt pair, such that $K_{X_0}+B_0 +  { {\bbeta}}_0$ is nef and big. Then $K_{X}  + B +  { {\bbeta}_X}$ is nef and big over $U$ for some open neighborhood $U\subset S$ of $0$. 
\end{conjecture}

\begin{remark}\label{Extbigandnef3fold}
Following the same proof as in Proposition  \ref{Extbigandnef}, one can show that the conjecture above follows from the transcendental base point free Conjecture \ref{Transbpf}. 

Since the transcendental base point free conjecture has been proved for K\"ahler threefolds (cf.~\cite{DH24}) and for cases where the numerical Kodaira dimension $\mathrm{nd}(K_X + \alpha) \le 3$ (cf.~\cite{JunSheng}), Conjecture \ref{conj: bignef} holds in these two cases.
\end{remark}


%%%%%%%%%%%%%%%%%%%%%%%%%%%%%%%%%%%%%%%%%%%%%%%%%%%%%%%%%%%%%%%%%%%%%%%%%%%%%%%%%%%%%%%%%%%%%%%%%%%%%%%%%%%%%%%%%%%%%%%%%%%%%%%%%
%%%%%%%%%%%%%%%%%%%%%%%%%%%%% Section Deformation of Pseudo-effectivity  %%%%%%%%%%%%%%%%%%%%%%%%%%%%%%%%%%%%%
%%%%%%%%%%%%%%%%%%%%%%%%%%%%%%%%%%%%%%%%%%%%%%%%%%%%%%%%%%%%%%%%%%%%%%%%%%%%%%%%%%%%%%%%%%%%%%%%%%%%%%%%%%%%%%%%%%%%%%%%%%%%%%%%%
\section{Pseudo-effectivity of canonical divisors under deformations}\label{Section: Deformation pseff}
The goal of this section is to study the global stability of the pseudo-effectivity of canonical divisors in K\"ahler families, under the assumption that the central fiber is projective with canonical singularities. Our first result involves the deformation closedness of pseudo-effectivity of adjoint classes.

%%%%%%%%%%%%%%%%%%%%%%%%%%%%%%%%%%%%%%%%%%%%%%%%%%%%%%%%%%%%%%%%%%%%%%%%%%%%%%%%%%%%%%%%%%%%%%%%%%%%%%%%%%%%%%%%%%%%%%%%%%%%%%%%%
%%%%%%%%%%%%%%%%%%%%%%%%%%%%%%%%%%%%%%%%%%%%%%%%%%%% Subsection Positivity in Kahler family %%%%%%%%%%%%%%%%%%%%%%%%%%%%%%%%%%%%%
%%%%%%%%%%%%%%%%%%%%%%%%%%%%%%%%%%%%%%%%%%%%%%%%%%%%%%%%%%%%%%%%%%%%%%%%%%%%%%%%%%%%%%%%%%%%%%%%%%%%%%%%%%%%%%%%%%%%%%%%%%%%%%%%%
%\subsection{Deformation closedness of pseudo-effective adjoint classes}

\begin{proposition}\label{Def-Close-Pseff}
Let $f: X \to S$ be a family from a generalized K\"ahler pair $(X, B+  { \boldsymbol{\beta}})$ onto a smooth, connected, and relatively compact curve $S$. Fix a point $0\in S$ and assume that the boundary divisor $B$ does not contain the central fiber $X_0$. 
Suppose that $(X_0,B_0 + { \boldsymbol{\beta}}_0)$ is projective, with canonical singularities, where $(K_X+X_0+B+\bbeta _X)|_{X_0}=K_{X_0}+B_0+\bbeta _0$, and $\lfloor B_0 \rfloor = 0$ with $N(K_{X_0 }+B_0 + \bbeta_0)\wedge B_0  = 0$. 

If there exists a sequence $\{t_i\}_{i\in \mathbb N}\subset S$ converging to $0$ such that $K_{X_{t_i}}+B_{t_i}+\bbeta _{t_i}$  is pseudo-effective for each $i$, then $K_{X_0} + B_0+  { \bbeta}_0$ is pseudo-effective.
\end{proposition}

\begin{proof}
We will argue by contradiction. Assume that $K_{X_0}+B_0+ { \bbeta}_0$ is not pseudo-effective. Theorem \ref{Theorem: TransBCHM} shows that the $(K_{X_0}+ B_0 + { \bbeta}_0)$-MMP with scaling will terminate with some Mori fiber space $X_0' \to Z_0$. By Proposition \ref{ExtendMMP}, there exists a neighborhood $U \subset S$ of $0$ such that that steps $(K_{X_0}+ B_0 + \bbeta_0)$-MMP will extend to a sequence of $(K_X+B+\bbeta)$-negative proper meromorphic maps over $U$
$$(X,B+ \bbeta)/U \dashrightarrow (X^{(1)}, B^{(1)} + \bbeta^{(1)})/U \dashrightarrow \cdots \dashrightarrow (X', B' + \bbeta')/U \dashrightarrow \cdots.$$
By Proposition \ref{extFano}, after shrinking $U$ the Mori fiber space $X'_0 \to Z_0$ extends to a relative Mori fiber space $X'\to Z/U$ (and thus the relative $(K_{X}+ B + \bbeta)$-MMP over $U$ terminates at $X'$). 
As $X'_{t_i} \to Z_{t_i}$ ($t_i \to 0$) is a Mori fiber space, Proposition \ref{NonPSEFFMFS} yields that the divisor
$K_{X'_{t_i}} + B'_{t_i} + \bbeta'_{t_i}$
is not pseudo-effective for $i \gg 1$. 
Since $K_{X_{t_i}}+B_{t_i} + {\bbeta}_{t_i}$ is pseudo-effective for a sequence of $t_i \to 0$ by assumption, Proposition \ref{PushPSEF} implies that $K_{X'_{t_i}}+B'_{t_i} +  { \bbeta}_{t_i}'$ is pseudo-effective as well, which is a contradiction.
\end{proof}


%%%%%%%%%%%%%%%%%%%%%%%%%%%%%%%%%%%%%%%%%%%%%%%%%%%%%%%%%%%%%%%%%%%%%%%%%%%%%%%%%%%%%%%%%%%%%%%%%%%%%%%%%%%%%%%%%%%%%%%%%%%%%%%%%
%%%%%%%%%%%%%%%%%%%%%%%%%%%%%% Subsection Deformation Openness of Pseudo-Effective Adjoint Classes%%%%%%%%%%%%%%%%%%%%%%%%%%%%%%%
%%%%%%%%%%%%%%%%%%%%%%%%%%%%%%%%%%%%%%%%%%%%%%%%%%%%%%%%%%%%%%%%%%%%%%%%%%%%%%%%%%%%%%%%%%%%%%%%%%%%%%%%%%%%%%%%%%%%%%%%%%%%%%%%%
%\subsection{Deformation openness of pseudo-effective s}
We next prove the deformation openness of the pseudo-effectivity of canonical divisors. At present, we can establish this result only for canonical divisors, rather than for adjoint classes, due to the absence of an analogue of Proposition~\ref{pseffthreshold} for adjoint classes.

\begin{theorem}\label{Deformation invariance of pseudo-effectivity}
Let $f: X \to S$ be a family from a  K\"ahler variety $X$ onto a smooth, connected, and relatively compact curve $S$. Fix a point $0\in S$ and assume that $X_0$ is projective with canonical singularities. Then the canonical divisor $K_{X_0}$ of $X_0$ is pseudo-effective if and only if there exists a Euclidean open neighborhood $U \subset S$ of $0$ such that $K_{X}$ is relatively pseudo-effective over $U$.
\end{theorem}
\begin{proof}By Proposition \ref{DefCanSing}, we may assume that ${X_t}$ and $X$ admit canonical singularities. Suppose that $K_{X_0}$ is pseudo-effective, then we must prove that $K_X$ is relatively pseudo-effective over some open neighborhood $U\subset S$. By contradiction, suppose that $K_X $ is not $f$-pseudo-effective over any (sufficiently small) open neighborhood $U\subset S$ of $0$. For a K\"ahler class $\omega$ on $X$, define the pseudo-effective threshold  $$\tau':= \inf\{t>0 \mid K_X+ t \omega \text{ is }f\text{-pseudo-effective over some neighborhood }U\subset S\}.$$ 
Since $K_X$ is not $f$-pseudo-effective over any neighborhood $U \subset S$, $\tau '>0$ by Proposition \ref{pseffthreshold}. 

Let $0< \tau < \tau '$, and consider the pair $(X_0,  \tau \boldsymbol{\omega}_{X_0})$ where $\boldsymbol{\omega}_{X}=\overline \omega$. By assumption $K_{X_0} $ is pseudo-effective, so that $K_{X_0} + \tau {{\omega}}_{X_0}$ is big. Since $X_0$ is projective, by Theorem \ref{Theorem: TransBCHM}, we can run the transcendental MMP on the central fiber. This MMP will terminate in $(X_0^{(n)} ,   \tau \boldsymbol{\omega}_{X_0^{(n)}})$ with nef $K_{X_0^{(n)} }+  \tau {\omega}_{X^{(n)}_0}$ 
$$(X_0,   \tau \boldsymbol{\omega}_{X_0}) \dashrightarrow (X^{(1)}_0 ,  \tau \boldsymbol{\omega}_{X_0^{(1)}}) \dashrightarrow \cdots \dashrightarrow (X^{(n)}_0,  \tau \boldsymbol{\omega}_{X_0^{(n)}}).$$ 
By Proposition \ref{ExtendMMP}, there exists a neighborhood $0\in U\subset S$, such that the $(K_{X_0}+ \tau {\omega}_{X_0})$-MMP extends to a proper $(K_X  +\tau{{\omega}}_X)$-negative map over $U$
$$(X, \tau \boldsymbol{\omega}_X) /U\dashrightarrow (X^{(1)},  \tau\boldsymbol{ \omega}_{X^{(1)}}) /U\dashrightarrow \cdots \dashrightarrow (X^{(n)}, \tau \boldsymbol{\omega}_{X^{(n)}})/U\dashrightarrow \cdots.$$
By Proposition \ref{PushPSEF}, $K_{X_0^{(n)}}  + \tau {\omega}_{X_0^{(n)}}$ is big and nef. By Proposition \ref{Extbigandnef}, we know that after further shrinking $U$, $K_{X^{(n)}}  + \tau {\omega}_{X^{(n)}}$ is nef (and hence also relatively pseudo-effective) over $U$. By Proposition \ref{PushPSEF}, $K_{X} + \tau {\omega}_X$ is relatively pseudo-effective over $U$, and hence $\tau' \le \tau < \tau'$, which is a contradiction. 

Finally, suppose that $K_X$ is relatively pseudo-effective over some open neighborhood $U\subset S$ of $0$, then it follows from Proposition \ref{Def-Close-Pseff} that $K_{X_0}$ is pseudo-effective.
\end{proof}


\begin{proposition}\label{pseflocalglobal}
Let $f:X\to S$ be a family from a K\"ahler variety onto a smooth, connected, and relatively compact curve $S$. Assume that $X_t$ admits canonical singularities for any $t\in S$ and there exists a Euclidean open subset $U\subset S $ such that $K_X$ is relatively pseudo-effective over $U$. Then $K_X$ is relatively pseudo-effective over $S$.
\end{proposition}

\begin{proof}
Since all the fibers admit canonical singularities, the restriction of the canonical divisor satisfies $K_X|_{X_{t}} \cong K_{X_{t}}$ for any $t\in S$. 

By contradiction, suppose that $K_X$ is not relatively pseudo-effective over $S$, then there is an uncountable set $T\subset  S$ such that $K_{X_t}$ is not pseudo-effective for any $t\in T$. By Proposition \ref{BDPPCanonical}, $X_t$ is uniruled for any $t\in T$. Therefore, by Lemma \ref{Lemma: NegCurve}, there is a family of $K_{X_t}$-negative rational curves dominating $X_t$ for each $t\in T$. Denote by $\mathcal{C}_t\subset \mathcal{D}(X/S)$ the corresponding representative of this family in the relative Douady space $\mathcal{D}(X/S)$. Since the relative Douady space has only countably many components, there is a component $\mathcal{D}\subset \mathcal{D}(X/S)$ and an uncountable subset $T' \subset T$ of $S$, such 
that $$\mathcal{C}_t\subset \mathcal{D}\ \text{for any } t\in T'.$$
Since $S$ is a relatively compact curve, the identity principle implies $\overline{T'} = S$ (where $\overline{T'}$ is the Zariski closure of $T'$). Therefore, $\mathcal{D} \to S$ is proper and surjective. 

Let $
\mathcal{C} = \bigcup_{s \in T'} \mathcal{C}_s \subset \mathcal{D}$, 
and $\overline{\mathcal{C}}$ the Zariski closure of $\mathcal{C}$ in $\mathcal{D}$. Since $T'$ is Zariski dense in $S$, $\overline{\mathcal{C}} \to S$ is proper surjective. Let $\mathcal{X}\to \overline{\mathcal{C}}$ be the restriction of the universal family over $\mathcal{D} \to S$ onto $\overline{\mathcal{C}}\to S$. We then have the commutative diagram
\begin{center}
\begin{tikzcd}[ampersand replacement=\&]
	{\mathcal{X}} \& X \\
	{\overline{\mathcal{C}}} \& {S.}
	\arrow["\Phi", from=1-1, to=1-2]
	\arrow[from=1-1, to=2-1]
	\arrow["f", from=1-2, to=2-2]
	\arrow["\pi"', from=2-1, to=2-2]
\end{tikzcd}
\end{center}

Since
$\Phi \colon \mathcal{X} \to X$ is proper, $Y := \Phi(\mathcal{X}) \subset X$ is closed. For any $t \in T'$, $\mathcal{C}_t$ gives a covering family of the fiber $X_t$, and therefore
$$
\bigcup_{t \in T'} X_t \subset Y.
$$
Since $T'$ is Zariski dense in $S$,
it follows that $\Phi \colon \mathcal{X} \to X$ is surjective, i.e.,  the image $Y = X$.

Since $K_X$ is relatively pseudo-effective over $U$, there exists some $s \in U$, such that $K_{X_s}$ is pseudo-effective. Then $K_{X_s} + \varepsilon \omega_s$ is big for any $\varepsilon > 0$. Let $[C] \in \overline{\mathcal{C}}$ be a general curve that lies in $X_s$, whose existence follows from the surjectivity of $\Phi :\mathcal{X} \to X$. Then we have
\[
0 < (K_{X_s} + \varepsilon \omega_s) \cdot C.
\]
On the other hand, for $t \in T'$
\[
K_{X_t} \cdot C' < 0, \quad \text{for } [C'] \in \mathcal{C}_t.
\]
Therefore, for $0 < \varepsilon \ll 1$, we deduce that
\[
0 < (K_{X_s} + \varepsilon \omega_s) \cdot C = (K_X + \varepsilon \omega) \cdot C = (K_X + \varepsilon \omega) \cdot C' = (K_{X_t} + \varepsilon \omega_t) \cdot C' < 0,
\]
where the second equality holds, since $[C], [C'] \in \overline{\mathcal{C}} \subset \mathcal{D}$. This leads to a contradiction. 
\end{proof}
Combining Theorem \ref{Deformation invariance of pseudo-effectivity} and Proposition \ref{pseflocalglobal}, we obtain the global deformation invariance of the pseudo-effectivity of canonical divisors. 
\begin{theorem}\label{GlobalPseff}
Let $f: X \to S$ be a family from a K\"ahler variety $X$ onto a smooth, connected, and relatively compact curve $S$. Assume that $X_0$ is projective, and $X_t$ have canonical singularities for all $t\in S$. Then $K_{X_0}$ is pseudo-effective if and only if $K_{X_t} $ are pseudo-effective for all $t \in S \setminus \{0\}$.
\end{theorem}
\begin{proof}
The `if part' follows from Proposition \ref{Def-Close-Pseff}, while the converse follows from Theorem \ref{Deformation invariance of pseudo-effectivity} and Proposition \ref{pseflocalglobal}.
\end{proof}

Based on the results above, we proved the global stability of uniruled structures.

\begin{corollary}\label{DefUniruled}
Let $f: X \to S$ be a family from a K\"ahler variety $X$ onto a smooth, connected, and relatively compact curve $S$. Assume that $X_0$ is projective, and $X_t$ have canonical singularities for all $t\in S$. Then $X_0$ is uniruled if and only if $X_t$ are uniruled for all $t \in S \setminus \{0\}$.
\end{corollary}

\begin{proof}%[Proof of Corollary \ref{DefUniruled}
Suppose that $X_0$ is uniruled, then $K_{X_0}$ is not pseudo-effective by Lemma \ref{BDPPCanonical}  and hence the $K_{X_0}$-MMP ends with a Mori fiber space. By Proposition \ref{ExtendMMP}  this MMP can be extended over a neighborhood $U\subset S$ of $0$. Therefore $X_t$ are uniruled for all $t\in U$.
Arguing as in the proof of Proposition \ref{pseflocalglobal}, $X_t$ are uniruled  for all $t\in S$.

The converse follows directly by Lemma \ref{BDPPCanonical} and Theorem  \ref{GlobalPseff}.
\end{proof}


\begin{remark}\label{Remark: Psef3fold}
Note that \cite[Theorem 2]{DefGeneralType} already implies that if $f:X\to S$ is a flat, proper morphism of reduced, complex analytic
spaces such that the special projective fiber $X_0$ has canonical singularities and is uniruled, then the same holds over some open neighborhood $0\in U\subset S$. It would be interesting to see to what extent Corollary \ref{DefUniruled} holds without the K\"ahler assumption.


Note that, if the MMP and the transcendental base point free conjecture hold on the central fiber, the same argument 
% as in the proof of Theorem \ref{Deformation invariance of pseudo-effectivity} 
shows that Theorem \ref{GlobalPseff} and Corollary \ref{DefUniruled} also hold.

Moreover, the MMP exists and the transcendental base point free theorem holds for K\"ahler threefolds (cf. \cite{DH24}). Therefore, Theorem \ref{GlobalPseff} and Corollary \ref{DefUniruled} also hold for K\"ahler families of threefolds, without assuming that $X_0$ is projective.
\end{remark}


%%%%%%%%%%%%%%%%%%%%%%%%%%%%%%%%%%%%%%%%%%%%%%%%%%%%%%%%%%%%%%%%%%%%%%%%%%%%%%%%%%%%%%%%%%%%%%%%%%%%%%%%%%%%%%%%%%%%%%%%%%%%%%%%%
%%%%%%%%%%%%%%%%%%%%%%%%%%%%%%%%%%%%%%%%%%%%%%%%%%%% Section Deformation Invariance of Volume %%%%%%%%%%%%%%%%%%%%%%%%%%%%%%%%%%%%%
%%%%%%%%%%%%%%%%%%%%%%%%%%%%%%%%%%%%%%%%%%%%%%%%%%%%%%%%%%%%%%%%%%%%%%%%%%%%%%%%%%%%%%%%%%%%%%%%%%%%%%%%%%%%%%%%%%%%%%%%%%%%%%%%%
\section{Invariance of the volumes of adjoint classes for K\"ahler families}\label{Section: DeformationVol}
In the final section, we study the deformation behavior of volume and plurigenera in K\"ahler families.

%%%%%%%%%%%%%%%%%%%%%%%%%%%%%%%%%%%%%%%%%%%%%%%%%%%%%%%%%%%%%%%%%%%%%%%%%%%%%%%%%%%%%%%%%%%%%%%%%%%%%%%%%%%%%%%%%%%%%%%%%%%%%%%%%
%%%%%%%%%%%%%%%% Subsection Deformation Invariance of Volume When the Adjoint Class on the Central Fiber Is Big%%%%%%%%%%%%%%%%
%%%%%%%%%%%%%%%%%%%%%%%%%%%%%%%%%%%%%%%%%%%%%%%%%%%%%%%%%%%%%%%%%%%%%%%%%%%%%%%%%%%%%%%%%%%%%%%%%%%%%%%%%%%%%%%%%%%%%%%%%%%%%%%%%
\subsection{Invariance of volume when central fiber has big adjoint class}
We first prove one of the main results of the paper: the deformation invariance of the volume when the central fiber is projective and the adjoint class is big. 
\begin{theorem}\label{InvarVol}
Let $f: X \to S$ be a smooth family from a K\"ahler manifold $X$ onto a smooth, connected, and relatively compact curve $S$. For a point $0\in S$, assume that $(X,B + \bbeta)$ is a generalized K\"ahler pair such that the boundary divisor $B$ does not contain the fiber $X_0$ and let $\omega$ be a K\"ahler class on $X$ with $\omega_{X_t}:=\omega|_{X_t}$ its restriction to the fiber $X_t$. 

Assume that there exists $\delta>0$ such that $K_{X_0} + B_0 +  \boldsymbol{\beta}_{X_0}+\delta {\omega}_{X_0}$ is big and $(X_0, B_0+  { \boldsymbol{\beta}}_{X_0}+ \delta\boldsymbol{\omega}_{X_0})$ is projective, with canonical singularities, $\lfloor{ B_0\rfloor} = 0$ and $N(K_{X_0 }+B_0 + \bbeta_{X_0})\wedge B_0  = 0$. Then the volume  function 
$$t \longmapsto \mathrm{vol}(K_{X_t}+B_t +  \boldsymbol{\beta}_{X_t} +\delta{\omega}_{X_t})$$ 
is constant on a Euclidean open neighborhood $U \subset S$ of $0$.
\end{theorem}

\begin{proof}
By Theorem \ref{Theorem: TransBCHM}, we can run the generalized $(K_{X_0}+B_0 + { \bbeta}_{X_0} +  \delta {\omega}_{X_0})$-MMP with scaling of $\omega_{X_0}$ on the central fiber $X_0$, and it will terminate at some minimal model
$$X_0 \dashrightarrow X_0^{(1)} \dashrightarrow X_0^{(2)} \dashrightarrow \cdots \dashrightarrow X_0^{(n)} \dashrightarrow X_0'.$$
Since $\omega _{X_0}$ is K\"ahler, $N(K_{X_0} + B_0 + \bbeta_{X_0}+ \delta{\omega}_{X_0}) \le  N(K_{X_0} + B_0 + \bbeta_{X_0})$. Therefore $N(K_{X_0} + B_0 + \bbeta_{X_0}+ \delta{\omega}_{X_0})\wedge B_0 = 0$. By Proposition \ref{ExtendMMP}, after shrinking the base to some (connected) open neighborhood $U \subset S$ of $0$, we can extend the steps of the MMP on $X_0$ to a sequence of $(K_X+B+\bbeta_X+ \delta{\omega}_X)$-negative proper bimeromorphic maps over $U$, $$X /U \dashrightarrow X^{(1)}/U \dashrightarrow \cdots \dashrightarrow X'/U\dashrightarrow \cdots.$$
Since $K_{X_0'}  + B_0' +  { \bbeta}_{X_0'}+ \delta {\omega}_{X_0'}$ is big and nef, by Proposition \ref{Extbigandnef}, after shrinking $U$, $K_{X'} +B' +  { \bbeta}_{X'} + \delta {\omega}_{X'}$ is relatively nef over $U$. By Lemma \ref{BigVol}, $$\mathrm{vol}(K_{X'_t}+ B_t' +  { \bbeta}_{X_t'}+\delta {\omega}_{X_t'}) = (K_{X'_t}+ B_t' +  { \bbeta}_{X_t'} + \delta {\omega}_{X_t'})^m,$$where  $m = \dim_{\mathbb{C}}(X/S)$. 
We claim that $t \longmapsto(K_{X'_t} + B_t' +  { \bbeta}_{X_t'}+ \delta {\omega}_{X_t'})^m$ is constant for $t\in U$. Pick an arbitrary point $s\in U$. There exists a real smooth connected curve $\gamma: [0,1] \to U$ with  $\gamma(0) = 0,\gamma(1) = s$, so that the inverse images $V = (f')^{-1}(\gamma(0,1))$ of $f':X'\to U$ and  $\overline{V} = (f')^{-1}(\gamma([0,1]))$ are analytic subvarieties, with topological boundary $\partial V = X_0 '\cup X_s'$. Set $$Z = \{x\in X' \mid f' \text{ is submersion at } x\in X' \},$$ and $M_V : = {V} - {V} \cap Z$ with $\partial   :=  \partial \overline{V} - \partial \overline{V} \cap Z$. By the constant rank theorem, the inverse image of the smooth connected curve under a submersion is a smooth manifold. Hence $M_V = (f'|_{X'- Z})^{-1}(\gamma (0,1))$
is smooth, such that $(M_V ,\partial )$ is a manifold with boundary. Since $X_0'$ and $X_s'$ are compact, we know that the volume of $\partial \overline{V}$ (and hence the volume of $\partial$) is finite. 
    
Since $f:X\to S$ is a smooth surjective morphism, $f$ is submersive at every point $x\in X$. Since $X'\dasharrow X$ is an isomorphism outside of a (complex) codimension 2 subset, it follows that
$\mathrm{codim}_{X'}Z\ge 2 $. By the singular Stokes' formula in Proposition \ref{stokes} below, $$(K_{X_s'}+ B'_s  + \bbeta_{X_s'}  + \delta{\omega}_{X_s'})^m-(K_{X_0'}+ B'_0 + \bbeta_{X_0'} + \delta{\omega}_{X_0'})^m  = \int_{M_V}d (K_{X'}+B'+\bbeta_{X'}+ \delta {\omega}_{X'})^m =  0.$$By the proof of Proposition \ref{ContInMMP} and \cite[Corollary 3.2]{CRC25}, 
$$\vol(K_{X_t}+ B_t + \bbeta_{X_t}+ \delta {\omega}_{X_t}) = \vol (K_{X_t'}+ B_t' + \bbeta_{X_t'} + \delta {\omega}_{X_t'}),\quad \text{ for }t\in U.$$ And therefore
$$t \longmapsto \mathrm{vol}(K_{X_t}+B_t +  \boldsymbol{\beta}_{X_t} +\delta{\omega}_{X_t})$$ is constant in the neighborhood $0\in U$. 
\end{proof}


\begin{proposition}[{\cite[p.10]{Stolzenberg}}]\label{stokes}
Let $M$ be a smooth, relatively compact, $n$-dimensional real manifold with finite $n$-volume, and  $\overline{M}$ its compactification. 
Assume that $\overline{M} \setminus M$ can be decomposed as $\partial \cup S$, where\\
\emph{(1)} $(M, \partial)$ is a manifold with boundary,\\
\emph{(2)} $\partial$ has finite $(n-1)$-volume,\\
\emph{(3)} $S$ is compact with $\operatorname{codim}_{\mathbb{R}}(S, M) \ge 2$.\\
Then for any smooth $(n-1)$-form $\alpha$,
\[
\int_{\partial} \alpha = \int_M d\alpha.
\]
\end{proposition}

Based on Theorem \ref{InvarVol}, we also obtain the following variant under the assumption that $(X,B)$ is log smooth over $S$, without assuming that $N(K_{X_0} + B_0 + \bbeta_{X_0})\wedge B_0  = 0$. Let $f : X \to S$ be a proper surjective morphism from a normal complex variety $X$ to a relatively compact complex variety $S$. We say that a log pair $(X,B)$ is \emph{log smooth over $S$} if $f : X \to S$ is a smooth morphism and, for each stratum $B_I$ of $B$, the restriction $B_I \to S$ is also a smooth morphism.
 
\begin{theorem}\label{InvarVol2}
Let $f:X\to S$ be a smooth family from a K\"ahler manifold $X$ onto a smooth, connected, relatively compact curve $S$. Assume that $(X,B+\bbeta)$ is a generalized klt pair such that $(X,B)$ is log smooth over $S$ and $\bbeta =\overline \beta $ where $\beta$ is nef over $S$, and $X_0$ is projective. If $K_{X_0}+B_0+\bbeta _{X_0}$ is big, then the volume function
\[t\longmapsto {\rm vol}(K_{X_t}+B_t+\bbeta _{X_t})\]
is constant on some Euclidean neighborhood $U\subset S$ of $0$.
\end{theorem}
\begin{proof}
After an \'etale base change, we may assume that $f|_Z$ has connected fibers for any stratum $Z$ of the support of $B$. After blowing up strata of $B$ via a morphism $\nu :X'\to X$ and writing $K_{X'}+B'=\nu^*(K_X+B)+E$ where $B',E\ge 0$ and $B'\wedge E=0$, we may replace $(X,B)$ by $(X',B')$. In particular, we may assume that the components of $B$ are disjoint and so $(X,B)$ is terminal. Let $0\leq B^\sharp\leq B$ be the unique $\R$-divisor such that  $$B^\sharp _0=B_0-B_0\wedge N(K_{X_0}+B_0+\bbeta_{X_0}).$$ Arguing as in the proof of Theorem \ref{InvarVol}, there is a $(K_{X_0}+B^\sharp_0+\bbeta _{X_0})$-minimal model $X_0\dasharrow X_0'$ which extends to a $(K_X+B^\sharp+\bbeta)$-negative proper bimeromorphic map $X\dasharrow X'/U$ over a neighborhood $0\in U\subset S$. Note that moreover, if $B^\varepsilon =B^\sharp+\varepsilon(B-B^\sharp)$, then $X\dasharrow X'$ is $(K_X+B^\varepsilon+\bbeta_X)$-negative for $0\leq \varepsilon \ll 1$. We then have \[{\rm vol}(K_{X_t}+B_t^\varepsilon+\bbeta_{X_t})={\rm vol}(K_{X'_t}+{B'_t}^\varepsilon+\bbeta_{X'_t})={\rm vol}(K_{X'_0}+{B'_0}^\varepsilon+\bbeta_{X'_0})={\rm vol}(K_{X_0}+B_0^\varepsilon+\bbeta_{X_0})\]
for $t\in U$ and $0\leq \varepsilon \ll 1$.
This immediately implies that 
\[{\rm vol}(K_{X_0}+B_0+\bbeta_{X_0})={\rm vol}(K_{X_0}+B^\sharp _0+\bbeta_{X_0})={\rm vol}(K_{X_t}+B^\sharp _t+\bbeta_{X_t})\leq {\rm vol}(K_{X_t}+B _t+\bbeta_{X_t}),\]where the first equality follows from Proposition \ref{VolZariski}.
Suppose that the above inequality is sharp for some $0\ne t_i\to 0$, then we also have 
\[
\begin{aligned}
{\rm vol}\!\bigl(K_{X_{t_i}} + B^\varepsilon_{t_i} + \bbeta_{X_{t_i}}\bigr)^{\frac{1}{m}}
&\ge (1 - \varepsilon)\, {\rm vol}\!\bigl(K_{X_{t_i}} + B^\sharp_{t_i} + \bbeta_{X_{t_i}}\bigr)^{\frac{1}{m}}
+ \varepsilon\, {\rm vol}\!\bigl(K_{X_{t_i}} + B_{t_i} + \bbeta_{X_{t_i}}\bigr)^{\frac{1}{m}} \\
&> {\rm vol}\!\bigl(K_{X_0} + B^\varepsilon_0 + \bbeta_{X_0}\bigr)^{\frac{1}{m}},
\end{aligned}
\]
where $m = \dim_{\mathbb{C}}(X/S)$, and the first inequality is implied by the log-concavity of the volume (cf. Proposition \ref{r-logc}). But then ${\rm vol}(K_{X_{t_i}}+B^\varepsilon _{t_i}+\bbeta_{X_{t_i}})>{\rm vol}(K_{X_0}+B^\varepsilon _0+\bbeta_{X_0})$ for $0<\varepsilon \ll 1$, contradicting what we have shown above.
\end{proof}

\begin{remark}
In \cite[Remark 4.10]{FH11}, de Fernex--Hacon showed that the deformation invariance of log plurigenera does not hold when the family is not log smooth. Consider the family of smooth quadric degenerating to the Hirzebruch surface $\mathbb{F}_2$. A line in one of the two rulings on the general fiber $X_t$ sweeps out under deformation over $T$, a divisor $R$ on $X$ with the following properties\\
(a) $\left.R\right|_{X_0}=E+F$, where $F$ is a fiber of the projection $\mathbb{F}^2 \rightarrow \mathbb{P}^1$, and $E$ is the tautological section on $\mathbb{F}_2$ whose intersection number is $E^2 = -2$. \\
(b) $R|_{X_t} = F_t$ for $t
 \ne 0$, where $F_t$ is one of the ruling of the quadrics.

For any $b \ge 1$, consider a divisor $L$ on $X$ such that $\left.L\right|_{X_0} \sim(b+1) E+2 b F$ and on the general fiber $L|_{X_t} = (b+1)F_0 + (b-1)F_1$ (where $F_0$ and $F_1$ are the rulings of the quadrics). Moreover, if we define $D:=\frac{1}{m}(H+2 R)$, for a general $H \in\left|-\frac{b-1+2 m}{2} K_X\right|$ with $m\gg 1$, then $$L \sim m\left(K_X+D\right).$$
The pair $(X,D)$ is klt but not log smooth since the divisor $R$ splits on the central fiber. In this case, the log plurigenera are not constant since on the general fiber $$ P_m (K_{X_t}+D_t)=h^0\left(\mathcal{O}_{X_t}(L)\right)=h^0\left(\mathcal{O}_{\mathbb{P}^1 \times \mathbb{P}^1}(b+1, b-1)\right)=b^2+2 b,$$while on the central fiber $|L|_{X_0}|=| b(E+2 F)+E|=| b(E+2 F) \mid+E$ and therefore $$P_m (K_{X_0} + D_0) = h^0\left(\mathcal{O}_{X_0}(L)\right)=\binom{b+3}{3}-\binom{b+1}{3}=b^2+2 b+1.$$The plurigenera jump on the central fiber, however the volume is constant in this family.
\end{remark}

%%%%%%%%%%%%%%%%%%%%%%%%%%%%%%%%%%%%%%%%%%%%%%%%%%%%%%%%%%%%%%%%%%%%%%%%%%%%%%%%%%%%%%%%%%%%%%%%%%%%%%%%%%%%%%%%%%%%%%%%%%%%%%%%%
%%%%%%%%%%%%%%%% Subsection Deformation Invariance of Volume When the Adjoint Class on the Central Fiber Is Big%%%%%%%%%%%%%%%%
%%%%%%%%%%%%%%%%%%%%%%%%%%%%%%%%%%%%%%%%%%%%%%%%%%%%%%%%%%%%%%%%%%%%%%%%%%%%%%%%%%%%%%%%%%%%%%%%%%%%%%%%%%%%%%%%%%%%%%%%%%%%%%%%
\subsection{Invariance of volume for families of K\"ahler threefolds}
In this subsection, we prove the deformation invariance of the plurigenera and the volume of an adjoint class for smooth families of K\"ahler threefolds.

\begin{theorem}\label{DefVol3fold}
Let $f: X \to S$ be a smooth family of compact K\"ahler threefolds over a smooth, connected, and relatively compact curve $S$. Then for any integer $m\ge 1$, the $m$-genus $P_m(X_t)$ is constant for any $t\in S$. Furthermore, assume that there exists a K\"ahler class $\omega_X$ on $X$ such that $\omega_{X_t} :  =  \omega_X |_{X_t}$. Then for every $\delta\ge 0$, the function $$t \longmapsto \operatorname{vol}(K_{X_t}+\delta\,\omega_{X_t})$$is constant for any $t\in S$.
\end{theorem}
\begin{proof}
We first consider the deformation invariance of plurigenera. Using the similar extension technique (cf. Proposition \ref{ExtendMMP}), \cite[Theorem 1.2]{DHP24} proved the existence of relative MMP for smooth family of K\"ahler threefolds. Then, the invariance of plurigenera follows from \cite[Theorem 11]{Naka86}, as the abundance conjecture is completely known for K\"ahler threefolds by \cite{CHP16, CHP23Erratum, DO24, DO24II}. 

We next consider the deformation invariance of the volume. It is easy to see that the case $\delta=0$ follows from the deformation invariance of plurigenera established above. Fix a point $t_0 \in S$, and let $
\tau := \inf\{ t \ge 0 \mid K_{X_{t_0}} + t\,\omega_{X_{t_0}} \text{ is pseudo-effective} \}$. We now assume that $\delta>0$ and divide the argument into three cases according to the relation between $\delta$ and $\tau$.


Case 1. If $\delta >  \tau$, then $K_{X_{t_0}}+ \delta{\omega}_{X_{t_0}}$ is big. Following a similar line of argument as in Theorem \ref{InvarVol}, and using \cite[Theorem 1.4]{DH24}, we can run the generalized $(K_{X_{t_0}} + \delta {\omega}_{X_{t_0}})$-MMP with scaling of $\omega_{X_{t_0}}$ on the central fiber, and it will terminate at some minimal model $$(X_{t_0},\delta \boldsymbol{\omega}_{X_{t_0}}) \dashrightarrow (X_{t_0}^{(1)} , \delta \boldsymbol{\omega}_{X_{t_0}^{(1)}})\dashrightarrow  \cdots \dashrightarrow (X_{t_0}^{(n)}, \delta \boldsymbol{\omega}_{X_{t_0}^{(n)}}) \dashrightarrow (X_{t_0}', \delta \boldsymbol{\omega}_{X_{t_0}'}).$$

By Remark \ref{Remark:ExtMMP} of Proposition \ref{ExtendMMP}, after shrinking the base to some (connected) open neighborhood $U \subset S$ of ${t_0}$, we can extend the steps of the MMP on $X_{t_0}$ to a sequence of $(K_X+ \delta{\omega}_X)$-negative proper bimeromorphic contractions over $U$, $$(X , \delta \boldsymbol{\omega}_X)/U \dashrightarrow (X^{(1)},\delta \boldsymbol{\omega}_{X^{(1)}})/U \dashrightarrow \cdots \dashrightarrow (X', \delta \boldsymbol{\omega}_{X'})/U\dashrightarrow \cdots.$$
Since $K_{X_{t_0}'} + \delta {\omega}_{X_{t_0}'}$ is big and nef, Remark \ref{Extbigandnef3fold} implies that after shrinking the base $K_{X'} + \delta {\omega}_{X'}$ is relatively nef over $U$. Then, by the singular Stokes' formula in Proposition \ref{stokes}, the volume function $$t \longmapsto \operatorname{vol}\!\left(K_{X_t}+\delta {\omega}_{X_t}\right) = (K_{X_t'}+ \delta{\omega}_{X_t'})^3$$
is constant for $t\in U$.

Case 2. If $0 < \delta < \tau$, then $K_{X_{t_0}}+\delta \omega_{X_{t_0}}$ is not pseudo-effective. We can run the generalized $(K_{X_{t_0}} + \delta {\omega}_{X_{t_0}})$-MMP with scaling of $\omega_{X_{t_0}}$ on the central fiber, and it will terminate at some Mori fiber space $X_{t_0}' \to Z_{t_0}'$ (with $-(K_{X_{t_0}'}+\delta \omega_{X'_{t_0}})$ relative K\"ahler over $Z_{t_0}'$). The same as in the Case 1, after shrinking the base to some (connected) open neighborhood $U \subset S$ of ${t_0}$, we can extend the steps of the MMP on $X_{t_0}$ to a sequence of $(K_X+ \delta{\omega}_X)$-negative proper bimeromorphic contractions over $U$, $$(X , \delta \boldsymbol{\omega}_X)/U \dashrightarrow (X^{(1)},\delta \boldsymbol{\omega}_{X^{(1)}})/U \dashrightarrow \cdots \dashrightarrow (X', \delta \boldsymbol{\omega}_{X'})/U\dashrightarrow \cdots.$$Then, by Proposition~\ref{extFano}, after shrinking $U$ again, one can extend the Mori fiber space $X'_{t_0} \to Z'_{t_0}$ to a relative Mori fiber space $X' \to Z'/U$ over $U$. It follows that $K_{X_t} + \delta \omega_{X_t}$ is not pseudo-effective for $t \in U$. Therefore
$$
\vol(K_{X_t} + \delta \omega_{X_t}) = \vol(K_{X_{t_0}} + \delta \omega_{X_{t_0}}) = 0,
$$
for $t \in U$.

Case 3. Suppose that $0<\delta = \tau$. Then $K_{X_{t_0}}$ is not pseudo-effective and thus, we can run a $K_{X_{t_0}}$-MMP with scaling of $\omega_{X_{t_0}}$ on $X_{t_0}$, which terminates with a $(K_{X_{t_0}'} + \tau \omega_{X_{t_0}'})$-trivial Mori fiber space $X_{t_0}' \to Z_{t_0}'$ (with $-K_{X_{t_0}'}$ relative ample over $Z_{t_0}'$). By Proposition~\ref{extFano}, this Mori fiber space on $X_{t_0}'$ extends to a relative Mori fiber space $X' \to Z'/U$ over some neighborhood $U$ of $t_0$. Since this Mori fiber space is projective, by Lemma~\ref{l-bir}, the induced Mori fiber spaces on the nearby fibers are $(K_{X_t'}+\tau \omega_{X_t'})$-trivial (and $K_{X_t'}$-negative). By Lemma \ref{Lemma: PSEFandMFS}, it follows that the pseudo-effective threshold is locally constant and equal to $\tau$ in a neighborhood of $t_0$. Consequently, the volume function $t \longmapsto \vol(K_{X_t} + \tau \omega_t)$ is equal to $0$ in a neighborhood of $t_0$ and locally constant.


In all three cases, the volume function $t\longmapsto \vol(K_{X_t} + \delta {\omega}_{X_t})$ is locally constant around any point $t_0 \in S$. Since $S$ is connected, the volume function is constant for any $t\in S$. 
\end{proof}

From the above results, it is natural to propose:
\begin{conjecture}
Let $f:X\to S$ be a smooth family from a K\"ahler manifold $X$ onto a smooth, connected, relatively compact curve $S$. Assume that $(X,B+\bbeta)$ is a generalized klt pair such that $(X,B)$ is log smooth over $S$ and $\bbeta =\overline \beta $ where $\beta$ is nef over $S$. Then the volume function
\[t\longmapsto {\rm vol}(K_{X_t}+B_t+\bbeta _{X_t})\]
is constant over $S$.
\end{conjecture}

%{\bf COI statement.} On behalf of all authors, the corresponding author states that there is no conflict of interest. 

%{\bf Data availability statement.} The authors declare that this manuscript has no associated data.


\bibliographystyle{amsalpha}
\bibliography{mybib.bib}
\end{document}
